\chapter{微分方程与振动}\label{ch:c2}

排水、绕绳和阻力运动都可以从过程的收支关系建立微分方程。状态变量随模型选取：液面高度可直接决定出流速率，质点运动通常还需同时给出位置与速度。初始条件确定解，约束边界确定这段解的适用区间。可分离变量时应保留平衡解，采用能量降阶时应保留速度的正负支。

本讲先由具体过程列式，再推导一阶线性方程和二阶常系数方程的解法。阻尼与受迫振动的解始终带入初值，分别计算瞬态、稳态以及位移峰和功率峰。接触条件或外源改变时，以边界状态连接下一段运动。


\section{由过程建立方程}\label{sec:c2-model}
微分方程来自局部过程的受力或守恒关系。状态变量、参考系和正方向确定后，对同一时间段或微元列式，再补上初始或边界条件。

\subsection{状态、初值与终止事件}
微分方程把同一时刻的状态与变化率联系起来。例如$\dot x=v$、$m\dot v=F(x,v,t)$构成两个一阶方程，也可合并为一个二阶方程。求解前须给出坐标方向、参数的物理含义和初始状态；到达边界、速度反向或接触消失时，方程的适用分支还可能改变。

对于$\dot y=f(y)$，若$f$在所讨论的分支上非零，可分离为$\dd t=\dd y/f(y)$，从而
\begin{equation}t-t_0=\int_{y_0}^{y(t)}\frac{\dd\eta}{f(\eta)}.
\end{equation}
积分上限随当前状态$y(t)$变化。还须单列$f(y_*)=0$对应的平衡解，因为分离时除以$f$已将它排除。边界处的积分是否收敛，还决定到达该边界所需时间是否有限。

二阶方程的初态通常需要位置与速度两个条件。若方程来自能量守恒，平方根只给出速度大小，正负支还须由初始方向和转向过程选择。对于无显式时间的方程，也可设$v=v(x)$，用$\ddot x=v\dd v/\dd x$降低阶数。

\Needspace{8\baselineskip}
\subsection{由收支关系建立状态方程}
状态变量应足以确定后续演化。质点在给定位置可以有不同速度，所以一般需要$(x,v)$两个量；水缸在准稳态模型中出流速率已由水深决定，故只需$h$。约束可减少独立变量，但消去变量时必须保留它带来的动能、流量或几何关系。

建立微分方程通常先写有限时间内的收支。若$Q$为某累计量，流入、流出率分别为$J_{\rm in},J_{\rm out}$，则
\[Q(t+\Delta t)-Q(t)=\int_t^{t+\Delta t}(J_{\rm in}-J_{\rm out})\dd s.\]
除以$\Delta t$取极限后得到$\dot Q=J_{\rm in}-J_{\rm out}$。水缸中$Q$可取体积，电容中可取电荷；对质点动量，右侧则是合力。先规定$Q$增加为正，流出项便以负号进入收支式。

若$Q=Q(y)$，方程应为$Q'(y)\dot y=J(y,t)$。截面变化的水缸有$\dd V/\dd h=A(h)$，因此液面变化率的系数是当前水深处的截面积，而非总容积除以高度。系统参数改变时，应回到这条局部守恒关系，不能只在已知答案中替换平均值。

运动模型还包含停止或换支规则。水深降至零后不再继续排水；绕绳自由段缩至零时，质点--轻绳模型到达其几何边界；相对绕入角速度变零时，则可能开始反向相对运动。方程及初值确定一段解，事件条件确定这段解何时终止。

\subsection{分离变量中的符号与唯一性}
在$f(y)\ne0$的连通区间内，函数$1/f(y)$有确定符号，时间可由状态积分表示。若解从$y_0$下降至$y_1<y_0$，则$\dot y<0$，分子$\dd y$和积分方向也为负，得到的用时仍为正。取平方根后只保留速率、却忘记实际运动方向，往往会在这一步得到负时间。

$f(y_*)=0$对应平衡解。分离变量的除法将它删去了，因此求通解后须单列。若$f$在$y_*$附近满足局部Lipschitz条件，经过同一初态的解局部唯一，非平衡解不能在有限时间穿过平衡点后继续到另一侧。习题~\ref{ex:c2-h5}中的$0<\nu<1$是需要单独讨论的边界：$y=0$附近导数无界，常规唯一性定理的假设不再满足。

判断到达时间只需检查$\int\dd y/f(y)$在目标边界的收敛性。有限到达不要求变化率保持非零；只要它趋零不太快，时间积分仍可有限。反过来，状态曲线在图上已十分接近目标值，并不表示解析解在有限时刻真正到达。水缸排空与线性阻力终速分别给出这两种情形。

\begin{example}[水缸泄水]\label{ex:c2-e1}
竖直水缸的水平截面积为常数$S_0$，初始水深$h_0$。底部开一面积$S\ll S_0$的小孔，两处均通大气，忽略黏性和收缩，采用准稳态理想流模型。求水深$h(t)$及排空时间$T$。 见图~\ref{fig:c2-e1-q}。
\end{example}

\begin{figure}[H]\centering
\begin{tikzpicture}[>=Stealth,line cap=round,line join=round]
\draw[thick](0,3.2)--(0,0)--(3,0)--(3,.35);\draw[thick](3,.55)--(3,3.2);\fill[gray!15](0,0)rectangle(3,2.5);\draw[curve](0,2.5)--(3,2.5);\draw[axis](3,.45)--(4,.45)node[right]{$v$};\draw[<->](-.4,0)--(-.4,2.5)node[midway,left]{$h$};\node at(1.5,1.5){$S_0$};\node[right]at(3.1,.9){孔面积$S$};
\end{tikzpicture}
\caption{高度向上量，孔口流速向外，液面下降。}\label{fig:c2-e1-q}
\end{figure}
\begin{solution}
未知量是液面高度$h(t)$。图~\ref{fig:c2-e1-q}中的小孔面积远小于液面面积，使液面动能可在首阶忽略；两处压强相同，重力势差主要转为出口动能。先由伯努利关系求当前水深下的出流速度，再以体积连续性把它转成液面下降率。


相同一段时间内，流出体积等于液面下降减少的体积。按图~\ref{fig:c2-e1-s}的方向，
\[S_0\dot h=-Sv,\qquad v\simeq\sqrt{2gh}.
\]
因此$\dot h=-(S/S_0)\sqrt{2gh}$，负号来自$h$向上而水面向下。分离变量并在两端同时代入初态，
\begin{align*}
\int_{h_0}^{h(t)}h^{-1/2}\dd h
&=[2\sqrt h]_{h_0}^{h(t)}\\
&=-\frac S{S_0}\sqrt{2g}\int_0^t\dd s
=-\frac S{S_0}\sqrt{2g}\,t,\\
\sqrt{h(t)}&=\sqrt{h_0}-\frac S{S_0}\sqrt{\frac g2}\,t.
\end{align*}
故在$0\le t\le T$上，
\[\boxed{h(t)=h_0(1-t/T)^2,\qquad
T=\frac{S_0}{S}\sqrt{\frac{2h_0}g}}.
\]
排空时间也可写为$T=\int_{h_0}^{0}\dd h/\dot h$，积分方向与负的$\dot h$共同给出正时间。

若不忽略液面速度，但仍采用同一理想准稳态模型，令 $\varepsilon=S/S_0<1$。连续性给出液面速度大小 $u=\varepsilon v$，伯努利关系 $v^2-u^2=2gh$ 给出 $v=\sqrt{2gh/(1-\varepsilon^2)}$，因此 $\dot h=-\varepsilon\sqrt{2gh/(1-\varepsilon^2)}$。分离变量后，时间积分只是原积分乘以 $\sqrt{1-\varepsilon^2}$，
\[
T_{\rm finite}=\frac{\sqrt{1-\varepsilon^2}}{\varepsilon\sqrt{2g}}
\int_0^{h_0}h^{-1/2}\dd h
=\frac{\sqrt{1-\varepsilon^2}}{\varepsilon}\sqrt{\frac{2h_0}{g}}.
\]
窄孔近似忽略的是相对 $O(\varepsilon^2)$ 的液面动能；真实孔口的损失另需流量系数，不能由此无损模型自行确定。


排空条件来自取平方根后的线性式，$\sqrt h$ 必须保持非负。令它第一次为零得到 $T$；此前 $\dot h<0$，以后若沿原平方式继续，会出现 $\dot h>0$，已违背原方程。

图~\ref{fig:c2-e1-s}中流出体积微元为$Sv\dd t$，液面减少体积为$-S_0\dd h$。底部附近$v\to0$，但$\dd t\propto h^{-1/2}\dd h$的奇性可积，所以排空仍在有限时间完成。整个过程的平均出流速度为初始出口速度的一半，因为排出总体积$S_0h_0$除以$ST$正好等于$\sqrt{2gh_0}/2$。

\begin{figure}[H]\centering
\begin{tikzpicture}[>=Stealth,line cap=round,line join=round]
\draw[thick](0,3.2)--(0,0)--(3,0)--(3,.35);\draw[thick](3,.55)--(3,3.2);\fill[gray!15](0,0)rectangle(3,2.5);\draw[curve](0,2.5)--(3,2.5);\draw[axis](3,.45)--(4,.45)node[right]{$v$};\draw[<->](-.4,0)--(-.4,2.5)node[midway,left]{$h$};\node at(1.5,1.5){$S_0$};\node[right]at(3.1,.9){孔面积$S$};\draw[axis](1.5,2.5)--(1.5,1.8)node[right]{$-\dot h$};\node[right]at(4.7,2.4){$S_0(-\dot h)=Sv$};
\end{tikzpicture}
\caption{体积流率联系液面下降与孔口流速。}\label{fig:c2-e1-s}
\end{figure}
\end{solution}

\Needspace{8\baselineskip}


\Needspace{8\baselineskip}


\section{线性方程与初始条件}\label{sec:c2-linear}
初态确定线性方程自由响应中的常数。变量分离和积分因子按方程形式选择，初值可直接放进积分上下限，同时保留除法可能丢失的常数解。

\subsection{线性、齐次与阶数}
未知函数及其各阶导数只以一次形式出现，且相乘的系数不依赖未知函数时，方程才是线性的。系数可以依赖自变量，例如$y^{\prime\prime}+a(x)y^{\prime}+b(x)y=f(x)$仍是线性方程；$(y^{\prime})^2$、$\sin y$或$yy^{\prime}$则是非线性项。最高出现的导数阶数决定方程的阶，而不是函数在式子中的次数。

设线性算子$\mathcal L[y]=y^{\prime\prime}+ay^{\prime}+by$。若$y_h$满足$\mathcal L[y_h]=0$，$y_p$满足$\mathcal L[y_p]=f$，则$y_h+y_p$满足原非齐次方程。两个非齐次解之差为齐次解，因此一般解可写为$y=y_h+y_p$。但任意倍数的非齐次解不再是同一外源下的解，不能把齐次叠加法则原样套用到它们上。

\subsection{一阶线性方程的积分因子}
对$y^{\prime}+p(x)y=q(x)$，令$\mu(x)=\exp(\int p(x)\dd x)$，则
\[(\mu y)'=\mu q,\qquad
y(x)=\frac{\mu(x_0)y_0+\int_{x_0}^x\mu(\xi)q(\xi)\dd\xi}{\mu(x)}.
\]
这个表达式已把初值放入积分边界。对于$q=0$，它退化为分离变量的指数衰减或增长式，并保留$y_0=0$对应的零解。

\Needspace{8\baselineskip}
\subsection{积分因子的构造与初值项}
一阶线性方程之所以可直接积分，是因为能找到一个因子把左侧凑成乘积导数。乘以未知的非零函数$\mu$后，要求
\[\mu y'+\mu p y=(\mu y)'=\mu y'+\mu' y.\]
比较$y$的系数得到$\mu'=p\mu$，于是$\mu=\exp(\int p\dd x)$。这个构造把求$y$的问题转成一个只含已知系数$p$的辅助方程。

为使初值项直接保留为$y_0$，取归一化积分因子
\[\mu(x)=\exp\left(\int_{x_0}^xp(\xi)\dd\xi\right),\quad \mu(x_0)=1.\]
在$[x_0,x]$积分得$\mu(x)y(x)-y_0=\int_{x_0}^x\mu(\eta)q(\eta)\dd\eta$。因此初态部分$y_0/\mu(x)$与外源累计部分分开出现。外源为零时仍包含$y_0=0$，不需要在求对数时另分符号。

若方程不显含未知函数本身，例如$y''+p(x)y'=q(x)$，可先设$u=y'$。第一步用积分因子求$u$，第二步从$x_0$积分$u$求$y$。两个初值分别确定导数和函数值，不能在第一次积分后就声称已完成二阶初值问题。反之，方程若只是不显含自变量，可尝试用$v(y)$降阶，但那时二阶导数变为$v\dd v/\dd y$，与令$u=y'$的一阶线性化不是同一个变换。

求得积分形式也是完整解的一种表达。已知$p,q$的积分不能化成初等函数时，带固定初值边界的积分式仍能确定唯一函数，并可据此研究符号、单调性和渐近行为。若另作近似，仍须保留积分边界，并估计它对初值解的误差。

\begin{example}[变化系数的衰减与降阶]\label{ex:c2-e4}
求解：(1)$y^{\prime}=-(a+bx)y$，$y(0)=y_0$；(2)$y^{\prime\prime}+p(x)y^{\prime}=0$，$y(x_0)=Y_0,y^{\prime}(x_0)=V_0$。讨论第一式中$b>0$时的图形，并说明第二式怎样降低阶数。 见图~\ref{fig:c2-e4-q}。
\end{example}

\begin{figure}[H]\centering
\begin{tikzpicture}[>=Stealth,line cap=round,line join=round]
\begin{scope}[x=1.00000cm,y=1.38462cm]\draw[axis](0,0)--(3.25,0)node[right]{$x$};\draw[axis](0,0)--(0,1.3)node[above]{$y/y_0$};\draw[curve,domain=0:3,samples=101]plot(\x,{exp(-.5*\x-.25*(\x)^2)});\end{scope}
\end{tikzpicture}
\caption{例示$a,b>0$时的非匀速指数衰减。}\label{fig:c2-e4-q}
\end{figure}
\begin{solution}
第一式可直接分离，但积分因子法能同时包含零初态。第二式中未知函数本身没有出现，适合把导数$y'$作为新未知量。图~\ref{fig:c2-e4-q}的示意曲线只代表某一组参数，单调性仍须由系数$a+bx$及$y$的符号判断。


第一式的积分因子为 $\mu(x)=\exp[\int_0^x(a+b\xi)\dd\xi]=\exp(ax+bx^2/2)$。按链式法则，$\mu'=(a+bx)\mu$，所以原方程乘以它后化为
\[
(\mu y)'=\mu y'+\mu'y=\mu[y'+(a+bx)y]=0.
\]
从 $0$ 积到 $x$，$\mu(x)y(x)-\mu(0)y_0=0$；利用 $\mu(0)=1$，得到
\[\boxed{y(x)=y_0\exp(-ax-bx^2/2)}.
\]
若由$\dd y/y=-(a+bx)\dd x$求解，要补回被除法排除的$y_0=0$解；积分因子法自然包含它。图~\ref{fig:c2-e4-s}画出对数幅值，$\ln|y/y_0|=-ax-bx^2/2$是抛物线，只有$b=0$时才是直线。若$a+bx$变号，$y$的单调性也随之改变，不能只见指数就判定处处衰减。

第二式不显含 $y$，令 $u=y'$，得到 $u'+p(x)u=0$。记 $P(x)=\int_{x_0}^xp(\xi)\dd\xi$，则 $P'=p$，乘积求导给出 $(\ee^Pu)'=\ee^P(u'+pu)=0$。积分并用 $P(x_0)=0,u(x_0)=V_0$，得 $\ee^{P(x)}u(x)=V_0$，即
\[u(x)=V_0\exp\left[-\int_{x_0}^x p(\xi)\dd\xi\right],
\]
\[\boxed{y(x)=Y_0+V_0\int_{x_0}^x
\exp\left[-\int_{x_0}^{\eta}p(\xi)\dd\xi\right]\dd\eta}.
\]
这里第二次积分用了 $y(x)-Y_0=\int_{x_0}^xu(\eta)\dd\eta$，两个初值分别确定两个常数。若 $p=\gamma>0$，外层积分可继续计算：
\[
y-Y_0=V_0\int_{x_0}^x\ee^{-\gamma(\eta-x_0)}\dd\eta
=V_0\left[-\frac1\gamma\ee^{-\gamma(\eta-x_0)}\right]_{x_0}^x
=\frac{V_0}{\gamma}[1-\ee^{-\gamma(x-x_0)}].
\]
导数趋于零而函数趋于另一个常数，这与直接把 $y$ 本身当作衰减指数不同。


对第一式，设 $y_0\ne0$，在不变号的解分支上分离变量积分，也有 $\ln|y(x)/y_0|=-ax-bx^2/2$。
指数式从而保留初态符号。$b>0$时可配方为$-b(x+a/b)^2/2+a^2/(2b)$；若允许整个实轴，幅度在$x=-a/b$最大，而不是总从$x=0$开始下降。图~\ref{fig:c2-e4-s}画出的对数幅值斜率为$-(a+bx)$，可以直接读出这种转折。

对第二式求导验证，变上限积分给出 $y'=V_0\ee^{-P(x)}$，再按链式法则求导得 $y''=-V_0\ee^{-P(x)}P'(x)=-p(x)y'$。因此即使 $p$ 没有初等原函数，积分式仍完整满足原方程和两个初值。

\begin{figure}[H]\centering
\begin{tikzpicture}[>=Stealth,line cap=round,line join=round]
\begin{scope}[x=1.00000cm,y=1.00000cm]\draw[axis](0,0)--(3.25,0)node[right]{$x$};\draw[axis](0,-4)--(0,0.3)node[above]{$\ln|y/y_0|$};\draw[curve,domain=0:3,samples=101]plot(\x,{-.5*\x-.25*(\x)^2});\end{scope}\node[right]at(3.5,-1){局部斜率$=-(a+bx)$};
\end{tikzpicture}
\caption{对数幅值的斜率由位置相关系数决定。}\label{fig:c2-e4-s}
\end{figure}
\end{solution}
\subsection{常系数齐次方程的三种根}
对$y^{\prime\prime}+ay^{\prime}+by=0$，设$y=\ee^{\lambda x}$，得到特征方程$\lambda^2+a\lambda+b=0$。若两根不等，通解为$C_1\ee^{\lambda_1x}+C_2\ee^{\lambda_2x}$。重根$\lambda=-a/2$时，第二个独立解不能仍取相同指数，而应为$x\ee^{\lambda x}$，所以
\[y=(C_1+C_2x)\ee^{-ax/2}.
\]
若$\lambda=-a/2\pm i\omega$，$\omega=\sqrt{b-a^2/4}>0$，实通解为
\[y=\ee^{-ax/2}(C_1\cos\omega x+C_2\sin\omega x).
\]
复指数形式与实三角形式完全等价，常数须由初态确定；当要求$y$实值时，复指数的两个系数互为共轭。不同根型由判别式$a^2-4b$决定，不能把根是否为实数与解是否为实函数混为一谈。

\Needspace{8\baselineskip}
\subsection{特征根与独立解的来源}
把指数$\ee^{\lambda t}$代入常系数线性方程，所有导数只产生$\lambda$的幂，因而微分方程化成特征多项式。这种简化依赖系数恒定；系数随$t$变化时，通常不能继续用同一个指数试探全部解。

若二阶算子可因式分解为$(D-\lambda_1)(D-\lambda_2)$，可先设$z=y'-\lambda_2y$，原方程成为$z'-\lambda_1z=0$。乘$\ee^{-\lambda_1t}$后，$(\ee^{-\lambda_1t}z)'=0$，故$z=A\ee^{\lambda_1t}$。再解关于$y$的一阶方程：
\[(\ee^{-\lambda_2t}y)'=A\ee^{(\lambda_1-\lambda_2)t}.\]
两根不同时，右侧的原函数为$A\ee^{(\lambda_1-\lambda_2)t}/(\lambda_1-\lambda_2)$，因此
\[y=\frac{A}{\lambda_1-\lambda_2}\ee^{\lambda_1t}+B\ee^{\lambda_2t}.\]
将$A/(\lambda_1-\lambda_2)$、$B$重记为两个独立常数，便得到两指数通解。两根相同为$\lambda$时，上式改为$(\ee^{-\lambda t}y)'=A$，积分得到$\ee^{-\lambda t}y=At+B$，即$y=(At+B)\ee^{\lambda t}$。重根第二解的时间因子来自一次独立积分。

也可从两根靠拢的极限看出重根结构：
\[\frac{\ee^{(\lambda+\epsilon)t}-\ee^{\lambda t}}{\epsilon}
\longrightarrow t\ee^{\lambda t}.\]
直接令两个指数相同会把原来的两维解空间错误压成一维，取这个差商则保留了第二个独立方向。

当根为$\sigma\pm i\omega$时，由Euler公式得到$\ee^{\sigma t}\cos\omega t$和$\ee^{\sigma t}\sin\omega t$两个实独立解。实部决定包络增长或衰减，虚部决定振荡频率。对阻尼振子，根为$-\gamma\pm\sqrt{\gamma^2-\omega_0^2}$。因此有无振荡取决于判别式，而能否从某一侧穿越零点还取决于两种模态的系数，也就是初值。

初值应在通解上代入求出两个常数。若写成$\ee^{-\gamma t}(A\cos\omega_dt+B\sin\omega_dt)$，零时刻导数为$-\gamma A+\omega_dB$，其中包络的导数不可遗漏。下面统一用同一初态比较三类阻尼，使根型变化与初态变化分开。

\begin{example}[同一初态下的三种阻尼]\label{ex:c2-e5}
振子满足$\ddot x+2\gamma\dot x+\omega_0^2x=0$，其中$\gamma\ge0$、$\omega_0>0$，初态$x(0)=0,\dot x(0)=v_0$。分别求$\gamma<\omega_0$、$\gamma=\omega_0$和$\gamma>\omega_0$时的解；包括无阻尼情形。见图~\ref{fig:c2-e5-q}。
\end{example}

\begin{figure}[H]\centering
\begin{tikzpicture}[>=Stealth,line cap=round,line join=round]
\draw[thick](0,-.05)--(0,2.1);
\draw[thick](0,1.4)--(.3,1.4)--(.45,1.65)--(.7,1.15)--(.95,1.65)--(1.2,1.15)--(1.45,1.65)--(1.7,1.15)--(1.85,1.4)--(2.5,1.4);
\draw(2.5,.9)rectangle(3.4,1.95);\node at(2.95,1.4){$m$};
\draw(0,.35)--(.75,.35);\draw(.75,.1)rectangle(1.45,.6);\draw[thick](1.05,.16)--(1.05,.54);
\draw(1.05,.35)--(2.95,.35)--(2.95,.9);
\node[above]at(1.1,1.8){$k=m\omega_0^2$};\node[below]at(1.15,0){$b=2m\gamma$};
\draw[axis](3.4,1.4)--(4.7,1.4)node[right]{$v_0$};
\draw[axis](2.5,-.7)--(4.9,-.7)node[right]{$x$};
\end{tikzpicture}
\caption{弹簧与黏性阻尼器并联，阻尼力为$-b\dot x$；图示初态$x=0,\dot x=v_0$。}\label{fig:c2-e5-q}
\end{figure}
\begin{solution}
三种阻尼使用相同初态$x(0)=0,\dot x(0)=v_0$，图~\ref{fig:c2-e5-q}中的阻尼器系数为$2m\gamma$。按特征根分别写出两个独立自由解，再用这两个初值确定常数；临界阻尼须采用重根通解。


特征方程为 $s^2+2\gamma s+\omega_0^2=0$，根为 $s=-\gamma\pm\sqrt{\gamma^2-\omega_0^2}$。欠阻尼时，记 $\omega_d=\sqrt{\omega_0^2-\gamma^2}$，实通解为 $x=\ee^{-\gamma t}(A\cos\omega_dt+B\sin\omega_dt)$。对乘积求导，
\[
\dot x=\ee^{-\gamma t}\left[-\gamma(A\cos\omega_dt+B\sin\omega_dt)
-A\omega_d\sin\omega_dt+B\omega_d\cos\omega_dt\right].
\]
由 $x(0)=A=0$、$\dot x(0)=-\gamma A+\omega_dB=v_0$，得 $B=v_0/\omega_d$，因此
\[\boxed{x(t)=\frac{v_0}{\omega_d}\ee^{-\gamma t}\sin\omega_dt}.
\]
当$\gamma=0$时退化为$x=(v_0/\omega_0)\sin\omega_0t$，这也可由两个纯虚特征根直接得到。

临界阻尼为重根 $-\omega_0$，通解 $(A+Bt)\ee^{-\omega_0t}$。乘积求导为 $\dot x=[B-\omega_0(A+Bt)]\ee^{-\omega_0t}$，初态给出 $A=0$、$B-\omega_0A=v_0$，故
\[\boxed{x(t)=v_0t\ee^{-\omega_0t}}.
\]
过阻尼时记 $\delta=\sqrt{\gamma^2-\omega_0^2}$，两根为 $-\gamma\pm\delta$。通解写成 $x=C_+\ee^{(-\gamma+\delta)t}+C_-\ee^{(-\gamma-\delta)t}$，代入两个初值，
\[
C_++C_-=0,\qquad
(-\gamma+\delta)C_++(-\gamma-\delta)C_-=v_0.
\]
先用 $C_-=-C_+$，第二式化为 $2\delta C_+=v_0$，所以 $C_+=v_0/(2\delta)$、$C_-=-v_0/(2\delta)$，得到
\[\boxed{x(t)=\frac{v_0}{2\delta}
\left[\ee^{(-\gamma+\delta)t}-\ee^{(-\gamma-\delta)t}\right]
=\frac{v_0}{\delta}\ee^{-\gamma t}\sinh\delta t}.
\]
因为$\delta<\gamma$，过阻尼式两项都随时间衰减。图~\ref{fig:c2-e5-s}用相同初始斜率画出三条曲线：欠阻尼允许穿过零点，临界和过阻尼在$v_0>0$时从右侧趋近原点。


两个根均为负，但较慢模态的衰减率是 $\gamma-\delta$。有理化得 $\gamma-\delta=\omega_0^2/(\gamma+\delta)$，在 $\gamma\gg\omega_0$ 时约为 $\omega_0^2/(2\gamma)$，并不随阻尼增大而一直加快。

欠阻尼第一正零点为$\pi/\omega_d$；临界和过阻尼在$v_0>0,t>0$时均为正。图~\ref{fig:c2-e5-s}因此具有共同初始斜率，却有不同的穿越性质。两个极限中$\sin(\omega_dt)/\omega_d$和$\sinh(\delta t)/\delta$都趋于$t$，保留了重根的第二独立解。

\begin{figure}[H]\centering
\begin{tikzpicture}[>=Stealth,line cap=round,line join=round]
\draw[axis](0,0)--(6.2,0)node[right]{$\omega_0t$};\draw[axis](0,-.6)--(0,1.2)node[above]{$\omega_0x/v_0$};\draw[curve]plot coordinates{(0.0000000,0.0000000) (0.0300000,0.0298162) (0.0600000,0.0592502) (0.0900000,0.0882800) (0.1200000,0.1168845) (0.1500000,0.1450434) (0.1800000,0.1727365) (0.2100000,0.1999448) (0.2400000,0.2266498) (0.2700000,0.2528334) (0.3000000,0.2784785) (0.3300000,0.3035686) (0.3600000,0.3280878) (0.3900000,0.3520211) (0.4200000,0.3753540) (0.4500000,0.3980728) (0.4800000,0.4201646) (0.5100000,0.4416169) (0.5400000,0.4624183) (0.5700000,0.4825579) (0.6000000,0.5020256) (0.6300000,0.5208120) (0.6600000,0.5389084) (0.6900000,0.5563069) (0.7200000,0.5730002) (0.7500000,0.5889818) (0.7800000,0.6042459) (0.8100000,0.6187874) (0.8400000,0.6326019) (0.8700000,0.6456857) (0.9000000,0.6580358) (0.9300000,0.6696500) (0.9600000,0.6805266) (0.9900000,0.6906647) (1.0200000,0.7000640) (1.0500000,0.7087250) (1.0800000,0.7166486) (1.1100000,0.7238367) (1.1400000,0.7302915) (1.1700000,0.7360160) (1.2000000,0.7410138) (1.2300000,0.7452892) (1.2600000,0.7488468) (1.2900000,0.7516922) (1.3200000,0.7538311) (1.3500000,0.7552702) (1.3800000,0.7560165) (1.4100000,0.7560776) (1.4400000,0.7554616) (1.4700000,0.7541772) (1.5000000,0.7522334) (1.5300000,0.7496399) (1.5600000,0.7464068) (1.5900000,0.7425445) (1.6200000,0.7380640) (1.6500000,0.7329768) (1.6800000,0.7272945) (1.7100000,0.7210294) (1.7400000,0.7141941) (1.7700000,0.7068014) (1.8000000,0.6988646) (1.8300000,0.6903973) (1.8600000,0.6814134) (1.8900000,0.6719271) (1.9200000,0.6619529) (1.9500000,0.6515055) (1.9800000,0.6405999) (2.0100000,0.6292514) (2.0400000,0.6174753) (2.0700000,0.6052873) (2.1000000,0.5927032) (2.1300000,0.5797391) (2.1600000,0.5664110) (2.1900000,0.5527351) (2.2200000,0.5387280) (2.2500000,0.5244060) (2.2800000,0.5097858) (2.3100000,0.4948839) (2.3400000,0.4797170) (2.3700000,0.4643019) (2.4000000,0.4486554) (2.4300000,0.4327942) (2.4600000,0.4167350) (2.4900000,0.4004945) (2.5200000,0.3840896) (2.5500000,0.3675367) (2.5800000,0.3508525) (2.6100000,0.3340534) (2.6400000,0.3171559) (2.6700000,0.3001762) (2.7000000,0.2831306) (2.7300000,0.2660350) (2.7600000,0.2489053) (2.7900000,0.2317573) (2.8200000,0.2146066) (2.8500000,0.1974684) (2.8800000,0.1803580) (2.9100000,0.1632904) (2.9400000,0.1462803) (2.9700000,0.1293423) (3.0000000,0.1124905) (3.0300000,0.0957392) (3.0600000,0.0791020) (3.0900000,0.0625925) (3.1200000,0.0462240) (3.1500000,0.0300094) (3.1800000,0.0139613) (3.2100000,-0.0019078) (3.2400000,-0.0175860) (3.2700000,-0.0330614) (3.3000000,-0.0483226) (3.3300000,-0.0633585) (3.3600000,-0.0781585) (3.3900000,-0.0927119) (3.4200000,-0.1070089) (3.4500000,-0.1210395) (3.4800000,-0.1347946) (3.5100000,-0.1482650) (3.5400000,-0.1614420) (3.5700000,-0.1743175) (3.6000000,-0.1868834) (3.6300000,-0.1991323) (3.6600000,-0.2110570) (3.6900000,-0.2226506) (3.7200000,-0.2339067) (3.7500000,-0.2448193) (3.7800000,-0.2553828) (3.8100000,-0.2655918) (3.8400000,-0.2754414) (3.8700000,-0.2849271) (3.9000000,-0.2940449) (3.9300000,-0.3027908) (3.9600000,-0.3111615) (3.9900000,-0.3191540) (4.0200000,-0.3267657) (4.0500000,-0.3339943) (4.0800000,-0.3408379) (4.1100000,-0.3472950) (4.1400000,-0.3533643) (4.1700000,-0.3590452) (4.2000000,-0.3643371) (4.2300000,-0.3692400) (4.2600000,-0.3737540) (4.2900000,-0.3778799) (4.3200000,-0.3816186) (4.3500000,-0.3849712) (4.3800000,-0.3879395) (4.4100000,-0.3905254) (4.4400000,-0.3927311) (4.4700000,-0.3945591) (4.5000000,-0.3960124) (4.5300000,-0.3970941) (4.5600000,-0.3978076) (4.5900000,-0.3981568) (4.6200000,-0.3981457) (4.6500000,-0.3977785) (4.6800000,-0.3970599) (4.7100000,-0.3959946) (4.7400000,-0.3945879) (4.7700000,-0.3928449) (4.8000000,-0.3907713) (4.8300000,-0.3883728) (4.8600000,-0.3856556) (4.8900000,-0.3826257) (4.9200000,-0.3792898) (4.9500000,-0.3756543) (4.9800000,-0.3717261) (5.0100000,-0.3675123) (5.0400000,-0.3630200) (5.0700000,-0.3582565) (5.1000000,-0.3532294) (5.1300000,-0.3479462) (5.1600000,-0.3424148) (5.1900000,-0.3366431) (5.2200000,-0.3306391) (5.2500000,-0.3244109) (5.2800000,-0.3179668) (5.3100000,-0.3113152) (5.3400000,-0.3044643) (5.3700000,-0.2974229) (5.4000000,-0.2901994) (5.4300000,-0.2828024) (5.4600000,-0.2752407) (5.4900000,-0.2675230) (5.5200000,-0.2596580) (5.5500000,-0.2516546) (5.5800000,-0.2435215) (5.6100000,-0.2352676) (5.6400000,-0.2269017) (5.6700000,-0.2184326) (5.7000000,-0.2098691) (5.7300000,-0.2012200) (5.7600000,-0.1924941) (5.7900000,-0.1837001) (5.8200000,-0.1748466) (5.8500000,-0.1659424) (5.8800000,-0.1569959) (5.9100000,-0.1480157) (5.9400000,-0.1390102) (5.9700000,-0.1299877) (6.0000000,-0.1209566)};\draw[dashed]plot coordinates{(0.0000000,0.0000000) (0.0300000,0.0291134) (0.0600000,0.0565059) (0.0900000,0.0822538) (0.1200000,0.1064305) (0.1500000,0.1291062) (0.1800000,0.1503486) (0.2100000,0.1702227) (0.2400000,0.1887907) (0.2700000,0.2061125) (0.3000000,0.2222455) (0.3300000,0.2372448) (0.3600000,0.2511635) (0.3900000,0.2640522) (0.4200000,0.2759597) (0.4500000,0.2869327) (0.4800000,0.2970160) (0.5100000,0.3062527) (0.5400000,0.3146841) (0.5700000,0.3223495) (0.6000000,0.3292870) (0.6300000,0.3355328) (0.6600000,0.3411219) (0.6900000,0.3460875) (0.7200000,0.3504616) (0.7500000,0.3542749) (0.7800000,0.3575567) (0.8100000,0.3603350) (0.8400000,0.3626368) (0.8700000,0.3644878) (0.9000000,0.3659127) (0.9300000,0.3669350) (0.9600000,0.3675772) (0.9900000,0.3678609) (1.0200000,0.3678068) (1.0500000,0.3674346) (1.0800000,0.3667632) (1.1100000,0.3658104) (1.1400000,0.3645937) (1.1700000,0.3631293) (1.2000000,0.3614331) (1.2300000,0.3595199) (1.2600000,0.3574041) (1.2900000,0.3550993) (1.3200000,0.3526186) (1.3500000,0.3499744) (1.3800000,0.3471784) (1.4100000,0.3442420) (1.4400000,0.3411760) (1.4700000,0.3379905) (1.5000000,0.3346952) (1.5300000,0.3312996) (1.5600000,0.3278123) (1.5900000,0.3242417) (1.6200000,0.3205959) (1.6500000,0.3168823) (1.6800000,0.3131083) (1.7100000,0.3092805) (1.7400000,0.3054055) (1.7700000,0.3014894) (1.8000000,0.2975380) (1.8300000,0.2935568) (1.8600000,0.2895511) (1.8900000,0.2855257) (1.9200000,0.2814854) (1.9500000,0.2774344) (1.9800000,0.2733771) (2.0100000,0.2693172) (2.0400000,0.2652586) (2.0700000,0.2612046) (2.1000000,0.2571585) (2.1300000,0.2531234) (2.1600000,0.2491023) (2.1900000,0.2450977) (2.2200000,0.2411122) (2.2500000,0.2371483) (2.2800000,0.2332080) (2.3100000,0.2292935) (2.3400000,0.2254067) (2.3700000,0.2215493) (2.4000000,0.2177231) (2.4300000,0.2139295) (2.4600000,0.2101700) (2.4900000,0.2064458) (2.5200000,0.2027582) (2.5500000,0.1991082) (2.5800000,0.1954969) (2.6100000,0.1919252) (2.6400000,0.1883938) (2.6700000,0.1849034) (2.7000000,0.1814549) (2.7300000,0.1780487) (2.7600000,0.1746853) (2.7900000,0.1713652) (2.8200000,0.1680888) (2.8500000,0.1648563) (2.8800000,0.1616681) (2.9100000,0.1585244) (2.9400000,0.1554252) (2.9700000,0.1523708) (3.0000000,0.1493612) (3.0300000,0.1463964) (3.0600000,0.1434763) (3.0900000,0.1406010) (3.1200000,0.1377704) (3.1500000,0.1349842) (3.1800000,0.1322424) (3.2100000,0.1295447) (3.2400000,0.1268910) (3.2700000,0.1242810) (3.3000000,0.1217145) (3.3300000,0.1191910) (3.3600000,0.1167105) (3.3900000,0.1142724) (3.4200000,0.1118765) (3.4500000,0.1095224) (3.4800000,0.1072098) (3.5100000,0.1049382) (3.5400000,0.1027072) (3.5700000,0.1005164) (3.6000000,0.0983654) (3.6300000,0.0962537) (3.6600000,0.0941810) (3.6900000,0.0921467) (3.7200000,0.0901504) (3.7500000,0.0881915) (3.7800000,0.0862698) (3.8100000,0.0843846) (3.8400000,0.0825354) (3.8700000,0.0807219) (3.9000000,0.0789435) (3.9300000,0.0771996) (3.9600000,0.0754899) (3.9900000,0.0738139) (4.0200000,0.0721709) (4.0500000,0.0705606) (4.0800000,0.0689825) (4.1100000,0.0674360) (4.1400000,0.0659206) (4.1700000,0.0644359) (4.2000000,0.0629814) (4.2300000,0.0615566) (4.2600000,0.0601610) (4.2900000,0.0587941) (4.3200000,0.0574555) (4.3500000,0.0561446) (4.3800000,0.0548611) (4.4100000,0.0536043) (4.4400000,0.0523740) (4.4700000,0.0511695) (4.5000000,0.0499905) (4.5300000,0.0488365) (4.5600000,0.0477070) (4.5900000,0.0466016) (4.6200000,0.0455199) (4.6500000,0.0444614) (4.6800000,0.0434258) (4.7100000,0.0424125) (4.7400000,0.0414212) (4.7700000,0.0404514) (4.8000000,0.0395028) (4.8300000,0.0385749) (4.8600000,0.0376674) (4.8900000,0.0367798) (4.9200000,0.0359117) (4.9500000,0.0350629) (4.9800000,0.0342328) (5.0100000,0.0334212) (5.0400000,0.0326277) (5.0700000,0.0318519) (5.1000000,0.0310934) (5.1300000,0.0303520) (5.1600000,0.0296272) (5.1900000,0.0289187) (5.2200000,0.0282263) (5.2500000,0.0275495) (5.2800000,0.0268880) (5.3100000,0.0262416) (5.3400000,0.0256099) (5.3700000,0.0249927) (5.4000000,0.0243895) (5.4300000,0.0238002) (5.4600000,0.0232244) (5.4900000,0.0226619) (5.5200000,0.0221123) (5.5500000,0.0215754) (5.5800000,0.0210509) (5.6100000,0.0205386) (5.6400000,0.0200382) (5.6700000,0.0195494) (5.7000000,0.0190720) (5.7300000,0.0186058) (5.7600000,0.0181504) (5.7900000,0.0177057) (5.8200000,0.0172715) (5.8500000,0.0168474) (5.8800000,0.0164333) (5.9100000,0.0160290) (5.9400000,0.0156343) (5.9700000,0.0152488) (6.0000000,0.0148725)};\draw[dotted]plot coordinates{(0.0000000,0.0000000) (0.0300000,0.0286853) (0.0600000,0.0548770) (0.0900000,0.0787672) (0.1200000,0.1005334) (0.1500000,0.1203397) (0.1800000,0.1383377) (0.2100000,0.1546676) (0.2400000,0.1694589) (0.2700000,0.1828312) (0.3000000,0.1948953) (0.3300000,0.2057533) (0.3600000,0.2154997) (0.3900000,0.2242218) (0.4200000,0.2320002) (0.4500000,0.2389095) (0.4800000,0.2450185) (0.5100000,0.2503908) (0.5400000,0.2550852) (0.5700000,0.2591561) (0.6000000,0.2626535) (0.6300000,0.2656239) (0.6600000,0.2681100) (0.6900000,0.2701514) (0.7200000,0.2717846) (0.7500000,0.2730434) (0.7800000,0.2739588) (0.8100000,0.2745596) (0.8400000,0.2748724) (0.8700000,0.2749218) (0.9000000,0.2747303) (0.9300000,0.2743188) (0.9600000,0.2737067) (0.9900000,0.2729117) (1.0200000,0.2719502) (1.0500000,0.2708375) (1.0800000,0.2695874) (1.1100000,0.2682129) (1.1400000,0.2667259) (1.1700000,0.2651373) (1.2000000,0.2634573) (1.2300000,0.2616951) (1.2600000,0.2598594) (1.2900000,0.2579580) (1.3200000,0.2559983) (1.3500000,0.2539869) (1.3800000,0.2519301) (1.4100000,0.2498335) (1.4400000,0.2477023) (1.4700000,0.2455413) (1.5000000,0.2433551) (1.5300000,0.2411475) (1.5600000,0.2389225) (1.5900000,0.2366833) (1.6200000,0.2344331) (1.6500000,0.2321749) (1.6800000,0.2299112) (1.7100000,0.2276444) (1.7400000,0.2253769) (1.7700000,0.2231105) (1.8000000,0.2208472) (1.8300000,0.2185887) (1.8600000,0.2163364) (1.8900000,0.2140917) (1.9200000,0.2118560) (1.9500000,0.2096304) (1.9800000,0.2074160) (2.0100000,0.2052136) (2.0400000,0.2030242) (2.0700000,0.2008485) (2.1000000,0.1986872) (2.1300000,0.1965409) (2.1600000,0.1944102) (2.1900000,0.1922956) (2.2200000,0.1901975) (2.2500000,0.1881162) (2.2800000,0.1860523) (2.3100000,0.1840058) (2.3400000,0.1819771) (2.3700000,0.1799664) (2.4000000,0.1779739) (2.4300000,0.1759997) (2.4600000,0.1740439) (2.4900000,0.1721067) (2.5200000,0.1701881) (2.5500000,0.1682882) (2.5800000,0.1664070) (2.6100000,0.1645444) (2.6400000,0.1627006) (2.6700000,0.1608754) (2.7000000,0.1590689) (2.7300000,0.1572810) (2.7600000,0.1555115) (2.7900000,0.1537606) (2.8200000,0.1520280) (2.8500000,0.1503137) (2.8800000,0.1486176) (2.9100000,0.1469395) (2.9400000,0.1452794) (2.9700000,0.1436372) (3.0000000,0.1420127) (3.0300000,0.1404058) (3.0600000,0.1388164) (3.0900000,0.1372443) (3.1200000,0.1356894) (3.1500000,0.1341515) (3.1800000,0.1326306) (3.2100000,0.1311264) (3.2400000,0.1296388) (3.2700000,0.1281677) (3.3000000,0.1267129) (3.3300000,0.1252743) (3.3600000,0.1238517) (3.3900000,0.1224449) (3.4200000,0.1210538) (3.4500000,0.1196783) (3.4800000,0.1183181) (3.5100000,0.1169732) (3.5400000,0.1156434) (3.5700000,0.1143285) (3.6000000,0.1130284) (3.6300000,0.1117430) (3.6600000,0.1104719) (3.6900000,0.1092153) (3.7200000,0.1079727) (3.7500000,0.1067442) (3.7800000,0.1055296) (3.8100000,0.1043287) (3.8400000,0.1031413) (3.8700000,0.1019674) (3.9000000,0.1008068) (3.9300000,0.0996593) (3.9600000,0.0985248) (3.9900000,0.0974031) (4.0200000,0.0962942) (4.0500000,0.0951978) (4.0800000,0.0941139) (4.1100000,0.0930423) (4.1400000,0.0919828) (4.1700000,0.0909354) (4.2000000,0.0898998) (4.2300000,0.0888760) (4.2600000,0.0878638) (4.2900000,0.0868632) (4.3200000,0.0858738) (4.3500000,0.0848958) (4.3800000,0.0839288) (4.4100000,0.0829729) (4.4400000,0.0820278) (4.4700000,0.0810935) (4.5000000,0.0801698) (4.5300000,0.0792566) (4.5600000,0.0783538) (4.5900000,0.0774612) (4.6200000,0.0765788) (4.6500000,0.0757065) (4.6800000,0.0748441) (4.7100000,0.0739915) (4.7400000,0.0731486) (4.7700000,0.0723153) (4.8000000,0.0714914) (4.8300000,0.0706770) (4.8600000,0.0698718) (4.8900000,0.0690758) (4.9200000,0.0682889) (4.9500000,0.0675109) (4.9800000,0.0667418) (5.0100000,0.0659814) (5.0400000,0.0652297) (5.0700000,0.0644866) (5.1000000,0.0637519) (5.1300000,0.0630256) (5.1600000,0.0623075) (5.1900000,0.0615976) (5.2200000,0.0608959) (5.2500000,0.0602021) (5.2800000,0.0595162) (5.3100000,0.0588381) (5.3400000,0.0581678) (5.3700000,0.0575050) (5.4000000,0.0568499) (5.4300000,0.0562022) (5.4600000,0.0555619) (5.4900000,0.0549288) (5.5200000,0.0543030) (5.5500000,0.0536843) (5.5800000,0.0530727) (5.6100000,0.0524680) (5.6400000,0.0518702) (5.6700000,0.0512792) (5.7000000,0.0506950) (5.7300000,0.0501174) (5.7600000,0.0495464) (5.7900000,0.0489819) (5.8200000,0.0484238) (5.8500000,0.0478721) (5.8800000,0.0473267) (5.9100000,0.0467875) (5.9400000,0.0462544) (5.9700000,0.0457274) (6.0000000,0.0452064)};
\end{tikzpicture}
\caption{实线、虚线、点线分别取$\gamma/\omega_0=0.2,1,1.5$，初始斜率相同。}\label{fig:c2-e5-s}
\end{figure}
\end{solution}

\section{非齐次方程与外力响应}\label{sec:c2-forcing}
外力决定特解的形式，初态决定齐次解的系数。先求响应结构，再用初始条件连接两部分，才能区分短时建立过程与长期稳态。

图~\ref{fig:c2-characteristic-roots}将三类根画在复平面上。根的实部决定指数衰减，虚部决定振荡；重复实根还须配上$t\ee^{rt}$才能补足第二个独立解。
\begin{figure}[H]\centering
\begin{tikzpicture}[>=Stealth,line cap=round,line join=round]
\begin{scope}
\draw[axis](-2,0)--(.5,0)node[right]{$\Re r$};\draw[axis](0,-1.15)--(0,1.25)node[above]{$\Im r$};
\fill(-1.7,0)circle(1.8pt)node[below]{$r_-$};\fill(-.6,0)circle(1.8pt)node[below]{$r_+$};\node at(-.8,-1.65){两个负实根};
\end{scope}
\begin{scope}[shift={(3.7,0)}]
\draw[axis](-2,0)--(.5,0)node[right]{$\Re r$};\draw[axis](0,-1.15)--(0,1.25)node[above]{$\Im r$};
\draw[thick](-1.2,0)circle(.10);\fill(-1.2,0)circle(1pt)node[below=4pt]{$-\omega_0$};\node at(-.8,-1.65){重负实根};
\end{scope}
\begin{scope}[shift={(7.4,0)}]
\draw[axis](-2,0)--(.5,0)node[right]{$\Re r$};\draw[axis](0,-1.15)--(0,1.25)node[above]{$\Im r$};
\draw[guide](-1,-.85)--(-1,.85);\fill(-1,.85)circle(1.8pt)node[left]{$-\gamma+\ii\omega_d$};\fill(-1,-.85)circle(1.8pt)node[left]{$-\gamma-\ii\omega_d$};\node at(-.8,-1.65){共轭复根};
\end{scope}
\end{tikzpicture}
\caption{阻尼振子的特征根：过阻尼、临界阻尼与欠阻尼。}\label{fig:c2-characteristic-roots}
\end{figure}

\subsection{待定系数与共振重数}
设$\mathcal L=D^2+aD+b$，$D=\dd/\dd x$。若外源为$P_n(x)\ee^{\eta x}$，令$y=\ee^{\eta x}Q(x)$。连续两次使用乘积法则，得到
\[
 y'=\ee^{\eta x}(Q'+\eta Q),\qquad
 y''=\ee^{\eta x}(Q''+2\eta Q'+\eta^2Q).
\]
将这两式代回算子，并按$Q''$、$Q'$、$Q$收集系数，得到
\begin{equation}\mathcal L[\ee^{\eta x}Q]=\ee^{\eta x}
\left[Q''+(2\eta+a)Q'+(\eta^2+a\eta+b)Q\right].
\end{equation}
当$\eta$不是特征根时，可取$Q$为$n$次多项式；当它是单根或重根时，微分算子分别消去一个或两个最低阶系数，需要把试探特解乘$x$或$x^2$。统一写成
\[y_p=x^s\ee^{\eta x}Q_n(x),\qquad
s=\eta\text{作为特征根的重数，取 }0,1,2.
\]
若外源是$\ee^{\eta x}[P_l(x)\cos\omega x+P_n(x)\sin\omega x]$，试探式应同时含正弦与余弦，两个多项式都取次数$\max(l,n)$，再乘$x^s$；这里$s$由复数$\eta+i\omega$的特征根重数决定。即使外源只有余弦，求导仍会产生正弦，通常不能只猜一个余弦项。

待定系数适用于这类有限维函数族。一般$f(x)$可用常数变易法：若$y_1,y_2$是两个齐次独立解，$W=y_1y_2'-y_1'y_2\ne0$，令$y_p=u_1y_1+u_2y_2$并附加$u_1'y_1+u_2'y_2=0$。于是
\[y_p'=u_1y_1'+u_2y_2',\qquad
y_p''=u_1'y_1'+u_2'y_2'+u_1y_1''+u_2y_2''.\]
代入方程时，含$u_j$的项因齐次方程相消，留下
\[u_1'y_1+u_2'y_2=0,\qquad u_1'y_1'+u_2'y_2'=f.\]
前式乘$y_2'$、后式乘$y_2$再相减，可解$u_1'$；消去$u_1'$同理解$u_2'$，得
\[u_1'=-\frac{y_2f}{W},\qquad u_2'=\frac{y_1f}{W}.
\]
积分常数并入齐次解，不能重复计算自由常数。外源为脉冲或分段函数时，按时间边界拼接往往比把每一段常数全部重写更清楚。

\Needspace{8\baselineskip}
\subsection{共振外源与试探式的重数}
待定系数法的实质是寻找一个在求导后仍留在同一有限维函数族中的试探式。多项式求导后次数降低，指数求导后保留指数，正弦余弦求导后互相转换，所以这些外源可用有限个未知系数求特解。任意函数外源一般不具有这种闭合性，需用积分响应或常数变易法。

代换$y=\ee^{\eta x}Q$使算子中的$D$整体替成$D+\eta$。若$\eta$不是特征根，常数项系数不为零，$n$次多项式外源可由$n$次试探多项式匹配。若是单根，算子至少含一个$D$因子，次数会降低一次，便须让$Q$多含一个$x$；若为二重根，则降低两次，须多含$x^2$。加入的低次项若恰为齐次解，其系数可以并入齐次常数，所以标准试探式写成$x^sQ_n(x)$即可。

三角外源可先看作$\Re[\ee^{(\eta+i\omega)x}P(x)]$。重合与否由复数$\eta+i\omega$判断；最后取实部恢复实解。即使外源只有余弦，试探式也通常要同时含正弦，因为阻尼项的一次导数会使二者耦合。

特解只需找到一个。加上任意齐次解仍是另一个特解，但不增加初值问题的自由度。计算时先求任一特解，再把两个自由常数留给齐次部分；若在特解里重复加入任意常数，便会使同一解被多种参数表示，妨碍初值确定。

\begin{example}[特解何时需要乘自变量]\label{ex:c2-e6}
分别构造下列方程的一个特解，并写出通解：(1)$y^{\prime\prime}-2y^{\prime}+y=\ee^x$；(2)$y^{\prime\prime}+y=\cos x$；(3)$y^{\prime\prime}+y=x^2\ee^x$。说明三式中指数或复指数与特征根的关系。 见图~\ref{fig:c2-e6-q}。
\end{example}

\begin{figure}[H]\centering
\begin{tikzpicture}[>=Stealth,line cap=round,line join=round]
\draw[axis](-1.5,0)--(2,0)node[right]{$\Re\lambda$};\draw[axis](0,-1.6)--(0,1.8)node[above]{$\Im\lambda$};\fill(1,0)circle(2pt)node[above]{$1$：重根};\fill(0,1)circle(1.5pt)node[left]{$i$};\fill(0,-1)circle(1.5pt)node[left]{$-i$};\node[right,align=left]at(2.7,.7){第一式：$\eta=1$\\[6pt]第二式：$\eta+i\omega=i$\\[6pt]第三式：$\eta=1\ne\pm i$};
\end{tikzpicture}
\caption{用外源对应的复指数与各方程的特征根比较，重数决定试探式的补因子。}\label{fig:c2-e6-q}
\end{figure}
\begin{solution}
先将外源对应的指数与特征根比较，再选试探函数。图~\ref{fig:c2-e6-q}中的第一式是实指数二重重合，第二式是复指数单根重合，第三式则没有重合。三种情况需要乘的$x$次数不同，取错后会把试探特解变成齐次解。


第一式的特征多项式是 $(\lambda-1)^2$，外源指数 $\eta=1$ 为二重根。图~\ref{fig:c2-e6-s}中设 $y=\ee^xu$，乘积求导为 $y'=\ee^x(u+u')$、$y''=\ee^x(u+2u'+u'')$。代回左端，括号内 $u+2u'+u''-2(u+u')+u=u''$，所以 $u''=1$。依次积分得 $u'=x+C_2$、$u=x^2/2+C_2x+C_1$，通解为
\[y=\ee^x\left(C_1+C_2x+\frac{x^2}{2}\right).
\]
第二式的特征根为 $\pm i$，外源 $\cos x$ 对应 $i$，是单根共振。取完整试探 $y_p=x(A\cos x+B\sin x)$，逐次求导，
\[
\begin{aligned}
y_p'&=A\cos x+B\sin x+x(-A\sin x+B\cos x),\\
y_p''&=2(-A\sin x+B\cos x)-x(A\cos x+B\sin x).
\end{aligned}
\]
因此 $y_p''+y_p=-2A\sin x+2B\cos x$，匹配外源得 $A=0,B=1/2$，通解为
\[y=C_1\cos x+C_2\sin x+\frac x2\sin x.
\]
第三式中 $1$ 不是 $\lambda^2+1$ 的根，可令 $y_p=\ee^xQ_2(x)$。乘积求导为 $y_p'=\ee^x(Q_2+Q_2')$、$y_p''=\ee^x(Q_2+2Q_2'+Q_2'')$，约去指数因子后得 $Q_2''+2Q_2'+2Q_2=x^2$。设 $Q_2=Ax^2+Bx+C$，则 $Q_2'=2Ax+B$、$Q_2''=2A$；代入并收集幂次为
\[
2Ax^2+(4A+2B)x+(2A+2B+2C)=x^2.
\]
比较系数，依次得到 $A=1/2$、$B=-2A=-1$、$C=-A-B=1/2$，因此
\[y=C_1\cos x+C_2\sin x+\ee^x\left(\frac{x^2}{2}-x+\frac12\right).
\]
这里第一式若只乘一次$x$，所得函数仍是齐次解，代回后左边为零；第二式若只猜$A\cos x$也有同样问题。所需的$x$次幂由特征根的重数决定。


图~\ref{fig:c2-e6-s}中的算子降阶与系数匹配给出的只是特解；通解中的两项齐次模态仍要保留，以便以后施加任意两个初值。

\begin{figure}[H]\centering
\begin{tikzpicture}[>=Stealth,line cap=round,line join=round]
\node[draw,align=center]at(0,1){$y=\ee^xu$\\$(D-1)^2y=\ee^xu^{\prime\prime}$};\node[draw,align=center]at(4.7,1){$y_p=Ax\sin x$\\$(D^2+1)y_p=2A\cos x$};\draw[axis](0,.3)--(0,-.7)node[below]{$u=x^2/2$};\draw[axis](4.7,.3)--(4.7,-.7)node[below]{$A=1/2$};
\end{tikzpicture}
\caption{把试探特解代回算子，分别剩下常数与余弦外源。}\label{fig:c2-e6-s}
\end{figure}
\end{solution}
\subsection{从初始条件构造因果响应}
对$\ddot x+2\gamma\dot x+\omega_0^2x=f(t)$，定义单位初速响应$G(t)$：$t>0$时满足齐次方程，$G(0)=0,G'(0)=1$。三种阻尼分别有
\begin{equation}\label{eq:c2-G}
G(t)=\begin{cases}
\ee^{-\gamma t}\sin(\omega_dt)/\omega_d,&\gamma<\omega_0,\\
t\ee^{-\omega_0t},&\gamma=\omega_0,\\
\ee^{-\gamma t}\sinh(\delta t)/\delta,&\gamma>\omega_0.
\end{cases}
\end{equation}
取$G(t)=0$于$t<0$，表达响应只依赖过去的外源。初态$x(0)=0,\dot x(0)=v_0$时，
\begin{equation}\label{eq:c2-duhamel}
x(t)=v_0G(t)+\int_0^tG(t-s)f(s)\dd s.
\end{equation}
对卷积求一次导数，边界项含$G(0)=0$；求第二次导数时出现$G'(0)f(t)=f(t)$，其余积分满足齐次算子关系。因此式\eqref{eq:c2-duhamel}既满足方程，也满足两个初值。若$x(0)=x_0\ne0$，再加$x_0[G'(t)+2\gamma G(t)]$即可。

\Needspace{8\baselineskip}
\subsection{因果卷积与完整初值解}
把$G$定义为“零位移、单位初速度”的自由解，是为了把一次短促输入引起的响应单独求出。时刻$s$的一小段加速度外源$f(s)\dd s$，相当于让速度增加$f(s)\dd s$；此后经过$t-s$的响应为$G(t-s)f(s)\dd s$。线性方程允许把所有过去的输入相加，得到卷积积分。

这个物理构造可以逐步代回方程验证。令$z(t)=\int_0^tG(t-s)f(s)\dd s$，在$f$连续的区间内，
\[z'=\int_0^tG'(t-s)f(s)\dd s,\]
\[z''=f(t)+\int_0^tG''(t-s)f(s)\dd s.\]
第一次求导的端点项因$G(0)=0$消失，第二次则因$G'(0)=1$留下$f(t)$。将三项线性组合，积分中的$G''+2\gamma G'+\omega_0^2G$为零，恰好恢复外源。与此同时$z(0)=z'(0)=0$，所以卷积部分不改变指定初态。

对于非零初位移，需要另一个自由解$K(t)=G'(t)+2\gamma G(t)$。系数为常量，对$G$的齐次方程求导，$G'$仍满足同一齐次方程，所以$K$也是自由解。将$G(0)=0$、$G'(0)=1$代入原方程，得$G''(0^+)=-2\gamma$；因此
\[
 K(0)=G'(0)+2\gamma G(0)=1,\qquad
 K'(0^+)=G''(0^+)+2\gamma G'(0)=0.
\]
它恰好承载单位初位移而不引入初速度，故完整解为
\[x(t)=x_0K(t)+v_0G(t)+\int_0^tG(t-s)f(s)\dd s.\]
三个部分分别承载初位移、初速度和外源。矩形脉冲、延迟开启等分段外源都可直接放入积分限，而不必为每段另设一组互相混杂的常数。

当$\gamma>0$时，齐次项逐渐衰减，周期特解描述最终稳态。给定$t=0$的初值后，仍须加入齐次项，使初位移和初速度与题设相符。图~\ref{fig:c2-e7-s}的完整曲线与稳态曲线在短时不同，正是这两项修正的结果。

\begin{example}[简谐外力下的完整初值解]\label{ex:c2-e7}
质量$m$的振子满足$\ddot x+2\gamma\dot x+\omega_0^2x=(F_0/m)\cos\Omega t$，初始$x(0)=0,\dot x(0)=v_0$。给出适用于三种阻尼的积分解，并在欠阻尼时写成瞬态加稳态的显式形式。另讨论无阻尼共振。 见图~\ref{fig:c2-e7-q}。
\end{example}

\begin{figure}[H]\centering
\begin{tikzpicture}[>=Stealth,line cap=round,line join=round]
\draw[thick](0,-.05)--(0,2.1);
\draw[thick](0,1.4)--(.3,1.4)--(.45,1.65)--(.7,1.15)--(.95,1.65)--(1.2,1.15)--(1.45,1.65)--(1.7,1.15)--(1.85,1.4)--(2.5,1.4);
\draw(2.5,.9)rectangle(3.4,1.95);\node at(2.95,1.4){$m$};
\draw(0,.35)--(.75,.35);\draw(.75,.1)rectangle(1.45,.6);\draw[thick](1.05,.16)--(1.05,.54);
\draw(1.05,.35)--(2.95,.35)--(2.95,.9);
\node[above]at(1.1,1.8){$k=m\omega_0^2$};\node[below]at(1.15,0){$b=2m\gamma$};
\draw[axis](3.4,1.4)--(4.7,1.4)node[right]{$F_0\cos\Omega t$};
\draw[axis](2.5,-.7)--(4.9,-.7)node[right]{$x$};
\end{tikzpicture}
\caption{右端外力驱动弹簧与黏性阻尼器并联的振子，阻尼系数$b=2m\gamma$。}\label{fig:c2-e7-q}
\end{figure}
\begin{solution}
图~\ref{fig:c2-e7-q}中的驱动从$t=0$施加，初位移和初速度也须同时满足。
由式\eqref{eq:c2-duhamel}，三种阻尼共用的积分解为
\[
 \boxed{x(t)=v_0G(t)+\frac{F_0}{m}
 \int_0^tG(t-s)\cos\Omega s\dd s}.
\]
这一卷积满足原方程，可由变上限求导直接看出。记 $I(t)=\int_0^tG(t-s)\cos\Omega s\dd s$，用 $G(0)=0,G'(0)=1$，有
\[
I'=\int_0^tG'(t-s)\cos\Omega s\dd s,\qquad
I''=\cos\Omega t+\int_0^tG''(t-s)\cos\Omega s\dd s.
\]
将 $I''+2\gamma I'+\omega_0^2I$ 合并，积分核为 $G''+2\gamma G'+\omega_0^2G=0$，只剩 $\cos\Omega t$；同时 $I(0)=I'(0)=0$，所以初速度由 $v_0G$ 给出。

为写出显式解，先设 $x_p=C\cos\Omega t+D\sin\Omega t$。求导得 $\dot x_p=-C\Omega\sin\Omega t+D\Omega\cos\Omega t$、$\ddot x_p=-C\Omega^2\cos\Omega t-D\Omega^2\sin\Omega t$，代回方程比较系数：
\[
 \begin{cases}
 (\omega_0^2-\Omega^2)C+2\gamma\Omega D=F_0/m,\\
 (\omega_0^2-\Omega^2)D-2\gamma\Omega C=0.
 \end{cases}
\]
记 $\Delta=\omega_0^2-\Omega^2$、$\mathcal D=\Delta^2+4\gamma^2\Omega^2$。第一式乘 $\Delta$，减去第二式乘 $2\gamma\Omega$，消去 $D$，得 $\mathcal DC=(F_0/m)\Delta$；第一式乘 $2\gamma\Omega$，加上第二式乘 $\Delta$，消去 $C$，得 $\mathcal DD=(F_0/m)2\gamma\Omega$。在 $\mathcal D>0$ 时，
\[
 C=\frac{(F_0/m)(\omega_0^2-\Omega^2)}{\mathcal D},\qquad
 D=\frac{(F_0/m)2\gamma\Omega}{\mathcal D}.
\]
对$\gamma>0$，也可写$x_p=X\cos(\Omega t-\delta)$，其中
\[
 X=\frac{F_0/m}{\sqrt{\mathcal D}},\qquad
 \delta=\operatorname{atan2}(2\gamma\Omega,\omega_0^2-\Omega^2).
\]
两自变量反正切保留正确象限，尤其在$\Omega>\omega_0$时不能丢失相位的$\pi$差别。
\begin{figure}[H]\centering
\begin{tikzpicture}[>=Stealth,line cap=round,line join=round]
\begin{scope}[x=.4cm,y=1.2cm]\draw[axis](0,0)--(16.5,0)node[right]{$\tau$};\draw[axis](0,-1.2)--(0,1.5)node[above]{$x$};\draw[curve]plot coordinates{(0.0000000,0.0000000) (0.0666667,0.0273059) (0.1333333,0.0557784) (0.2000000,0.0852456) (0.2666667,0.1155306) (0.3333333,0.1464525) (0.4000000,0.1778278) (0.4666667,0.2094707) (0.5333333,0.2411946) (0.6000000,0.2728129) (0.6666667,0.3041398) (0.7333333,0.3349915) (0.8000000,0.3651869) (0.8666667,0.3945484) (0.9333333,0.4229031) (1.0000000,0.4500832) (1.0666667,0.4759271) (1.1333333,0.5002801) (1.2000000,0.5229947) (1.2666667,0.5439319) (1.3333333,0.5629611) (1.4000000,0.5799615) (1.4666667,0.5948217) (1.5333333,0.6074407) (1.6000000,0.6177285) (1.6666667,0.6256057) (1.7333333,0.6310045) (1.8000000,0.6338688) (1.8666667,0.6341539) (1.9333333,0.6318275) (2.0000000,0.6268688) (2.0666667,0.6192695) (2.1333333,0.6090330) (2.2000000,0.5961747) (2.2666667,0.5807218) (2.3333333,0.5627130) (2.4000000,0.5421987) (2.4666667,0.5192401) (2.5333333,0.4939093) (2.6000000,0.4662889) (2.6666667,0.4364715) (2.7333333,0.4045591) (2.8000000,0.3706632) (2.8666667,0.3349036) (2.9333333,0.2974083) (3.0000000,0.2583127) (3.0666667,0.2177592) (3.1333333,0.1758965) (3.2000000,0.1328788) (3.2666667,0.0888655) (3.3333333,0.0440201) (3.4000000,-0.0014901) (3.4666667,-0.0474949) (3.5333333,-0.0938219) (3.6000000,-0.1402972) (3.6666667,-0.1867457) (3.7333333,-0.2329927) (3.8000000,-0.2788636) (3.8666667,-0.3241852) (3.9333333,-0.3687861) (4.0000000,-0.4124977) (4.0666667,-0.4551543) (4.1333333,-0.4965942) (4.2000000,-0.5366602) (4.2666667,-0.5752000) (4.3333333,-0.6120670) (4.4000000,-0.6471206) (4.4666667,-0.6802270) (4.5333333,-0.7112594) (4.6000000,-0.7400988) (4.6666667,-0.7666338) (4.7333333,-0.7907618) (4.8000000,-0.8123886) (4.8666667,-0.8314293) (4.9333333,-0.8478081) (5.0000000,-0.8614590) (5.0666667,-0.8723257) (5.1333333,-0.8803617) (5.2000000,-0.8855309) (5.2666667,-0.8878071) (5.3333333,-0.8871744) (5.4000000,-0.8836272) (5.4666667,-0.8771702) (5.5333333,-0.8678183) (5.6000000,-0.8555962) (5.6666667,-0.8405390) (5.7333333,-0.8226913) (5.8000000,-0.8021075) (5.8666667,-0.7788514) (5.9333333,-0.7529956) (6.0000000,-0.7246221) (6.0666667,-0.6938211) (6.1333333,-0.6606910) (6.2000000,-0.6253381) (6.2666667,-0.5878760) (6.3333333,-0.5484256) (6.4000000,-0.5071141) (6.4666667,-0.4640749) (6.5333333,-0.4194471) (6.6000000,-0.3733747) (6.6666667,-0.3260068) (6.7333333,-0.2774961) (6.8000000,-0.2279994) (6.8666667,-0.1776761) (6.9333333,-0.1266884) (7.0000000,-0.0752003) (7.0666667,-0.0233774) (7.1333333,0.0286141) (7.2000000,0.0806074) (7.2666667,0.1324362) (7.3333333,0.1839349) (7.4000000,0.2349391) (7.4666667,0.2852866) (7.5333333,0.3348173) (7.6000000,0.3833740) (7.6666667,0.4308031) (7.7333333,0.4769546) (7.8000000,0.5216832) (7.8666667,0.5648481) (7.9333333,0.6063139) (8.0000000,0.6459508) (8.0666667,0.6836351) (8.1333333,0.7192495) (8.2000000,0.7526833) (8.2666667,0.7838331) (8.3333333,0.8126029) (8.4000000,0.8389043) (8.4666667,0.8626569) (8.5333333,0.8837882) (8.6000000,0.9022344) (8.6666667,0.9179399) (8.7333333,0.9308580) (8.8000000,0.9409507) (8.8666667,0.9481888) (8.9333333,0.9525519) (9.0000000,0.9540289) (9.0666667,0.9526173) (9.1333333,0.9483236) (9.2000000,0.9411634) (9.2666667,0.9311608) (9.3333333,0.9183487) (9.4000000,0.9027685) (9.4666667,0.8844702) (9.5333333,0.8635116) (9.6000000,0.8399589) (9.6666667,0.8138859) (9.7333333,0.7853739) (9.8000000,0.7545115) (9.8666667,0.7213942) (9.9333333,0.6861242) (10.0000000,0.6488100) (10.0666667,0.6095659) (10.1333333,0.5685119) (10.2000000,0.5257733) (10.2666667,0.4814801) (10.3333333,0.4357667) (10.4000000,0.3887716) (10.4666667,0.3406365) (10.5333333,0.2915068) (10.6000000,0.2415300) (10.6666667,0.1908562) (10.7333333,0.1396372) (10.8000000,0.0880259) (10.8666667,0.0361763) (10.9333333,-0.0157573) (11.0000000,-0.0676207) (11.0666667,-0.1192602) (11.1333333,-0.1705230) (11.2000000,-0.2212578) (11.2666667,-0.2713153) (11.3333333,-0.3205483) (11.4000000,-0.3688124) (11.4666667,-0.4159663) (11.5333333,-0.4618724) (11.6000000,-0.5063970) (11.6666667,-0.5494106) (11.7333333,-0.5907885) (11.8000000,-0.6304111) (11.8666667,-0.6681641) (11.9333333,-0.7039389) (12.0000000,-0.7376330) (12.0666667,-0.7691501) (12.1333333,-0.7984005) (12.2000000,-0.8253013) (12.2666667,-0.8497765) (12.3333333,-0.8717574) (12.4000000,-0.8911829) (12.4666667,-0.9079992) (12.5333333,-0.9221602) (12.6000000,-0.9336277) (12.6666667,-0.9423715) (12.7333333,-0.9483691) (12.8000000,-0.9516064) (12.8666667,-0.9520769) (12.9333333,-0.9497825) (13.0000000,-0.9447329) (13.0666667,-0.9369459) (13.1333333,-0.9264470) (13.2000000,-0.9132697) (13.2666667,-0.8974551) (13.3333333,-0.8790520) (13.4000000,-0.8581163) (13.4666667,-0.8347114) (13.5333333,-0.8089076) (13.6000000,-0.7807821) (13.6666667,-0.7504185) (13.7333333,-0.7179068) (13.8000000,-0.6833432) (13.8666667,-0.6468295) (13.9333333,-0.6084730) (14.0000000,-0.5683861) (14.0666667,-0.5266861) (14.1333333,-0.4834947) (14.2000000,-0.4389377) (14.2666667,-0.3931446) (14.3333333,-0.3462485) (14.4000000,-0.2983852) (14.4666667,-0.2496933) (14.5333333,-0.2003134) (14.6000000,-0.1503881) (14.6666667,-0.1000612) (14.7333333,-0.0494775) (14.8000000,0.0012176) (14.8666667,0.0518785) (14.9333333,0.1023599) (15.0000000,0.1525175) (15.0666667,0.2022076) (15.1333333,0.2512883) (15.2000000,0.2996196) (15.2666667,0.3470638) (15.3333333,0.3934856) (15.4000000,0.4387532) (15.4666667,0.4827380) (15.5333333,0.5253152) (15.6000000,0.5663642) (15.6666667,0.6057688) (15.7333333,0.6434177) (15.8000000,0.6792047) (15.8666667,0.7130290) (15.9333333,0.7447954) (16.0000000,0.7744147)};\draw[dashed]plot coordinates{(0.0000000,0.6215470) (0.0666667,0.6574781) (0.1333333,0.6915396) (0.2000000,0.7236345) (0.2666667,0.7536715) (0.3333333,0.7815653) (0.4000000,0.8072365) (0.4666667,0.8306121) (0.5333333,0.8516256) (0.6000000,0.8702173) (0.6666667,0.8863343) (0.7333333,0.8999307) (0.8000000,0.9109680) (0.8666667,0.9194146) (0.9333333,0.9252467) (1.0000000,0.9284476) (1.0666667,0.9290082) (1.1333333,0.9269269) (1.2000000,0.9222096) (1.2666667,0.9148698) (1.3333333,0.9049283) (1.4000000,0.8924134) (1.4666667,0.8773607) (1.5333333,0.8598130) (1.6000000,0.8398201) (1.6666667,0.8174390) (1.7333333,0.7927333) (1.8000000,0.7657733) (1.8666667,0.7366355) (1.9333333,0.7054030) (2.0000000,0.6721644) (2.0666667,0.6370143) (2.1333333,0.6000528) (2.2000000,0.5613848) (2.2666667,0.5211203) (2.3333333,0.4793740) (2.4000000,0.4362644) (2.4666667,0.3919141) (2.5333333,0.3464494) (2.6000000,0.2999994) (2.6666667,0.2526963) (2.7333333,0.2046745) (2.8000000,0.1560708) (2.8666667,0.1070232) (2.9333333,0.0576712) (3.0000000,0.0081553) (3.0666667,-0.0413839) (3.1333333,-0.0908053) (3.2000000,-0.1399686) (3.2666667,-0.1887337) (3.3333333,-0.2369622) (3.4000000,-0.2845168) (3.4666667,-0.3312623) (3.5333333,-0.3770658) (3.6000000,-0.4217970) (3.6666667,-0.4653287) (3.7333333,-0.5075371) (3.8000000,-0.5483022) (3.8666667,-0.5875080) (3.9333333,-0.6250431) (4.0000000,-0.6608007) (4.0666667,-0.6946792) (4.1333333,-0.7265821) (4.2000000,-0.7564188) (4.2666667,-0.7841045) (4.3333333,-0.8095603) (4.4000000,-0.8327139) (4.4666667,-0.8534995) (4.5333333,-0.8718579) (4.6000000,-0.8877369) (4.6666667,-0.9010915) (4.7333333,-0.9118835) (4.8000000,-0.9200824) (4.8666667,-0.9256647) (4.9333333,-0.9286147) (5.0000000,-0.9289239) (5.0666667,-0.9265914) (5.1333333,-0.9216240) (5.2000000,-0.9140356) (5.2666667,-0.9038480) (5.3333333,-0.8910900) (5.4000000,-0.8757979) (5.4666667,-0.8580153) (5.5333333,-0.8377927) (5.6000000,-0.8151876) (5.6666667,-0.7902643) (5.7333333,-0.7630936) (5.8000000,-0.7337529) (5.8666667,-0.7023256) (5.9333333,-0.6689010) (6.0000000,-0.6335743) (6.0666667,-0.5964457) (6.1333333,-0.5576211) (6.2000000,-0.5172106) (6.2666667,-0.4753294) (6.3333333,-0.4320964) (6.4000000,-0.3876347) (6.4666667,-0.3420706) (6.5333333,-0.2955337) (6.6000000,-0.2481564) (6.6666667,-0.2000734) (6.7333333,-0.1514214) (6.8000000,-0.1023388) (6.8666667,-0.0529652) (6.9333333,-0.0034410) (7.0000000,0.0460930) (7.0666667,0.0954959) (7.1333333,0.1446273) (7.2000000,0.1933474) (7.2666667,0.2415176) (7.3333333,0.2890010) (7.4000000,0.3356626) (7.4666667,0.3813696) (7.5333333,0.4259921) (7.6000000,0.4694032) (7.6666667,0.5114794) (7.7333333,0.5521010) (7.8000000,0.5911526) (7.8666667,0.6285231) (7.9333333,0.6641063) (8.0000000,0.6978009) (8.0666667,0.7295110) (8.1333333,0.7591467) (8.2000000,0.7866234) (8.2666667,0.8118633) (8.3333333,0.8347943) (8.4000000,0.8553514) (8.4666667,0.8734761) (8.5333333,0.8891168) (8.6000000,0.9022290) (8.6666667,0.9127756) (8.7333333,0.9207264) (8.8000000,0.9260588) (8.8666667,0.9287578) (8.9333333,0.9288156) (9.0000000,0.9262321) (9.0666667,0.9210145) (9.1333333,0.9131779) (9.2000000,0.9027443) (9.2666667,0.8897436) (9.3333333,0.8742126) (9.4000000,0.8561955) (9.4666667,0.8357437) (9.5333333,0.8129152) (9.6000000,0.7877749) (9.6666667,0.7603944) (9.7333333,0.7308515) (9.8000000,0.6992302) (9.8666667,0.6656205) (9.9333333,0.6301179) (10.0000000,0.5928234) (10.0666667,0.5538430) (10.1333333,0.5132876) (10.2000000,0.4712726) (10.2666667,0.4279174) (10.3333333,0.3833452) (10.4000000,0.3376830) (10.4666667,0.2910604) (10.5333333,0.2436101) (10.6000000,0.1954671) (10.6666667,0.1467682) (10.7333333,0.0976519) (10.8000000,0.0482579) (10.8666667,-0.0012733) (10.9333333,-0.0508009) (11.0000000,-0.1001841) (11.0666667,-0.1492823) (11.1333333,-0.1979560) (11.2000000,-0.2460668) (11.2666667,-0.2934778) (11.3333333,-0.3400542) (11.4000000,-0.3856636) (11.4666667,-0.4301762) (11.5333333,-0.4734656) (11.6000000,-0.5154085) (11.6666667,-0.5558857) (11.7333333,-0.5947820) (11.8000000,-0.6319870) (11.8666667,-0.6673947) (11.9333333,-0.7009046) (12.0000000,-0.7324212) (12.0666667,-0.7618549) (12.1333333,-0.7891222) (12.2000000,-0.8141453) (12.2666667,-0.8368532) (12.3333333,-0.8571813) (12.4000000,-0.8750718) (12.4666667,-0.8904737) (12.5333333,-0.9033434) (12.6000000,-0.9136441) (12.6666667,-0.9213467) (12.7333333,-0.9264292) (12.8000000,-0.9288771) (12.8666667,-0.9286835) (12.9333333,-0.9258489) (13.0000000,-0.9203814) (13.0666667,-0.9122966) (13.1333333,-0.9016174) (13.2000000,-0.8883742) (13.2666667,-0.8726047) (13.3333333,-0.8543537) (13.4000000,-0.8336732) (13.4666667,-0.8106218) (13.5333333,-0.7852652) (13.6000000,-0.7576755) (13.6666667,-0.7279312) (13.7333333,-0.6961168) (13.8000000,-0.6623228) (13.8666667,-0.6266452) (13.9333333,-0.5891857) (14.0000000,-0.5500507) (14.0666667,-0.5093514) (14.1333333,-0.4672036) (14.2000000,-0.4237273) (14.2666667,-0.3790459) (14.3333333,-0.3332867) (14.4000000,-0.2865796) (14.4666667,-0.2390576) (14.5333333,-0.1908557) (14.6000000,-0.1421111) (14.6666667,-0.0929624) (14.7333333,-0.0435493) (14.8000000,0.0059876) (14.8666667,0.0555076) (14.9333333,0.1048696) (15.0000000,0.1539335) (15.0666667,0.2025595) (15.1333333,0.2506096) (15.2000000,0.2979470) (15.2666667,0.3444371) (15.3333333,0.3899476) (15.4000000,0.4343493) (15.4666667,0.4775158) (15.5333333,0.5193243) (15.6000000,0.5596560) (15.6666667,0.5983961) (15.7333333,0.6354346) (15.8000000,0.6706660) (15.8666667,0.7039902) (15.9333333,0.7353124) (16.0000000,0.7645436)};\foreach\t in{4,8,12,16}{\draw(\t,.04)--(\t,-.04)node[below]{$\t$};}\node at(8,1.55){含初值解与稳态};\end{scope}\begin{scope}[shift={(8.2,0)},x=.3cm,y=1.2cm]\draw[axis](0,0)--(12.5,0)node[right]{$\tau$};\draw[axis](0,-1.2)--(0,1.5)node[above]{$x$};\draw[curve]plot coordinates{(0.0000000,0.0000000) (0.0600000,0.0002878) (0.1200000,0.0011492) (0.1800000,0.0025780) (0.2400000,0.0045639) (0.3000000,0.0070925) (0.3600000,0.0101455) (0.4200000,0.0137008) (0.4800000,0.0177323) (0.5400000,0.0222107) (0.6000000,0.0271028) (0.6600000,0.0323726) (0.7200000,0.0379806) (0.7800000,0.0438846) (0.8400000,0.0500400) (0.9000000,0.0563995) (0.9600000,0.0629139) (1.0200000,0.0695320) (1.0800000,0.0762012) (1.1400000,0.0828674) (1.2000000,0.0894758) (1.2600000,0.0959707) (1.3200000,0.1022963) (1.3800000,0.1083966) (1.4400000,0.1142160) (1.5000000,0.1196994) (1.5600000,0.1247927) (1.6200000,0.1294432) (1.6800000,0.1335994) (1.7400000,0.1372121) (1.8000000,0.1402341) (1.8600000,0.1426205) (1.9200000,0.1443295) (1.9800000,0.1453222) (2.0400000,0.1455628) (2.1000000,0.1450192) (2.1600000,0.1436631) (2.2200000,0.1414700) (2.2800000,0.1384198) (2.3400000,0.1344966) (2.4000000,0.1296889) (2.4600000,0.1239900) (2.5200000,0.1173979) (2.5800000,0.1099152) (2.6400000,0.1015497) (2.7000000,0.0923141) (2.7600000,0.0822257) (2.8200000,0.0713072) (2.8800000,0.0595859) (2.9400000,0.0470941) (3.0000000,0.0338688) (3.0600000,0.0199517) (3.1200000,0.0053891) (3.1800000,-0.0097684) (3.2400000,-0.0254660) (3.3000000,-0.0416449) (3.3600000,-0.0582423) (3.4200000,-0.0751920) (3.4800000,-0.0924247) (3.5400000,-0.1098677) (3.6000000,-0.1274459) (3.6600000,-0.1450816) (3.7200000,-0.1626953) (3.7800000,-0.1802054) (3.8400000,-0.1975292) (3.9000000,-0.2145830) (3.9600000,-0.2312825) (4.0200000,-0.2475430) (4.0800000,-0.2632800) (4.1400000,-0.2784098) (4.2000000,-0.2928495) (4.2600000,-0.3065174) (4.3200000,-0.3193338) (4.3800000,-0.3312210) (4.4400000,-0.3421041) (4.5000000,-0.3519108) (4.5600000,-0.3605724) (4.6200000,-0.3680237) (4.6800000,-0.3742036) (4.7400000,-0.3790555) (4.8000000,-0.3825272) (4.8600000,-0.3845719) (4.9200000,-0.3851479) (4.9800000,-0.3842191) (5.0400000,-0.3817553) (5.1000000,-0.3777324) (5.1600000,-0.3721325) (5.2200000,-0.3649441) (5.2800000,-0.3561625) (5.3400000,-0.3457896) (5.4000000,-0.3338343) (5.4600000,-0.3203121) (5.5200000,-0.3052457) (5.5800000,-0.2886649) (5.6400000,-0.2706060) (5.7000000,-0.2511126) (5.7600000,-0.2302350) (5.8200000,-0.2080301) (5.8800000,-0.1845615) (5.9400000,-0.1598993) (6.0000000,-0.1341194) (6.0600000,-0.1073042) (6.1200000,-0.0795414) (6.1800000,-0.0509243) (6.2400000,-0.0215514) (6.3000000,0.0084742) (6.3600000,0.0390449) (6.4200000,0.0700490) (6.4800000,0.1013713) (6.5400000,0.1328933) (6.6000000,0.1644938) (6.6600000,0.1960494) (6.7200000,0.2274346) (6.7800000,0.2585229) (6.8400000,0.2891868) (6.9000000,0.3192987) (6.9600000,0.3487311) (7.0200000,0.3773571) (7.0800000,0.4050515) (7.1400000,0.4316904) (7.2000000,0.4571527) (7.2600000,0.4813199) (7.3200000,0.5040771) (7.3800000,0.5253131) (7.4400000,0.5449213) (7.5000000,0.5628000) (7.5600000,0.5788527) (7.6200000,0.5929890) (7.6800000,0.6051246) (7.7400000,0.6151821) (7.8000000,0.6230910) (7.8600000,0.6287886) (7.9200000,0.6322198) (7.9800000,0.6333376) (8.0400000,0.6321038) (8.1000000,0.6284886) (8.1600000,0.6224713) (8.2200000,0.6140404) (8.2800000,0.6031936) (8.3400000,0.5899381) (8.4000000,0.5742905) (8.4600000,0.5562771) (8.5200000,0.5359339) (8.5800000,0.5133062) (8.6400000,0.4884490) (8.7000000,0.4614266) (8.7600000,0.4323126) (8.8200000,0.4011896) (8.8800000,0.3681493) (8.9400000,0.3332917) (9.0000000,0.2967253) (9.0600000,0.2585665) (9.1200000,0.2189394) (9.1800000,0.1779752) (9.2400000,0.1358119) (9.3000000,0.0925941) (9.3600000,0.0484718) (9.4200000,0.0036007) (9.4800000,-0.0418591) (9.5400000,-0.0877430) (9.6000000,-0.1338830) (9.6600000,-0.1801079) (9.7200000,-0.2262445) (9.7800000,-0.2721176) (9.8400000,-0.3175511) (9.9000000,-0.3623684) (9.9600000,-0.4063935) (10.0200000,-0.4494512) (10.0800000,-0.4913680) (10.1400000,-0.5319728) (10.2000000,-0.5710977) (10.2600000,-0.6085785) (10.3200000,-0.6442553) (10.3800000,-0.6779734) (10.4400000,-0.7095840) (10.5000000,-0.7389444) (10.5600000,-0.7659192) (10.6200000,-0.7903804) (10.6800000,-0.8122083) (10.7400000,-0.8312917) (10.8000000,-0.8475289) (10.8600000,-0.8608278) (10.9200000,-0.8711064) (10.9800000,-0.8782935) (11.0400000,-0.8823286) (11.1000000,-0.8831627) (11.1600000,-0.8807584) (11.2200000,-0.8750901) (11.2800000,-0.8661442) (11.3400000,-0.8539197) (11.4000000,-0.8384276) (11.4600000,-0.8196916) (11.5200000,-0.7977477) (11.5800000,-0.7726445) (11.6400000,-0.7444431) (11.7000000,-0.7132166) (11.7600000,-0.6790506) (11.8200000,-0.6420423) (11.8800000,-0.6023006) (11.9400000,-0.5599461) (12.0000000,-0.5151100)};\draw[guide](0,0)--(12,.96) (0,0)--(12,-.96);\node at(6,1.55){无阻尼共振};\end{scope}
\end{tikzpicture}
\caption{左图取$\omega_0=1,\gamma=.25,\Omega=.8,F_0/m=.5,v_0=.4$，实线含初值，虚线为稳态；右图为$.08\tau\sin\tau$，$\tau=\omega_0t$。}\label{fig:c2-e7-s}
\end{figure}
记齐次部分的初值为
$A=-x_p(0)=-C$、$w=v_0-\dot x_p(0)=v_0-\Omega D$。
在欠阻尼情形，取$\omega_d=\sqrt{\omega_0^2-\gamma^2}$，有
\[
 x=x_p+\ee^{-\gamma t}(A\cos\omega_dt+B\sin\omega_dt).
\]
零时刻齐次部分的导数为$-\gamma A+\omega_dB=w$，因此
\[
 \boxed{B=\frac{w+\gamma A}{\omega_d}}.
\]
图~\ref{fig:c2-e7-s}中的完整解由给定初态出发，衰减的齐次部分消失后才接近周期曲线。
另两种阻尼使用相同的$A,w$，分别为
\[
 \begin{aligned}
 \gamma=\omega_0:\quad
 x&=x_p+\ee^{-\omega_0t}[A+(w+\omega_0A)t],\\
 \gamma>\omega_0:\quad
 x&=x_p+\ee^{-\gamma t}
 \left[A\cosh(dt)+\frac{w+\gamma A}{d}\sinh(dt)\right],
 \quad d=\sqrt{\gamma^2-\omega_0^2}.
 \end{aligned}
\]
这两式的初值系数可分别从齐次通解求出：临界解 $\ee^{-\omega_0t}(A+Bt)$ 的零时刻导数为 $B-\omega_0A=w$，故 $B=w+\omega_0A$；过阻尼解 $\ee^{-\gamma t}(A\cosh dt+B\sinh dt)$ 的零时刻导数为 $dB-\gamma A=w$，故 $B=(w+\gamma A)/d$。

无阻尼时，若$\Omega\ne\omega_0$，由同样的初值修正得到
\[
 x(t)=\frac{v_0}{\omega_0}\sin\omega_0t+
 \frac{F_0}{m(\omega_0^2-\Omega^2)}
 (\cos\Omega t-\cos\omega_0t).
\]
共振处须将余弦差与分母一起取极限。对$\Omega$用一次洛必达法则，
\[
 \lim_{\Omega\to\omega_0}
 \frac{\cos\Omega t-\cos\omega_0t}{\omega_0^2-\Omega^2}
 =\frac{-t\sin\omega_0t}{-2\omega_0}.
\]
所以
\[
 \boxed{x(t)=\frac{v_0}{\omega_0}\sin\omega_0t+
 \frac{F_0}{2m\omega_0}t\sin\omega_0t}.
\]
此时$\mathcal D=0$，有界周期试探式失效，响应具有随时间增长的包络；有限阻尼下的周期稳态则有有限振幅。
\end{solution}


\begin{example}[缓慢增大后突然撤去的外力]\label{ex:c2-v1}

质量$m$的无阻尼弹簧振子，劲度系数$k$，$\omega=\sqrt{k/m}$，初始静止于平衡点。外力在$0<t<\tau$为$F(t)=F_0t/\tau$，随后为零，$F_0>0$。求撤力前后的运动，并说明$t=\tau$处各阶运动量的连接关系。 见图~\ref{fig:c2-v1-q}。
\end{example}

\begin{figure}[H]\centering
\begin{tikzpicture}[>=Stealth,line cap=round,line join=round]

\draw[axis](-.3,0)--(4.3,0)node[right]{$t$};\draw[axis](0,-.2)--(0,2)node[above]{$F$};
\draw[curve](0,0)--(2.4,1.5)--(2.4,0)--(4,0);\node[below]at(2.4,0){$\tau$};\node[left]at(0,1.5){$F_0$};\draw[guide](0,1.5)--(2.4,1.5);

\end{tikzpicture}
\caption{外力在开启时连续，从零线性增大，在$\tau$时由$F_0$降至零。}\label{fig:c2-v1-q}
\end{figure}
\begin{solution}
图~\ref{fig:c2-v1-q}的外力有限，故撤力不会产生速度冲量。先在$0<t<\tau$解完整初值问题：
\[\ddot x+\omega^2x=\frac{F_0}{m\tau}t,
\quad x(0)=\dot x(0)=0.\]
取一次函数特解 $x_p=ct$，$\ddot x_p=0$，比较 $\omega^2ct=F_0t/(m\tau)$ 得 $c=F_0/(k\tau)$。通解 $x=ct+C\cos\omega t+D\sin\omega t$ 求导为 $\dot x=c-C\omega\sin\omega t+D\omega\cos\omega t$；两个初值分别给出 $C=0$、$c+D\omega=0$。因此
\[x(t)=\frac{F_0}{k}\left[\frac t\tau-\frac{\sin\omega t}{\omega\tau}\right],
\quad\dot x(t)=\frac{F_0}{k\tau}(1-\cos\omega t).\]
记$b=\omega\tau$，撤力瞬间
\[x_\tau=\frac{F_0}{k}\left(1-\frac{\sin b}{b}\right),\qquad
v_\tau=\frac{F_0}{k\tau}(1-\cos b).\]
撤力后只剩齐次方程，以这两个量为新初值，故
\[\boxed{x(t)=x_\tau\cos[\omega(t-\tau)]
+\frac{v_\tau}{\omega}\sin[\omega(t-\tau)]},\quad t\ge\tau.\]
图~\ref{fig:c2-v1-s}在同一坐标中连接两段，接点的值和切线相同。由原方程直接求跳变，
\[x(\tau^+)=x(\tau^-),\quad
\dot x(\tau^+)=\dot x(\tau^-),\quad
\ddot x(\tau^+)-\ddot x(\tau^-)=-\frac{F_0}{m}.\]
有限力跃变导致加速度跃变；如果输入改成冲量，才需允许速度跳变。

也可直接按响应核叠加。对单位加速度输入，$G(u)=\sin(\omega u)/\omega$（$u\ge0$，负时间为零），因果积分为
\[x(t)=\frac{F_0}{m\tau\omega}
\int_0^{\min(t,\tau)}s\sin[\omega(t-s)]\dd s.\]
积分上限选择尚未撤力或已经撤力的时间区间。分部积分时，对 $s$ 求微分，对正弦积分，得到
\[\int s\sin[\omega(t-s)]\dd s
=\frac{s\cos[\omega(t-s)]}{\omega}
+\frac{\sin[\omega(t-s)]}{\omega^2},\]
其中 $\int\sin[\omega(t-s)]\dd s=\cos[\omega(t-s)]/\omega$，相位的负导数确定了此处正号。令 $a=\min(t,\tau)$，上下限代入为
\[
\int_0^as\sin[\omega(t-s)]\dd s
=\frac{a\cos[\omega(t-a)]}{\omega}
+\frac{\sin[\omega(t-a)]-\sin\omega t}{\omega^2}.
\]
当 $a=t$ 时，第一项为 $t/\omega$；当 $a=\tau$ 时，用正弦加角公式把 $\sin\omega t$ 展开，可读出余弦、正弦系数分别为上面的 $x_\tau$、$v_\tau/\omega$。第四讲继续计算这一输入的 Laplace 形式，其中的延迟因子对应同一撤力边界。
\begin{figure}[H]\centering
\begin{tikzpicture}[>=Stealth,line cap=round,line join=round]

\begin{scope}[x=.65cm,y=1.25cm]
\draw[axis](0,0)--(10,0)node[right]{$\omega t$};\draw[axis](0,-1.4)--(0,1.5)node[above]{$kx/F_0$};
\draw[curve,domain=0:3.14159265,samples=81]plot(\x,{(\x-sin(\x r))/pi});
\draw[curve,domain=3.14159265:9.5,samples=121]plot(\x,{cos((\x-pi) r)+2/pi*sin((\x-pi) r)});
\draw[guide](3.14159265,-1.3)--(3.14159265,1.35);\fill(3.14159265,1)circle(1.5pt);\node[below]at(3.14159265,0){$b$};
\end{scope}

\end{tikzpicture}
\caption{取$b=\omega\tau=\pi$；撤力时位移和速度连续，曲线的曲率发生跳变。}\label{fig:c2-v1-s}
\end{figure}
\end{solution}



\section{约束运动中的降阶与终止条件}
运动方程不显含某个变量时，可以利用第一积分降阶。降阶得到的分支仍受接触、可达区域和过程终止条件限制，转折点附近尤其需要保留原方程的含义。

绕筒运动的初值解还受到接触几何和可达边界的限制。不可伸长条件联系自由绳长与绕入角，能量、角动量和张力方程则用于消去其余变量。固定筒与自由筒的守恒量不同，不能只在同一公式中替换一个参数。以下两题完整讨论首次接触至完全绕入或相对转向这一段过程。
\begin{example}[质点绕固定圆筒]\label{ex:c2-e2}
光滑水平面上固定一半径$r$的圆筒。轻而不可伸长的绳一端固定在筒上，初始自由直段长度$\ell_0$，另一端系质量$m$的质点。绳始终绷紧并逐渐绕上筒，质点初速$v_0$垂直绳，方向使自由段缩短。求完全绕上的理想极限时间。 见图~\ref{fig:c2-e2-q}。
\end{example}

\begin{figure}[H]\centering
\begin{tikzpicture}[>=Stealth,line cap=round,line join=round]
\draw[curve](0,0)circle(1.1);\draw[guide](0,0)--(0,1.55);\fill(0,0)circle(1.4pt)node[below left]{$O$};\draw[guide](0,0)--(0.749803,0.804858)node[midway,below]{$r$};\draw[thick](0.749803,0.804858)--(2.140011,-0.490256)node[midway,right]{$\ell$};\fill(0.749803,0.804858)circle(1.4pt)node[above right]{$P$};\fill(2.140011,-0.490256)circle(2pt)node[right]{$m$};\draw(0,.55)arc(90:47.028165:.55);\node at(.22,.75){$q$};
\end{tikzpicture}
\caption{切点$P$随绳绕入而移动，质点沿自由绳段的垂直方向运动。}\label{fig:c2-e2-q}
\end{figure}
\begin{solution}
固定筒不向系统输入能量，绳张力与质点速度正交，所以质点速率恒定。变化的是接触几何：绳自由段越短，切点转角变化越快。图~\ref{fig:c2-e2-q}中$P$是几何切点，质点沿圆筒渐开线运动，并不是以筒心为圆心作匀速圆周运动。


令$q$为切点半径从初始半径起转过的角，$\ell=\ell_0-rq$。质点的位置可写成
\[\vct R=r\vct e_r(q)+\ell\vct e_t(q),\qquad
\frac{\dd\vct e_r}{\dd q}=\vct e_t,\quad
\frac{\dd\vct e_t}{\dd q}=-\vct e_r.
\]
对位置按乘积法则求导，
\[
\dot{\vct R}=r\dot q\vct e_t+\dot\ell\vct e_t-\ell\dot q\vct e_r.
\]
图~\ref{fig:c2-e2-s}中，前两项沿 $\vct e_t$ 的速度由 $\dot\ell=-r\dot q$ 相消，留下 $\vct v=-\ell\dot q\vct e_r$。因此 $\dot q=v_0/\ell$，
\[T_1=\int_0^{\ell_0/r}\frac{\ell_0-rq}{v_0}\dd q
=\frac1{v_0}\left[\ell_0q-\frac{rq^2}{2}\right]_0^{\ell_0/r}
=\boxed{\frac{\ell_0^2}{2rv_0}}.
\]
也可直接求绳长：$\ell\dot\ell=-rv_0$，分离变量积分得 $[\ell^2/2]_{\ell_0}^{\ell(t)}=-rv_0[t]_0^t$，所以 $\ell(t)=\sqrt{\ell_0^2-2rv_0t}$，取正根直到 $T_1$。自由段趋零时角速度发散，但质点速率仍为 $v_0$，到达时间有限。

由于 $\ell\dot q=v_0$ 恒定，对 $\vct v=-v_0\vct e_r$ 再求导，$\vct a=-v_0\dot q\vct e_t=-(v_0^2/\ell)\vct e_t$，故张力为 $mv_0^2/\ell$。完全绕上只作为轻绳与质点模型的极限；有限强度、绳厚或有限质点尺寸会在此前改变运动。若初速度方向相反，自由段先增大，此时不能仍把上述积分上限作为到达事件。


剩余速度垂直自由绳，张力沿 $-\vct e_t$，两者点积为零，与开头使用的能量守恒一致。图~\ref{fig:c2-e2-s}显示，自由段越短，速度方向转动越快，因而角速度与张力虽在终点发散，质点速率仍不发散。

\begin{figure}[H]\centering
\begin{tikzpicture}[>=Stealth,line cap=round,line join=round]
\draw[curve](0,0)circle(1.1);\draw[guide](0,0)--(0,1.55);\fill(0,0)circle(1.4pt)node[below left]{$O$};\draw[guide](0,0)--(0.749803,0.804858)node[midway,below]{$r$};\draw[thick](0.749803,0.804858)--(2.140011,-0.490256)node[midway,right]{$\ell$};\fill(0.749803,0.804858)circle(1.4pt)node[above right]{$P$};\fill(2.140011,-0.490256)circle(2pt)node[right]{$m$};\draw(0,.55)arc(90:47.028165:.55);\node at(.22,.75){$q$};\node[right,align=left]at(3.6,.1){$r\dot q+\dot\ell=0$\\[5pt]$v=\ell\dot q$\\[5pt]$\mathcal T=mv_0^2/\ell$};
\end{tikzpicture}
\caption{切点移动与自由段缩短共同决定质点速度，不能将$r\dot q$直接当作质点速率。}\label{fig:c2-e2-s}
\end{figure}
\end{solution}

\begin{example}[可自由转动的圆筒与绕入条件]\label{ex:c2-e3}
将上题的圆筒改为绕固定竖直轴自由转动的薄壁圆筒，质量也为$m$，惯量$I=mr^2$，初始静止。其余初态相同。绳在筒上不滑动。求完全绕上的时间；若并非任意$\ell_0$都能绕完，给出可达条件。 见图~\ref{fig:c2-e3-q}。
\end{example}

\begin{figure}[H]\centering
\begin{tikzpicture}[>=Stealth,line cap=round,line join=round]
\draw[curve](0,0)circle(1.1);\draw[guide](0,0)--(0,1.55);\fill(0,0)circle(1.4pt)node[below left]{$O$};\draw[guide](0,0)--(0.925618,0.594333)node[midway,below]{$r$};\draw[thick](0.925618,0.594333)--(1.898162,-0.920315)node[midway,right]{$\ell$};\fill(0.925618,0.594333)circle(1.4pt)node[above right]{$P$};\fill(1.898162,-0.920315)circle(2pt)node[right]{$m$};\draw(0,.55)arc(90:32.704220:.55);\node at(.22,.75){$q$};\fill(0.325072,1.050870)circle(1.5pt)node[above]{$A$};\node[right,align=left]at(3.5,1){$\theta=q-\alpha$\\[5pt]$\ell=\ell_0-r\theta$};
\end{tikzpicture}
\caption{$A$是随筒转动的材料标记；$q$为切点的绝对角，$\alpha$为圆筒转角。}\label{fig:c2-e3-q}
\end{figure}
\begin{solution}
图~\ref{fig:c2-e3-q}中，材料标记$A$随筒转过角$\alpha$，绝对切点角为$q$，相对绕入角为$\theta=q-\alpha$，均取顺时针为正。自由绳长及其导数为
\[
 \ell=\ell_0-r\theta,\qquad
 \dot\ell=-r(\dot q-\dot\alpha).
\]
仍以 $\vct R=r\vct e_r+\ell\vct e_t$ 描述质点位置，求导得 $\dot{\vct R}=(r\dot q+\dot\ell)\vct e_t-\ell\dot q\vct e_r$。将绳长导数代入，第一项系数为 $r\dot q-r(\dot q-\dot\alpha)=r\dot\alpha$，得到图~\ref{fig:c2-e3-s}中的两个正交速度分量：
\[
 \vct v=r\dot\alpha\vct e_t-\ell\dot q\vct e_r,\qquad
 v^2=r^2\dot\alpha^2+\ell^2\dot q^2.
\]
圆筒转动动能与质点切向动能各贡献$mr^2\dot\alpha^2/2$；
质点对轴角动量为$m(r^2\dot\alpha+\ell^2\dot q)$，
圆筒自身再贡献$mr^2\dot\alpha$。初始$\dot\alpha=0$、
$\dot q=v_0/\ell_0$，能量和轴角动量守恒分别写成
\begin{align}
 2r^2\dot\alpha^2+\ell^2\dot q^2&=v_0^2,
 \label{eq:c2-windE}\\
 2r^2\dot\alpha+\ell^2\dot q&=\ell_0v_0.
 \label{eq:c2-windL}
\end{align}
\begin{figure}[H]\centering
\begin{tikzpicture}[>=Stealth,line cap=round,line join=round]
\begin{scope}[shift={(-2.6,0)}]
\draw[axis](0,0)--(0,1.4)node[above]{$r\dot\alpha\vct e_t$};
\draw[axis](0,0)--(-1.8,0)node[above]{$-\ell\dot q\vct e_r$};
\draw[axis](0,0)--(-1.8,1.4)node[above left]{$\vct v$};
\draw[guide](-1.8,0)--(-1.8,1.4)--(0,1.4);\node[below]at(0,0){$m$};
\node at(-.7,-.8){局部正交分解};\end{scope}
\begin{scope}[shift={(1.9,0)},x=1.1cm,y=.85cm]
\draw[axis](0,0)--(3.15,0)node[right]{$\ell$};\draw[axis](0,-.45)--(0,2.5)node[above]{$A+\ell^2-\ell_0^2$};
\draw[curve,domain=.3:2.5,samples=101]plot(\x,{.38*((\x)^2-.64)});
\draw[guide](.8,-.4)--(.8,.3);\fill(.8,0)circle(1.5pt)node[below right]{$\ell_{\min}$};
\node at(1.6,2.1){$\ell_0>\sqrt2r$};\end{scope}
\end{tikzpicture}
\caption{左图将局部基矢旋转到便于分解的位置；右图的负值段不可达，零点是第一次相对转向边界。}\label{fig:c2-e3-s}
\end{figure}
记$A=2r^2$、$p=\ell_0v_0$，用$\dot q=\dot\alpha+\dot\theta$消元：
\[
 \dot\alpha=\frac{p-\ell^2\dot\theta}{A+\ell^2},\qquad
 \dot q=\frac{p+A\dot\theta}{A+\ell^2}.
\]
代入能量式，分子先展开为
\[
\begin{aligned}
&A(p-\ell^2\dot\theta)^2+\ell^2(p+A\dot\theta)^2\\
&\quad=Ap^2-2Ap\ell^2\dot\theta+A\ell^4\dot\theta^2
+\ell^2p^2+2Ap\ell^2\dot\theta+A^2\ell^2\dot\theta^2\\
&\quad=(A+\ell^2)p^2+A\ell^2(A+\ell^2)\dot\theta^2.
\end{aligned}
\]
再除以共同分母 $(A+\ell^2)^2$，得到
\[
 v_0^2=\frac{p^2}{A+\ell^2}
 +\frac{A\ell^2}{A+\ell^2}\dot\theta^2,\qquad
 \boxed{\dot\theta^2=
 \frac{v_0^2(A+\ell^2-\ell_0^2)}{A\ell^2}}.
\]
右端非负决定相对绕入的可达范围。取绕入支$\dot\theta>0$，
结合 $\dot\ell=-r\dot\theta$，根式按链式法则求导，
\[
 \frac{\dd}{\dd t}\sqrt{A+\ell^2-\ell_0^2}
 =\frac{2\ell\dot\ell}{2\sqrt{A+\ell^2-\ell_0^2}}
 =-\frac{rv_0}{\sqrt A}.
\]
初始根式等于 $\sqrt A$，所以两边从 $0$ 积到 $t$，右侧为 $-(rv_0/\sqrt A)[s]_0^t$，得到
\[
 \sqrt{A+\ell^2-\ell_0^2}
 =\sqrt A-\frac{rv_0}{\sqrt A}t,\qquad
 \boxed{\ell^2(t)=\ell_0^2-2rv_0t+\frac{v_0^2t^2}{2}}.
\]
第一式要求右端非负，故本绕入支持续到完全绕上，或$t=2r/v_0$处首次相对转向。

若$\ell_0\le\sqrt2r$，长度平方在转向前可达到零。取其第一个零点，
\[
 \boxed{T_2=\frac{2r}{v_0}
 \left(1-\sqrt{1-\frac{\ell_0^2}{2r^2}}\right)},
 \qquad \ell_0\le\sqrt2r.
\]
临界$\ell_0=\sqrt2r$时，$\dot\theta=v_0/(\sqrt2r)$恒定，
$T_2=2r/v_0$，没有时间发散。
若$\ell_0>\sqrt2r$，在$t=2r/v_0$先达到
\[
 \ell_{\min}=\sqrt{\ell_0^2-2r^2}>0,\qquad
 \theta_{\max}=\frac{\ell_0-\ell_{\min}}r.
\]
此时$\dot\theta=0$，但$\dot q=\dot\alpha=p/A\ne0$，筒与切点仍一起转动；之后相对绕入角反向，不能由本初态完全绕上。

还须确认轻绳没有提前松弛。圆筒转矩及质点沿绳方向的方程为
\[
 I\ddot\alpha=r\mathcal T,\qquad
 mr\ddot\alpha-m\ell\dot q^2=-\mathcal T.
\]
代入$I=mr^2$得$\mathcal T=m\ell\dot q^2/2\ge0$，轻绳在上述阶段保持绷紧。
最后令$z=\ell_0^2/(2r^2)\ll1$。由$\sqrt{1-z}=1-z/2+O(z^2)$，时间式成为
\[
 T_2=\frac{2r}{v_0}\left[\frac z2+O(z^2)\right]
 =\frac{\ell_0^2}{2rv_0}
 +O\!\left(\frac{\ell_0^4}{r^3v_0}\right).
\]
首项等于固定筒的$T_1$，两种模型在短绳极限相合。
\end{solution}


\section{解的结构与非线性方程的局部行为}
局部存在性和唯一性决定同一初态能否对应多条解。在线性化之前，先辨认平衡点、奇异点及解的延拓范围，再讨论近似解能维持多长时间。

图~\ref{fig:c2-causal-superposition}中驱动发生在$\tau$，响应传播的时间是$t-\tau$。积分变量只遍历已发生的驱动；每一窄条先乘驱动幅值与宽度，再乘相应传播时间的核，最后相加。
\begin{figure}[H]\centering
\begin{tikzpicture}[>=Stealth,line cap=round,line join=round]
\draw[axis](-.2,0)--(6.2,0)node[right]{$\tau$};\draw[axis](0,-.1)--(0,2.3)node[above]{$f(\tau)$};
\fill[gray!16,domain=0:4.8,samples=50](0,0)--plot(\x,{.85+.5*sin(38*\x)})--(4.8,0)--cycle;
\draw[curve,domain=0:5.7,samples=60]plot(\x,{.85+.5*sin(38*\x)});
\fill[gray!50](1.75,0)rectangle(1.92,1.308);\draw[guide](4.8,0)--(4.8,1.8);\node[below]at(0,0){$0$};\node[below]at(1.83,0){$\tau$};\node[below]at(4.8,0){$t$};
\draw[->](1.85,1.53)to[bend left=22](4.8,1.53);\node[above]at(3.3,1.85){传播时间$t-\tau$};
\node[below]at(2.7,-.65){$\dd x(t)=G(t-\tau)f(\tau)\dd\tau$};
\end{tikzpicture}
\caption{因果卷积把每一次驱动的响应传播到同一个观察时刻$t$。}\label{fig:c2-causal-superposition}
\end{figure}

\subsection{一阶自治方程的相线}
对$\dot y=f(y)$，先找$f(y_*)=0$，再在各平衡点之间标出$f$的符号。$f>0$时状态向较大$y$移动，$f<0$时向较小$y$移动。导数为零的平衡不能只靠线性项判断，例如$\dot y=-y^3$仍稳定，但衰减为幂律；$\dot y=y^2$在零点两侧的行为又不相同。

局部线性化不能代替全局可达性。解接近另一个平衡、触及边界或状态变得很大时，需要回到原方程。分离变量积分$\int\dd y/f(y)$既能给出轨迹，也能判断到达特定边界是有限还是无限时间。

相线还可与扰动大小定量联系。若$f$在平衡附近二阶光滑，写$y=y_*+\eta$、$c=f'(y_*)$，则$\dot\eta=c\eta+O(\eta^2)$。线性近似下，
\[\frac{\dd}{\dd t}(\ee^{-ct}\eta)=\ee^{-ct}(\dot\eta-c\eta)=0,\]
积分并用初值$\eta(t_0)=\eta_0$，得$\eta(t)=\eta_0\ee^{c(t-t_0)}$。因此$c<0$时线性弛豫时间为$-1/c$，$c>0$时扰动增长；线性判据依赖$c\ne0$，导数为零时应直接考察$f$在两侧的符号。

稳定与渐近稳定也有区别。稳定要求初态足够接近时，解在以后仍留在平衡附近；渐近稳定还要求解趋近平衡。对一阶自治方程，相线两侧箭头都指向平衡时，通常可结合唯一性与单调有界性证明渐近趋近；一侧指向、一侧离开时，不能称为两侧稳定。

解的显式表达常能给出线性化看不到的时间尺度。习题~\ref{ex:c2-h15}的三次衰减在原点没有非零线性项，却仍从两侧向零移动，长时为幂律。在饱和模型中，初态位于容量值上方、下方或恰在容量值处，对应不同解支。分离变量时保留对数绝对值，可同时处理容量值两侧。

\subsection{Wronskian与两个独立解}
对$y''+p(t)y'+q(t)y=0$，设$y_1,y_2$为两解，定义$W=y_1y_2'-y_1'y_2$。乘积求导中两个$y_1'y_2'$抵消，再代入各解的方程：
\[\begin{aligned}
W'&=y_1y_2''-y_1''y_2\\
&=y_1(-py_2'-qy_2)-(-py_1'-qy_1)y_2\\
&=-p(y_1y_2'-y_1'y_2)=-pW.
\end{aligned}\]
令$P(t)=\int_{t_0}^tp(s)\dd s$，则$(\ee^PW)'=\ee^P(W'+pW)=0$。从$t_0$积分，得到
\[W(t)=W(t_0)\exp\!\left[-\int_{t_0}^tp(s)\dd s\right].\]
只要系数在同一连通区间连续，$W$一旦在某点非零，就在整段区间非零，两个解线性独立。

若已知非零解$y_1$，在$y_1\ne0$的区间令$y=y_1u$。两次求导为
\[y'=y_1'u+y_1u',\qquad
y''=y_1''u+2y_1'u'+y_1u''.\]
代入原方程，含$u$的系数因$y_1$满足齐次方程而消失，余下
\[y_1u''+(2y_1'+py_1)u'=0.\]
把$w=u'$视作未知函数。积分因子为
\[\exp\!\left(\int\left[2\frac{y_1'}{y_1}+p\right]\dd t\right)
=y_1^2\ee^{\int p\dd t},\]
其中$\int y_1'/y_1\dd t=\ln|y_1|$，常数因子可约去。因此
$(y_1^2\ee^{\int p\dd t}w)'=0$，即$w=C\ee^{-\int p\dd t}/y_1^2$。再积分，
\[y=y_1\left(C_0+C\int\frac{\ee^{-\int p\dd t}}{y_1^2}\dd t\right).
\]
取$C=1,C_0=0$便得到第二个独立解。$C_0y_1$属于已有解，不是第三个自由方向。该公式解释了重根的时间因子，也适用于某些变系数方程；$y_1$有零点时先分区间使用，再依据原方程系数的正则性讨论延拓。

\subsection{稳态、瞬态与共振极限}
对阻尼振子，齐次解由初态决定并随时间衰减，周期特解由驱动频率决定。只有在衰减后的长时间区间内，才可把完整响应近似为稳态。弱阻尼时这个等待时间可以很长。先令阻尼趋零与先令时间趋于无穷，通常不能交换：前者保留自由振动和共振增长，后者先抹去瞬态。

无阻尼近共振响应含$[\cos\Omega t-\cos\omega_0t]/(\omega_0^2-\Omega^2)$。把分子化为两个正弦之积，可看出拍频；在固定$t$下令$\Omega\to\omega_0$则得到$t\sin\omega_0t$。频率差很小时，若时间已长到其倒数量级，就不能再把拍包络近似成线性增长。

\subsection{耗散、共振与频率响应的两种峰}
对质量$m$的阻尼振子，将方程乘$m\dot x$，得到
\[\frac{\dd}{\dd t}\left(\tfrac12m\dot x^2+\tfrac12m\omega_0^2x^2\right)
=F(t)\dot x-2m\gamma\dot x^2.\]
左侧为机械能变化率，右侧为外力输入功率减去黏性耗散。自由运动时能量单调减小，但位置不一定单调；受驱运动时，瞬时输入功率可正可负，只有周期平均才在稳态下与平均耗散相等。

设稳态位移为$X\cos(\Omega t-\delta)$。复指数计算给出响应分母$\omega_0^2-\Omega^2+2i\gamma\Omega$，其模决定振幅，其辐角决定相位滞后。低频时位移近似跟随力，高频时位移与力近乎反相，$\Omega=\omega_0$时相差四分之一周期。用普通反正切而不确定象限，会把高频区相位写错。

位移峰最大化$X$，功率峰最大化$\Omega^2X^2$，两个目标函数不同。因此“共振频率”须注明观察量；习题~\ref{ex:c2-h8}和~\ref{ex:c2-h9}分别求这两种峰。自由运动的欠阻尼条件也不是频响有内部位移峰的条件，不能把判别式的临界值直接套用到频率最优化。

近共振与精确共振的时间范围同样要说明。频率差$\Delta$很小时，响应有快振荡与慢包络；只有$|\Delta|t\ll1$时慢包络可线性展开，才近似精确共振的$t$增长。若观测到$t\sim1/|\Delta|$，拍的回落已不可忽略。零阻尼极限与长时间极限改变的是不同过程，取极限的先后顺序应与问题的观测条件相符。


\section{习题与解答}\label{sec:c2-exercises}

\Needspace{8\baselineskip}
\begin{exercise}[变截面水缸]\label{ex:c2-h1}
水缸的水平截面积随高度变为$A(h)=A_0(1+h/H)$，初水深$H$，底孔面积$S\ll A_0$。沿用理想准稳态出流模型，求排空时间，并与截面积恒为$A_0$的水缸比较。 见图~\ref{fig:c2-h1-q}。
\end{exercise}

\begin{figure}[H]\centering
\begin{tikzpicture}[>=Stealth,line cap=round,line join=round]
\draw[thick](-1,0)--(-1.5,3) (1,0)--(1.5,3);\draw[thick](-1,0)--(.75,0);\draw[curve](-1.4,2.4)--(1.4,2.4);\draw[axis](1,0)--(2,0)node[right]{$v$};\node at(0,1.3){$A(h)$};\draw[<->](-1.9,0)--(-1.9,2.4)node[midway,left]{$H$};
\end{tikzpicture}
\caption{缸体向上扩张，图示侧形只表达截面积随高度增加。}\label{fig:c2-h1-q}
\end{figure}
\begin{solution}
图~\ref{fig:c2-h1-q}中缸体上部较宽，液面下降$\dd h$所排出的体积为$-A(h)\dd h$。连续性方程为$A(h)\dot h=-S\sqrt{2gh}$。图~\ref{fig:c2-h1-s}按水深积分求总时间，得到
\[T=\frac{A_0}{S\sqrt{2g}}\int_0^H\frac{1+h/H}{\sqrt h}\dd h.
\]
对 $h^{-1/2}+h^{1/2}/H$ 逐项积分，
\[
T=\frac{A_0}{S\sqrt{2g}}\left[2\sqrt h+\frac{2h^{3/2}}{3H}\right]_0^H
=\frac{A_0\sqrt H}{S\sqrt{2g}}\left(2+\frac23\right),
\]
所以
\[\boxed{T=\frac43\frac{A_0}{S}\sqrt{\frac{2H}g}}.
\]
恒截面模型只含第一个积分，所以时间比为$(2+2/3)/2=4/3$。图~\ref{fig:c2-h1-s}的时间权重$h^{-1/2}$在低水深处更大；高处新增面积所占权重较小，不能用平均截面积$3A_0/2$直接乘原时间。若要求隐式水深关系，积分从$h$到$H$可写为
\[t=\frac{A_0}{S\sqrt{2g}}\left[2(\sqrt H-\sqrt h)+\frac{2}{3H}(H^{3/2}-h^{3/2})\right].
\]
这已给出单调的时间--水深对应，无需为求显式根而引入不必要的高次方程。


隐式 $h(t)$ 式在 $h=H$ 时给 $t=0$，在 $h=0$ 时给 $t=T$。对它关于 $h$ 求导，
\[
\frac{\dd t}{\dd h}=\frac{A_0}{S\sqrt{2g}}\left[-h^{-1/2}-\frac{h^{1/2}}H\right]
=-\frac{A(h)}{S\sqrt{2gh}}.
\]
在 $h>0$ 的下降分支取倒数，恢复 $\dot h=-S\sqrt{2gh}/A(h)$，所以隐式关系已包含整个下降过程，求排空时间无需再解三次方程。

\begin{figure}[H]\centering
\begin{tikzpicture}[>=Stealth,line cap=round,line join=round]
\begin{scope}[x=2.39316cm,y=1.00000cm]\draw[axis](0,0)--(1.25,0)node[right]{$h/H$};\draw[axis](0,0)--(0,4.5)node[above]{$(1+h/H)/\sqrt{h/H}$};\draw[curve,domain=0.08:1,samples=101]plot(\x,{(1+\x)/sqrt(\x)});\node[right]at(1.5,2){时间积分的权重};\end{scope}
\end{tikzpicture}
\caption{单位水深对应的时间权重在低水深处增大，但积分仍收敛。}\label{fig:c2-h1-s}
\end{figure}
\end{solution}

\Needspace{8\baselineskip}
\begin{exercise}[线性阻力下落]\label{ex:c2-h2}
质点由静止下落，向下为正，重力$mg$，阻力为$-bv$，$b>0$。求$v(t)$、下落距离$y(t)$和终端速度，并讨论短时极限。 见图~\ref{fig:c2-h2-q}。
\end{exercise}

\begin{figure}[H]\centering
\begin{tikzpicture}[>=Stealth,line cap=round,line join=round]
\draw[axis](0,2.5)--(0,-.3)node[below]{$y$};\fill(1.3,1.3)circle(2pt)node[right]{$m$};\draw[axis](1.3,1.3)--(1.3,.1)node[below]{$mg$};\draw[axis](1.3,1.3)--(1.3,2.2)node[above]{$bv$};
\end{tikzpicture}
\caption{阻力在质点向下运动时向上，坐标正向向下。}\label{fig:c2-h2-q}
\end{figure}
\begin{solution}
向下为正，重力项为正，线性阻力项在下落中为负。图~\ref{fig:c2-h2-q}的两个力反向，初始阻力为零，因此初始加速度仍为$g$。先求一阶速度方程，再从速度积分出位移，两个初态分别用于两步计算。


方程 $m\dot v=mg-bv$ 是一阶非齐次线性方程。记 $\tau=m/b$、$v_T=mg/b=g\tau$，乘以积分因子 $\ee^{t/\tau}$，按乘积法则有
\[
\frac{\dd}{\dd t}(\ee^{t/\tau}v)=\ee^{t/\tau}\dot v+\frac1\tau\ee^{t/\tau}v=g\ee^{t/\tau}.
\]
从 $0$ 积到 $t$，初值 $v(0)=0$ 给出 $\ee^{t/\tau}v=g[\tau\ee^{s/\tau}]_0^t=g\tau(\ee^{t/\tau}-1)$，所以
\[\boxed{v(t)=v_T(1-\ee^{-t/\tau})}.
\]
再由 $y(0)=0$ 积分，
\[
y(t)=v_T\int_0^t(1-\ee^{-s/\tau})\dd s
=v_T[s+\tau\ee^{-s/\tau}]_0^t,
\]
得到
\[\boxed{y(t)=v_T\{t-\tau(1-\ee^{-t/\tau})\}}.
\]
图~\ref{fig:c2-h2-s}中速度渐近水平线$v_T$，到达有限终端速度不等于在有限时间内严格取到它。短时令$s=t/\tau$，先展开$\ee^{-s}=1-s+s^2/2-s^3/6+O(s^4)$，所以$1-\ee^{-s}=s-s^2/2+s^3/6+O(s^4)$。速度乘以$v_T=g\tau$，位移中的两个一次项$s-(1-\ee^{-s})$则相消，得到
\[v=gt-\frac{g}{2\tau}t^2+O(t^3),\qquad
y=\frac12gt^2-\frac{g}{6\tau}t^3+O(t^4).
\]
初始速度与位移分别从一次、二次开始，阻力修正晚一阶出现。长时间则$y\simeq v_T(t-\tau)$；把这一渐近直线直接延长到$t=0$，会违反初态。


位移式中的常数项保证 $y(0)=0$，不能把指数的原函数只取上端值而漏掉下端。

终端速度由$mg-bv_T=0$得到；图~\ref{fig:c2-h2-s}中速度差$v_T-v=v_T\ee^{-t/\tau}$在有限时间始终为正。若希望达到终速的$1-\epsilon$倍，需要$t=\tau\ln(1/\epsilon)$。这给出接近终速的实际时间尺度，也说明严格到达终速需要无限时间。

\begin{figure}[H]\centering
\begin{tikzpicture}[>=Stealth,line cap=round,line join=round]
\begin{scope}[x=1.00000cm,y=1.38462cm]\draw[axis](0,0)--(4.25,0)node[right]{$t/\tau$};\draw[axis](0,0)--(0,1.3)node[above]{$v/v_T$};\draw[curve,domain=0:4,samples=101]plot(\x,{1-exp(-\x)});\draw[guide](0,1)--(4,1);\node[right]at(4,1){$1$};\end{scope}
\end{tikzpicture}
\caption{速度单调趋近终端值，原点斜率对应初始加速度$g$。}\label{fig:c2-h2-s}
\end{figure}
\end{solution}

\Needspace{8\baselineskip}
\begin{exercise}[平方阻力与双曲函数]\label{ex:c2-h3}
质点向下运动时满足$m\dot v=mg-cv^2$，$c>0$。求从静止开始的速度与距离，并给出初速$v_i>v_T=\sqrt{mg/c}$时的速度分支。 见图~\ref{fig:c2-h3-q}。
\end{exercise}

\begin{figure}[H]\centering
\begin{tikzpicture}[>=Stealth,line cap=round,line join=round]
\begin{scope}[x=1.69697cm,y=1.00000cm]\draw[axis](0,0)--(1.65,0)node[right]{$v/v_T$};\draw[axis](0,-1.2)--(0,1.3)node[above]{$\dot v/g$};\draw[curve,domain=0:1.4,samples=101]plot(\x,{1-(\x)^2});\end{scope}
\end{tikzpicture}
\caption{相线在$v=v_T$处变号；高于终端速率的粒子减速。}\label{fig:c2-h3-q}
\end{figure}
\begin{solution}
平方阻力使速度方程非线性，但因右侧只含$v$，仍可分离。平衡值$v_T$把允许向下速度分成两侧：低于它时加速，高于它时减速。图~\ref{fig:c2-h3-q}的加速度曲线在$v_T$变号，反双曲函数的分支应与这个区间划分一致。


令 $u=v/v_T$、$\tau=v_T/g$，方程变成 $\dot u=(1-u^2)/\tau$。先分解 $1/(1-u^2)=\tfrac12[1/(1+u)+1/(1-u)]$，原函数为 $\tfrac12\ln|(1+u)/(1-u)|$。静止释放对应 $0\le u<1$，从初态积分，
\[
\int_0^{u(t)}\frac{\dd w}{1-w^2}=\frac1\tau\int_0^t\dd s,\qquad
\left[\frac12\ln\frac{1+w}{1-w}\right]_0^{u(t)}=\frac t\tau.
\]
左侧是 $\artanh u$，所以
\[\artanh u=t/\tau,\qquad \boxed{v=v_T\tanh(t/\tau)}.
\]
求距离时令 $s=t/\tau$，并用 $\tanh s=\sinh s/\cosh s$。分子正是分母的导数，因此
\[
y(t)=v_T\tau\int_0^{t/\tau}\frac{\sinh s}{\cosh s}\dd s
=v_T\tau[\ln\cosh s]_0^{t/\tau},
\]
利用 $\cosh0=1$，得到
\[\boxed{y=\frac{v_T^2}{g}\ln\cosh(t/\tau)}.
\]
若 $v_i>v_T$，不能继续用实数 $\artanh(v_i/v_T)$。此时令 $u=\coth z$，由 $\dd(\coth z)/\dd z=1-\coth^2z$，代入原方程得 $(1-u^2)\dot z=(1-u^2)/\tau$。因 $u>1$，可约去非零因子，得 $\dot z=1/\tau$、$z=z_0+t/\tau$，初值确定 $z_0=\operatorname{arccoth}(v_i/v_T)$，故
\[\boxed{v(t)=v_T\coth\left(\frac t\tau+\operatorname{arccoth}\frac{v_i}{v_T}\right)}.
\]
图~\ref{fig:c2-h3-s}中两支分别从终端值下方和上方趋近同一水平线，$v=v_T$本身也是平衡解。平方阻力的方向应写成$-cv|v|$；本题限定向下，才可化为$-cv^2$。若运动中允许反向，必须相应更换符号，不能用同一个二次式跨过$v=0$解释上抛全过程。


上方分支的积分常数为
\[z_0=\operatorname{arccoth}(v_i/v_T)
=\tfrac12\ln\frac{v_i+v_T}{v_i-v_T}>0.\]
若也从 $y(0)=0$ 计距离，令 $w=t/\tau+z_0$，则
\[
y(t)=v_T\tau\int_{z_0}^{t/\tau+z_0}\frac{\cosh w}{\sinh w}\dd w
=v_T\tau[\ln\sinh w]_{z_0}^{t/\tau+z_0}
=v_T\tau\ln\frac{\sinh(t/\tau+z_0)}{\sinh z_0}.
\]
求导可恢复$coth$速度；$t=0$时距离为零。图中的上、下两支分别满足$\dot v<0$与$\dot v>0$，它们以平衡解$v=v_T$为界，各自使用相应的反双曲分支。

同一实对数原函数在 $0\le u<1$ 时是 $\artanh u$，在 $u>1$ 时用 $\operatorname{arccoth}u$ 表示。前者反解为 $\tanh$，后者反解为 $\coth$，不能用同一个实反函数越过 $u=1$。

静止分支的短时距离还须展开对数与双曲余弦的复合函数。先用$\cosh s=1+s^2/2+s^4/24+O(s^6)$，记$z=\cosh s-1$，再用$\ln(1+z)=z-z^2/2+O(z^3)$，得到
\[
 \ln\cosh s=\frac{s^2}{2}
 +\left(\frac1{24}-\frac18\right)s^4+O(s^6)
 =\frac{s^2}{2}-\frac{s^4}{12}+O(s^6).
\]
代入$s=t/\tau=gt/v_T$及$y=(v_T^2/g)\ln\cosh s$，有
\[
 y=\frac12gt^2-\frac{g^3t^4}{12v_T^2}
 +O\!\left(\frac{g^5t^6}{v_T^4}\right).
\]
阻力首个修正比线性阻力情形晚一阶，因为本题初期$v^2\sim g^2t^2$。图~\ref{fig:c2-h3-s}的两支在长时趋于同一终端值，都不在有限时刻穿过它。

\begin{figure}[H]\centering
\begin{tikzpicture}[>=Stealth,line cap=round,line join=round]
\begin{scope}[x=1.00000cm,y=1.00000cm]\draw[axis](0,0)--(4.25,0)node[right]{$t/\tau$};\draw[axis](0,0)--(0,2)node[above]{$v/v_T$};\draw[curve,domain=0:4,samples=101]plot(\x,{tanh(\x)});\draw[dashed,domain=0:4,samples=101]plot(\x,{1/tanh(\x+.7)});\draw[guide](0,1)--(4,1);\end{scope}
\end{tikzpicture}
\caption{实线从静止趋近终端值，虚线从更大初速衰减至终端值。}\label{fig:c2-h3-s}
\end{figure}
\end{solution}

\Needspace{8\baselineskip}
\begin{exercise}[电容充电的时间边界]\label{ex:c2-h4}
串联$R,C$电路在$t=0$接入恒电压$U$，电容初电荷为$q_0$。取电源正极向电容正极的电流为正，求$q(t)$、$I(t)$，并计算从未充电状态到充满时电阻耗散的能量。 见图~\ref{fig:c2-h4-q}。
\end{exercise}

\begin{figure}[H]\centering
\begin{tikzpicture}[>=Stealth,line cap=round,line join=round]
\draw[curve](0,0)--(0,.95) (0,1.45)--(0,2.4)--(1,2.4);\draw[thick](-.35,1.45)--(.35,1.45) (-.2,.95)--(.2,.95);\node[left]at(-.35,1.5){$+$};\node[left]at(-.35,.9){$-$};\draw(1,2.2)rectangle(2,2.6);\draw[curve](2,2.4)--(3.3,2.4)--(3.3,1.5) (3.3,1.1)--(3.3,0)--(0,0);\draw[thick](2.9,1.5)--(3.7,1.5) (2.9,1.1)--(3.7,1.1);\node at(1.5,2.9){$R$};\node[right]at(3.7,1.3){$C$};\node[left]at(-.5,1.2){$U$};\draw[axis](.2,2.9)--(.8,2.9)node[midway,above]{$I$};
\end{tikzpicture}
\caption{正电流增加电容正极板的电荷。}\label{fig:c2-h4-q}
\end{figure}
\begin{solution}
电容电压是$q/C$，电源与它的电压差驱动电阻电流。图~\ref{fig:c2-h4-q}规定正电流流向电容正极，故回路方程为$R\dot q=U-q/C$。初电荷可以大于最终值，这时同一方程自然给出反向电流，不必另改符号约定。


回路方程为 $RI+q/C=U$，$I=\dot q$。令 $\tau=RC$、$z=q-CU$，则 $\dot z+z/\tau=0$。乘以积分因子，$(\ee^{t/\tau}z)'=0$；从 $0$ 积到 $t$，$\ee^{t/\tau}z(t)=q_0-CU$。因此 $q=CU+(q_0-CU)\ee^{-t/\tau}$，求导时指数产生 $-1/\tau$，得到
\[\boxed{q(t)=CU+(q_0-CU)\ee^{-t/\tau},\qquad
I(t)=\frac{CU-q_0}{RC}\ee^{-t/\tau}}.
\]
图~\ref{fig:c2-h4-s}取$q_0=0$，电荷渐近$CU$，电流由$U/R$衰减至零。若$q_0>CU$，初始电流反向，公式无需更换分支；电流正向定义并不等于每个初态下电流都为正。

未充电时，电阻总耗散为
\[W_R=\int_0^\infty RI^2\dd t
=\frac{U^2}{R}\left[-\frac{RC}{2}\ee^{-2t/(RC)}\right]_0^\infty
=\boxed{\frac12CU^2}.
\]
电源做功 $U\int_0^\infty I\dd t=U[q(t)]_0^\infty=CU^2$，其余 $CU^2/2$ 储于电容。理想模型里，耗散总量与 $R$ 无关，电阻只改变过程的快慢；将 $R\to0$ 与忽略所有电磁暂态同时使用，则超出了该单纯 RC 模型。


偏差 $z=q-CU$ 的指数衰减说明，初态高于或低于平衡值只改变偏差的符号，不改变时间常数。

一般初电荷下，电阻耗散为
\[W_R=R\left(\frac{CU-q_0}{RC}\right)^2\int_0^\infty\ee^{-2t/(RC)}\dd t
=R\left(\frac{CU-q_0}{RC}\right)^2\frac{RC}{2}
=\frac{(CU-q_0)^2}{2C}.\]
电源做功为$U(CU-q_0)$，电容能量增加为$[(CU)^2-q_0^2]/(2C)$，两者之差正好是上式。图~\ref{fig:c2-h4-s}选择$q_0=0$，于是输入能量的一半储能、另一半耗散；这一比例不应套用于任意$q_0$。

\begin{figure}[H]\centering
\begin{tikzpicture}[>=Stealth,line cap=round,line join=round]
\begin{scope}[x=1.00000cm,y=1.38462cm]\draw[axis](0,0)--(4.25,0)node[right]{$t/(RC)$};\draw[axis](0,0)--(0,1.3)node[above]{$q/(CU)$};\draw[curve,domain=0:4,samples=101]plot(\x,{1-exp(-\x)});\draw[dashed,domain=0:4,samples=101]plot(\x,{exp(-\x)});\node[right]at(4,.5){虚线：$IR/U$};\end{scope}
\end{tikzpicture}
\caption{未充电初态下的电荷与电流，二者具有同一时间常数。}\label{fig:c2-h4-s}
\end{figure}
\end{solution}

\Needspace{8\baselineskip}
\begin{exercise}[是否在有限时间到达零]\label{ex:c2-h5}
非负量$y(t)$满足$\dot y=-ky^\nu$，$k>0,\nu>0$，$y(0)=y_0>0$。分类讨论到达$y=0$的时间，并说明分离变量时遗漏了什么解。 见图~\ref{fig:c2-h5-q}。
\end{exercise}

\begin{figure}[H]\centering
\begin{tikzpicture}[>=Stealth,line cap=round,line join=round]
\begin{scope}[x=1.24444cm,y=1.00000cm]\draw[axis](0,0)--(2.25,0)node[right]{$y$};\draw[axis](0,-1.8)--(0,0.3)node[above]{$\dot y/k$};\draw[curve,domain=0:2,samples=101]plot(\x,{-(\x)^.5});\end{scope}
\end{tikzpicture}
\caption{图示$\nu=1/2$，在零点处右端函数的导数不有界。}\label{fig:c2-h5-q}
\end{figure}
\begin{solution}
在$y>0$阶段可除以$y^\nu$分离变量，直到状态到达$y=0$为止。图~\ref{fig:c2-h5-q}中变化率随$y$趋零；到达所需时间是否有限，要由相应积分的收敛性判断。


对 $\nu\ne1$，在正值分支上分离变量并代入初态，
\[
\int_{y_0}^{y(t)}u^{-\nu}\dd u=-k\int_0^t\dd s,\qquad
\left[\frac{u^{1-\nu}}{1-\nu}\right]_{y_0}^{y(t)}=-kt.
\]
乘以 $1-\nu$ 并移项，得到
\[y^{1-\nu}=y_0^{1-\nu}-(1-\nu)kt.
\]
若$0<\nu<1$，右边在$T=y_0^{1-\nu}/[(1-\nu)k]$变零；到达前
\[y(t)=[y_0^{1-\nu}-(1-\nu)kt]^{1/(1-\nu)}.
\]
由于模型限定 $y\ge0$ 且随时间不增加，可在 $t\ge T$ 接上 $y=0$。若 $\nu=1$，分离变量的积分改为 $[\ln u]_{y_0}^{y(t)}=-kt$，取指数得 $y=y_0\ee^{-kt}$，只在无穷远趋零。若 $\nu>1$，将正值关系写为 $y^{-(\nu-1)}=y_0^{-(\nu-1)}+(\nu-1)kt$，两边取 $-1/(\nu-1)$ 次幂，得到
\[y(t)=\frac{y_0}{[1+(\nu-1)ky_0^{\nu-1}t]^{1/(\nu-1)}},
\]
仍在无限时间趋零。图~\ref{fig:c2-h5-s}比较这些过程，有限到达由$\int_0^{y_0}y^{-\nu}\dd y$是否收敛决定。

分离时除以$y^\nu$，会排除恒为零的平衡解。对于$0<\nu<1$，零点附近不满足常用的局部Lipschitz条件，讨论从零开始的反向延拓时还须另行说明唯一性；正值衰减分支及本题的非负终止规则则已明确。


用时也可直接写成下端截断的极限：
\[
T=\frac1k\lim_{\delta\downarrow0}\int_\delta^{y_0}y^{-\nu}\dd y
=\frac1k\lim_{\delta\downarrow0}
\begin{cases}[y^{1-\nu}/(1-\nu)]_\delta^{y_0},&\nu\ne1,\\
[\ln y]_\delta^{y_0},&\nu=1.
\end{cases}
\]
$0<\nu<1$ 时为前面的有限值，$\nu=1$ 时对数发散，$\nu>1$ 时幂次发散。

对于有限到达分支，近$T$时$y=[(1-\nu)k(T-t)]^{1/(1-\nu)}$。因为指数大于$1$，导数在$T^-$也趋零，接恒零段后函数的一阶导连续，仍满足原方程。图~\ref{fig:c2-h5-s}中的虚线正体现这一平滑接续；把正值公式继续到使括号为负的时间，则会违反非负状态的限制。

\begin{figure}[H]\centering
\begin{tikzpicture}[>=Stealth,line cap=round,line join=round]
\begin{scope}[x=1.1cm,y=2cm]
\draw[axis](0,0)--(4.3,0)node[right]{$\tau$};\draw[axis](0,0)--(0,1.3)node[above]{$y/y_0$};
\draw[curve,domain=0:4,samples=101]plot(\x,{exp(-\x)});
\draw[dashed,domain=0:2,samples=101]plot(\x,{(1-.5*\x)^2});\draw[dashed](2,0)--(4,0);
\draw[dotted,domain=0:4,samples=101]plot(\x,{1/(1+\x)});
\node[below]at(2,0){$2$};
\node[align=left,right]at(4.55,.7){实线：$\nu=1$\\虚线：$\nu=1/2$\\点线：$\nu=2$};\end{scope}
\end{tikzpicture}
\caption{每条曲线均取无量纲时间$\tau=ky_0^{\nu-1}t$。虚线在$\tau=2$到达零并接恒零段，其余两条只渐近趋零。}\label{fig:c2-h5-s}
\end{figure}
\end{solution}

\Needspace{8\baselineskip}
\begin{exercise}[临界阻尼的最大位移]\label{ex:c2-h6}
临界阻尼振子由$x(0)=0,\dot x(0)=v_0>0$启动。求最大位移、到达时刻以及该时刻剩余机械能相对于初始能量的比例。 见图~\ref{fig:c2-h6-q}。
\end{exercise}

\begin{figure}[H]\centering
\begin{tikzpicture}[>=Stealth,line cap=round,line join=round]
\begin{scope}[x=1.00000cm,y=3.00000cm]\draw[axis](0,0)--(4.25,0)node[right]{$\omega_0t$};\draw[axis](0,0)--(0,0.6)node[above]{$\omega_0x/v_0$};\draw[curve,domain=0:4,samples=101]plot(\x,{\x*exp(-\x)});\end{scope}
\end{tikzpicture}
\caption{位置只出现一个正的最大值，之后从右侧回到平衡。}\label{fig:c2-h6-q}
\end{figure}
\begin{solution}
临界阻尼解含$t\ee^{-\omega_0t}$，一开始随$t$增长，随后指数衰减占优，所以只有一次正峰。图~\ref{fig:c2-h6-q}的峰对应速度为零；这一瞬间机械能全为势能，但此前阻尼已经耗散了一部分初能。


由$x=v_0t\ee^{-\omega_0t}$，速度为$\dot x=v_0(1-\omega_0t)\ee^{-\omega_0t}$。因此图~\ref{fig:c2-h6-s}中的唯一正转向时刻为$t_*=1/\omega_0$，
\[\boxed{x_{\max}=\frac{v_0}{\ee\omega_0}}.
\]
此刻速度为零，机械能全为弹性势能，
\[E_* =\frac12m\omega_0^2x_{\max}^2=\frac12mv_0^2\ee^{-2},
\qquad \boxed{E_*/E_0=\ee^{-2}}.
\]
阻尼在运动期间耗能，因而不能把初始动能全部等同于转向点势能。由$\dot E=-2m\omega_0\dot x^2$积分也能得到相同损失。这里的单次转向并不意味着振子周期运动：转向后$x(t)$仍为正，只在$t\to\infty$趋零。


能量损失也可算成完整积分。设$s=\omega_0t$，则
\[\frac{E(t)}{E_0}=\ee^{-2s}\{(1-s)^2+s^2\}.\]
其中第一项来自速度平方，第二项来自位移平方。在图中$s=1$的峰点，第一项为零。另一方面，将耗散功率$2m\omega_0\dot x^2$积分后除以初能$E_0=mv_0^2/2$，再用$\dd t=\dd s/\omega_0$换元，得到
\[
\frac{W_{\rm diss}(0,t_*)}{E_0}=4\int_0^1(1-2s+s^2)\ee^{-2s}\dd s.
\]
积分中的一次、二次多项式分别作分部积分：
\[
\int s\ee^{-2s}\dd s=-\frac{s}{2}\ee^{-2s}+\frac12\int\ee^{-2s}\dd s
=-\left(\frac s2+\frac14\right)\ee^{-2s}+C,
\]
\[
\int s^2\ee^{-2s}\dd s=-\frac{s^2}{2}\ee^{-2s}+\int s\ee^{-2s}\dd s
=-\left(\frac{s^2}{2}+\frac s2+\frac14\right)\ee^{-2s}+C.
\]
组合三项原函数并代入上下限，得到
\[
\frac{W_{\rm diss}(0,t_*)}{E_0}
=[(-2s^2+2s-1)\ee^{-2s}]_0^1=1-\ee^{-2}.
\]
峰点剩余机械能与此前耗散之和等于初始能量。

把时间无量纲化为$s=\omega_0t$，位移为$x=(v_0/\omega_0)s\ee^{-s}$。导数符号由$1-s$决定，$0<s<1$时正、$s>1$时负，因此峰点不只是局部驻点，也是正时间轴上的全局最大值。

\begin{figure}[H]\centering
\begin{tikzpicture}[>=Stealth,line cap=round,line join=round]
\begin{scope}[x=1.00000cm,y=3.00000cm]\draw[axis](0,0)--(4.25,0)node[right]{$\omega_0t$};\draw[axis](0,0)--(0,0.6)node[above]{$\omega_0x/v_0$};\draw[curve,domain=0:4,samples=101]plot(\x,{\x*exp(-\x)});\fill(1,.3678794)circle(1.5pt)node[above]{$E$};\draw[guide](1,0)--(1,.3678794);\node[below]at(1,0){$1$};\end{scope}
\end{tikzpicture}
\caption{最大位移发生于$\omega_0t=1$，高度为$1/\ee$。}\label{fig:c2-h6-s}
\end{figure}
\end{solution}

\Needspace{8\baselineskip}
\begin{exercise}[过阻尼运动的转向时刻]\label{ex:c2-h7}
取$\gamma>\omega_0$，振子由$x(0)=0,\dot x(0)=v_0>0$启动。求位移唯一的正极大值时刻，并说明临界阻尼极限。 见图~\ref{fig:c2-h7-q}。
\end{exercise}

\begin{figure}[H]\centering
\begin{tikzpicture}[>=Stealth,line cap=round,line join=round]
\begin{scope}[x=1.00000cm,y=3.00000cm]\draw[axis](0,0)--(5.25,0)node[right]{$\omega_0t$};\draw[axis](0,0)--(0,0.6)node[above]{$\omega_0x/v_0$};\draw[curve,domain=0:5,samples=101]plot(\x,{exp(-1.25*\x)*sinh(.75*\x)/.75});\end{scope}
\end{tikzpicture}
\caption{示意取$\gamma/\omega_0=1.25$。}\label{fig:c2-h7-q}
\end{figure}
\begin{solution}
过阻尼运动由两个不同的衰减指数叠加。图~\ref{fig:c2-h7-q}中的曲线初始斜率为正，长时回到零；峰点须保留两个指数共同求导，单个慢模态只决定长时行为。


令 $\delta=\sqrt{\gamma^2-\omega_0^2}$。过阻尼初值解为 $x=(v_0/\delta)\ee^{-\gamma t}\sinh\delta t$，乘积求导的两项是 $-\gamma\ee^{-\gamma t}\sinh\delta t$ 和 $\delta\ee^{-\gamma t}\cosh\delta t$，故
\[\dot x=\frac{v_0}{\delta}\ee^{-\gamma t}
 [\delta\cosh(\delta t)-\gamma\sinh(\delta t)].\]
图~\ref{fig:c2-h7-s}的水平切线满足$\tanh(\delta t_*)=\delta/\gamma$。由于$0<\delta/\gamma<1$，而$\tanh$在正半轴严格增加，该方程只有一个解：
\[\boxed{t_* =\frac1\delta\artanh\frac\delta\gamma
 =\frac1{2\delta}\ln\frac{\gamma+\delta}{\gamma-\delta}}.\]
为计算峰值，先将双曲函数化成参数。记 $z=\delta t_*$，则 $\sinh z=(\delta/\gamma)\cosh z$；代入 $\cosh^2z-\sinh^2z=1$，得 $\cosh^2z(1-\delta^2/\gamma^2)=1$。因为 $z>0$、$\gamma^2-\delta^2=\omega_0^2$，取正根有
\[
\cosh z=\frac\gamma{\omega_0},\qquad \sinh z=\frac\delta{\omega_0}.
\]
再代回 $x(t_*)=(v_0/\delta)\ee^{-\gamma t_*}\sinh z$，得到
\[\boxed{x_* =\frac{v_0}{\omega_0}\ee^{-\gamma t_*}}.\]
$t<t_*$时速度为正，之后为负；位置始终为正，故这是全局最大值。令$\delta\to0$，用$\artanh u=u+O(u^3)$，得到$t_*\to1/\omega_0$与$x_*\to v_0/(\ee\omega_0)$，连续接上临界阻尼结果。

强阻尼下 $\delta=\gamma\sqrt{1-\omega_0^2/\gamma^2}\simeq\gamma-\omega_0^2/(2\gamma)$。用 $(\gamma+\delta)(\gamma-\delta)=\omega_0^2$，时间式改写为 $t_*=\ln[(\gamma+\delta)/\omega_0]/\delta$，于是 $t_*\simeq\ln(2\gamma/\omega_0)/\gamma$。较大的阻尼可以让第一次转向更早，却使随后靠近原点的慢模态衰减率 $\gamma-\delta$ 更小；第一次转向时刻与长时间弛豫率应分别比较。


图~\ref{fig:c2-h7-s}的水平切线前后分别对应正、负速度。临界极限中$\delta$趋零，但$\artanh(\delta/\gamma)$也同阶趋零，两者的比值有限；因此必须一起取极限，而不能直接把$1/\delta$判为发散。

\begin{figure}[H]\centering
\begin{tikzpicture}[>=Stealth,line cap=round,line join=round]
\begin{scope}[x=1.00000cm,y=3.00000cm]\draw[axis](0,0)--(5.25,0)node[right]{$\omega_0t$};\draw[axis](0,0)--(0,0.6)node[above]{$\omega_0x/v_0$};\draw[curve,domain=0:5,samples=101]plot(\x,{exp(-1.25*\x)*sinh(.75*\x)/.75});\draw[guide](.9242,0)--(.9242,.31498);\draw[dashed](.3,.31498)--(1.65,.31498);\node[below]at(.9242,0){$\omega_0t_*$};\end{scope}
\end{tikzpicture}
\caption{正极大值的切线水平；运动不会穿越平衡点。}\label{fig:c2-h7-s}
\end{figure}
\end{solution}

\Needspace{8\baselineskip}
\begin{exercise}[位移共振峰是否存在]\label{ex:c2-h8}
受迫阻尼振子中$\gamma>0$，固定驱动力幅度$F_0$，改变频率$\Omega\ge0$。求稳态位移振幅的最大值，并讨论没有内部共振峰的情形。 见图~\ref{fig:c2-h8-q}。
\end{exercise}

\begin{figure}[H]\centering
\begin{tikzpicture}[>=Stealth,line cap=round,line join=round]
\begin{scope}[x=1.01818cm,y=1.00000cm]\draw[axis](0,0)--(2.75,0)node[right]{$\Omega/\omega_0$};\draw[axis](0,0)--(0,3)node[above]{$m\omega_0^2X/F_0$};\draw[curve,domain=0:2.5,samples=101]plot(\x,{1/sqrt((1-(\x)^2)^2+.16*(\x)^2)});\end{scope}
\end{tikzpicture}
\caption{图示$\gamma/\omega_0=0.2$的位移振幅。}\label{fig:c2-h8-q}
\end{figure}
\begin{solution}
驱动力幅度固定时，最大位移振幅对应响应分母的最小值。驱动频率限制为非负，图~\ref{fig:c2-h8-q}的内部峰可能随阻尼增加移到端点并消失。求极值时需同时处理内部驻点与零频端点。


振幅的平方与
\[D(\Omega)=(\omega_0^2-\Omega^2)^2+4\gamma^2\Omega^2\]
成反比。求极值可改为求$D$在非负半轴的最小值：
\[D'=2(\omega_0^2-\Omega^2)(-2\Omega)+8\gamma^2\Omega
=4\Omega(\Omega^2-\omega_0^2+2\gamma^2).\]
令$u=\Omega^2$，先把分母配方：
\[D=(u-\omega_0^2+2\gamma^2)^2+4\gamma^2(\omega_0^2-\gamma^2),\qquad u\ge0.\]
当$0<\gamma<\omega_0/\sqrt2$时，顶点$u_r=\omega_0^2-2\gamma^2$位于可行半轴内。将它代回，
\[D(\Omega_r)=4\gamma^4+4\gamma^2(\omega_0^2-2\gamma^2)
=4\gamma^2(\omega_0^2-\gamma^2).\]
由$X=F_0/(m\sqrt D)$得到
\[\boxed{\Omega_r=\sqrt{\omega_0^2-2\gamma^2},\quad
 X_{\max}=\frac{F_0}{2m\gamma\sqrt{\omega_0^2-\gamma^2}}}.\]
$D'$在驻点之前为负、之后为正；又有$X\sim F_0/(m\Omega^2)\to0$，故图~\ref{fig:c2-h8-s}实线的内部峰是全局最大值。上述存在条件也保证根号均为实正。

当$\gamma\ge\omega_0/\sqrt2$时，抛物线顶点在$u\le0$，可行半轴上的最小点为$u=0$。此时$D'$在$\Omega>0$为正，振幅单调下降，最大值是静态响应$F_0/(m\omega_0^2)$，如图~\ref{fig:c2-h8-s}虚线所示。位移峰消失的临界值$\gamma=\omega_0/\sqrt2$给出$D=\omega_0^4+\Omega^4$；原点附近曲线较平，仍没有非零频率的峰。不能在根号变成虚数后继续报告一个位移“共振频率”。

弱阻尼时$\Omega_r\simeq\omega_0$，但二者通常并不严格相等。欠阻尼条件仅为$\gamma<\omega_0$；在$\omega_0/\sqrt2\le\gamma<\omega_0$内，自由运动仍有振荡，位移频率响应却没有内部峰。

\begin{figure}[H]\centering
\begin{tikzpicture}[>=Stealth,line cap=round,line join=round]
\begin{scope}[x=1.01818cm,y=1.00000cm]\draw[axis](0,0)--(2.75,0)node[right]{$\Omega/\omega_0$};\draw[axis](0,0)--(0,3)node[above]{$m\omega_0^2X/F_0$};\draw[curve,domain=0:2.5,samples=101]plot(\x,{1/sqrt((1-(\x)^2)^2+.16*(\x)^2)});\draw[dashed,domain=0:2.5,samples=101]plot(\x,{1/sqrt((1-(\x)^2)^2+2.56*(\x)^2)});\node[right]at(3,2.4){实线：$\gamma/\omega_0=.2$};\node[right]at(3,1.8){虚线：$\gamma/\omega_0=.8$};\end{scope}
\end{tikzpicture}
\caption{自由振荡存在与位移共振峰存在，阈值不同。}\label{fig:c2-h8-s}
\end{figure}
\end{solution}

\Needspace{8\baselineskip}
\begin{exercise}[吸收功率的峰值]\label{ex:c2-h9}
仍取稳态受迫阻尼振子。求驱动力每周期的平均功率，确定功率最大时的驱动频率，并与位移共振比较。 见图~\ref{fig:c2-h9-q}。
\end{exercise}

\begin{figure}[H]\centering
\begin{tikzpicture}[>=Stealth,line cap=round,line join=round]
\draw[axis](0,0)--(5,0)node[right]{$\Omega t$};\draw[axis](0,-1.3)--(0,1.5);\draw[curve,domain=0:4.7]plot(\x,{cos(\x r)});\draw[dashed,domain=0:4.7]plot(\x,{cos(\x r-.8 r)});\node at(3.8,1.4){实线$F$，虚线$v$};
\end{tikzpicture}
\caption{平均功率取决于力与速度之间的相位差。}\label{fig:c2-h9-q}
\end{figure}
\begin{solution}
图~\ref{fig:c2-h9-q}给出力与速度的相位关系。写$x=X\cos(\Omega t-\delta)$，则
\[v=\Omega X[\sin\delta\cos\Omega t-\cos\delta\sin\Omega t].\]
当 $\Omega>0$，令驱动周期 $T=2\pi/\Omega$。与 $F_0\cos\Omega t$ 相乘后，用到两个周期积分：
\[
\frac1T\int_0^T\cos^2\Omega t\dd t
=\frac1T\left[\frac t2+\frac{\sin2\Omega t}{4\Omega}\right]_0^T=\frac12,
\]
\[
\frac1T\int_0^T\sin\Omega t\cos\Omega t\dd t
=\frac1T\left[\frac{\sin^2\Omega t}{2\Omega}\right]_0^T=0.
\]
因此
\[\overline P=\frac12F_0\Omega X\sin\delta.\]
记$D(\Omega)=(\omega_0^2-\Omega^2)^2+4\gamma^2\Omega^2$，代入稳态振幅$X=F_0/(m\sqrt D)$及$\sin\delta=2\gamma\Omega/\sqrt D$，得到
\[\boxed{\overline P=m\gamma\Omega^2X^2
=\frac{F_0^2}{m}\frac{\gamma\Omega^2}{D(\Omega)}}.\]
稳态平均机械能不变，故输入功率也等于阻尼耗散$2m\gamma\overline{v^2}$。因此两种功率算法给出相同结果。

令 $u=\Omega^2$，则 $D=u^2+(4\gamma^2-2\omega_0^2)u+\omega_0^4$，$D'=2u+4\gamma^2-2\omega_0^2$。商法则的分子展开为 $D-uD'=u^2+(4\gamma^2-2\omega_0^2)u+\omega_0^4-u(2u+4\gamma^2-2\omega_0^2)=\omega_0^4-u^2$，所以
\[\frac{\dd}{\dd u}\frac uD=\frac{D-uD'}{D^2}
=\frac{\omega_0^4-u^2}{D^2}.\]
分子在$0\le u<\omega_0^2$为正，之后为负，因此图~\ref{fig:c2-h9-s}的唯一正峰位于
\[\boxed{\Omega=\omega_0,\qquad\overline P_{\max}=\frac{F_0^2}{4m\gamma}}.\]
此处$\delta=\pi/2$，速度与驱动力同相，位移落后驱动力四分之一周期。位移峰优化的是$1/D$，通常位于更低频率，甚至可能没有内部峰；功率峰优化$u/D$，对任意$\gamma>0$均在$\omega_0$。

当 $\Omega=0$，驱动力恒定，稳态速度为零，平均功率为零，与非零频率公式的连续极限一致。$\gamma\to0$ 时峰值发散，建立稳态所需的衰减时间也趋于无穷。无阻尼精确共振应按随时间增长的初值解处理，不能从此峰值式得到有限振幅稳态。

\begin{figure}[H]\centering
\begin{tikzpicture}[>=Stealth,line cap=round,line join=round]
\begin{scope}[x=1.01818cm,y=1.38462cm]\draw[axis](0,0)--(2.75,0)node[right]{$\Omega/\omega_0$};\draw[axis](0,0)--(0,1.3)node[above]{$\overline P/\overline P_{\max}$};\draw[curve,domain=0:2.5,samples=101]plot(\x,{.16*(\x)^2/((1-(\x)^2)^2+.16*(\x)^2)});\draw[guide](1,0)--(1,1);\fill(1,1)circle(1.5pt);\node[below]at(1,0){$1$};\end{scope}
\end{tikzpicture}
\caption{归一化功率在驱动频率等于无阻尼固有频率时最大。}\label{fig:c2-h9-s}
\end{figure}
\end{solution}

\Needspace{8\baselineskip}
\begin{exercise}[有限时长的驱动与冲量极限]\label{ex:c2-h10}
欠阻尼振子初始静止，外力在$0<t<\tau$为$F_0$，其余时刻为零。用响应函数写出全时间解，并取$F_0\tau=J$固定、$\tau\to0$的极限。 见图~\ref{fig:c2-h10-q}。
\end{exercise}

\begin{figure}[H]\centering
\begin{tikzpicture}[>=Stealth,line cap=round,line join=round]
\draw[axis](-.4,0)--(4.5,0)node[right]{$t$};\draw[axis](0,-.3)--(0,2)node[above]{$F$};\draw[curve](-.3,0)--(0,0)--(0,1.4)--(2.5,1.4)--(2.5,0)--(4,0);\node[above]at(1.2,1.4){$F_0$};\node[below]at(2.5,0){$\tau$};
\end{tikzpicture}
\caption{矩形脉冲有明确的开启与关闭边界。}\label{fig:c2-h10-q}
\end{figure}
\begin{solution}
图~\ref{fig:c2-h10-q}规定外力只在$0<s<\tau$输入。取因果响应核
\[G(t)=\frac{\ee^{-\gamma t}\sin\omega_dt}{\omega_d},\quad t\ge0,
\qquad G(t)=0,\quad t<0.\]
零初态的响应为
\[x(t)=\frac{F_0}{m}\int_0^{\min(t,\tau)}G(t-s)\dd s.\]
定义阶跃响应 $S(t)=\int_0^tG(u)\dd u$，并令负时间 $S=0$。为求原函数，取 $F(u)=\ee^{-\gamma u}(A\cos\omega_du+B\sin\omega_du)$，其导数为
\[
F'=\ee^{-\gamma u}[(-\gamma A+\omega_dB)\cos\omega_du
+(-\omega_dA-\gamma B)\sin\omega_du].
\]
与 $G(u)$ 比较系数，要求 $-\gamma A+\omega_dB=0$、$-\omega_dA-\gamma B=1/\omega_d$。用 $\gamma^2+\omega_d^2=\omega_0^2$ 消元，得 $A=-1/\omega_0^2$、$B=-\gamma/(\omega_d\omega_0^2)$，所以
\[
S(t)=\left[-\frac{\ee^{-\gamma u}}{\omega_0^2}
\left(\cos\omega_du+\frac\gamma{\omega_d}\sin\omega_du\right)\right]_0^t.
\]
代入两端，得到
\[S(t)=\frac{1-\ee^{-\gamma t}[\cos\omega_dt+(\gamma/\omega_d)\sin\omega_dt]}{\omega_0^2},\quad t\ge0.\]
当$t<\tau$，换元后积分为$S(t)$；当$t>\tau$，积分区间为$[t-\tau,t]$，故统一写为
\[\boxed{x(t)=\frac{F_0}{m}[S(t)-S(t-\tau)]}.\]
它也可解释为开启阶跃与延迟关闭阶跃的响应之差。$S$具有时间平方的量纲，与表示单位阶跃的无量纲函数$H$不同。

图~\ref{fig:c2-h10-s}的关闭点处，$S(0)=0$、$S'(0)=G(0)=0$保证位置和速度连续。原方程则给出加速度减少$F_0/m$。逐段求解时，也应以关闭点的连续位置、速度为自由振动的新初态。

固定 $J=F_0\tau$ 并取 $\tau\to0$。在任意固定 $t>0$ 处，Taylor 展开给出 $S(t-\tau)=S(t)-\tau S'(t)+\tau^2S''(t)/2+O(\tau^3)$，再用 $S'=G$，得到
\[S(t)-S(t-\tau)=\tau G(t)+O(\tau^2),\qquad
\boxed{x(t)\longrightarrow\frac JmG(t)}.\]
于是$x(0^+)=0$，$\dot x(0^+)=J/m$。冲量极限中力幅按$\tau^{-1}$增大，跨零积分留下有限动量变化；它允许速度跳变，与有限力撤去时只发生加速度跳变的情况不同。

\begin{figure}[H]\centering
\begin{tikzpicture}[>=Stealth,line cap=round,line join=round]
\begin{scope}[x=1.60000cm,y=1.05882cm]\draw[axis](0,0)--(1.75,0);\draw[axis](0,0)--(0,1.7)node[above]{$m\omega_0^2x/F_0$};\draw[curve,domain=0:1.5,samples=101]plot(\x,{(1-exp(-.2*\x)*(cos(.9797959*\x r)+.204124*sin(.9797959*\x r)))});\draw[curve,domain=1.5:6,samples=101]plot(\x,{exp(-.2*(\x-1.5))*(cos(.9797959*(\x-1.5) r)+.204124*sin(.9797959*(\x-1.5) r))-exp(-.2*\x)*(cos(.9797959*\x r)+.204124*sin(.9797959*\x r))});\draw[guide](1.5,0)--(1.5,1.3);\node[below]at(1.5,0){$\omega_0\tau$};\draw[axis](1.5,0)--(6.2,0)node[right]{$\omega_0t$};\end{scope}
\end{tikzpicture}
\caption{取$\gamma/\omega_0=.2,\ \omega_0\tau=1.5$；撤去外力后继续自由运动。}\label{fig:c2-h10-s}
\end{figure}
\end{solution}

\Needspace{8\baselineskip}
\begin{exercise}[变系数方程的对数坐标]\label{ex:c2-h11}
在$t>0$上求$t^2y^{\prime\prime}+at y^\prime+by=0$的通解，讨论重复根，并解释为什么不能直接把$t=0$作为普通初值点。 见图~\ref{fig:c2-h11-q}。
\end{exercise}

\begin{figure}[H]\centering
\begin{tikzpicture}[>=Stealth,line cap=round,line join=round]
\draw[axis](-.4,0)--(4,0)node[right]{$t$};\draw[axis](0,-1.5)--(0,1.8)node[above]{$u=\ln t$};\draw[curve,domain=.25:3.5,samples=101]plot(\x,{ln(\x)});\fill(1,0)circle(1.5pt)node[below]{$1$};
\end{tikzpicture}
\caption{对数坐标把伸缩变为平移，但不包含$t=0$。}\label{fig:c2-h11-q}
\end{figure}
\begin{solution}
系数中$t^2,t$的幂次提示对数坐标。令$u=\ln t$后，相邻导数带来的$t$因子恰好被这些系数抵消。图~\ref{fig:c2-h11-q}画出$t>0$到整个$u$实轴的对应，零点被推到$u=-\infty$而非一个普通有限初值点。


取$u=\ln t$，$Y(u)=y(\ee^u)$。链式法则给出
\[y'=Y'\frac{\dd u}{\dd t}=\frac{Y'}t,\qquad
y''=\frac1tY''\frac{\dd u}{\dd t}-\frac{Y'}{t^2}=\frac{Y''-Y'}{t^2},\]
故方程化为常系数形式 $Y''+(a-1)Y'+bY=0$。特征方程为 $\lambda^2+(a-1)\lambda+b=0$，根为 $\lambda_{1,2}=[1-a\pm\sqrt{(a-1)^2-4b}]/2$。不同实根时，$Y=C_1\ee^{\lambda_1u}+C_2\ee^{\lambda_2u}$；用 $\ee^u=t$，得到
\[y=C_1t^{\lambda_1}+C_2t^{\lambda_2};\]
重复实根 $\lambda$ 时，令 $Y=\ee^{\lambda u}Z$，则 $Y'=\ee^{\lambda u}(\lambda Z+Z')$、$Y''=\ee^{\lambda u}(\lambda^2Z+2\lambda Z'+Z'')$。代入方程后，$Z$ 的系数是特征多项式在重根处的值，$Z'$ 的系数是 $2\lambda+a-1$，两者均为零，所以 $Z''=0$。两次积分得 $Z=C_1+C_2u$，还原为
\[\boxed{y=t^\lambda(C_1+C_2\ln t)};\]
共轭复根$\sigma\pm\mathrm i\eta$时
\[y=t^\sigma[C_1\cos(\eta\ln t)+C_2\sin(\eta\ln t)].\]
图~\ref{fig:c2-h11-s}采用$t^2y^{\prime\prime}+ty^{\prime}+y=0$，画出$y=\cos\ln t$。振荡在$\ln t$上等间距，在$t$上则具有等比尺度。

在$t=0$处最高阶系数为零，除以$t^2$后系数发散，常规二阶初值定理不适用。若题目要求延拓至零，须另以解的有界性、可微性或物理边界条件选择分支，不能机械代入两个初值。


也可先试探 $y=t^\lambda$，求导得 $y'=\lambda t^{\lambda-1}$、$y''=\lambda(\lambda-1)t^{\lambda-2}$。代入后约去非零 $t^\lambda$，得到指标方程
\[\lambda(\lambda-1)+a\lambda+b=\lambda^2+(a-1)\lambda+b=0.\]
与对数坐标的特征方程完全相同。若有重根，第二解从$u\ee^{\lambda u}$变回$t^\lambda\ln t$；复根的振荡则发生在$\ln t$而非$t$中。

图~\ref{fig:c2-h11-s}的相邻同相位点满足$\eta\ln t$相差$2\pi$，因此$t_{n+1}/t_n=\ee^{2\pi/\eta}$。因此这些时刻在原时间轴上按等比数列排布。讨论$t\to0^+$的有界性时，需同时看幂因子$t^\sigma$与对数振荡，不能仅凭有正弦余弦就认为解有普通周期。

\begin{figure}[H]\centering
\begin{tikzpicture}[>=Stealth,line cap=round,line join=round]
\begin{scope}[x=1.00000cm,y=1.00000cm]\draw[axis](0,0)--(5.25,0)node[right]{$t$};\draw[axis](0,-1.3)--(0,1.4)node[above]{$y$};\draw[curve,domain=0.15:5,samples=101]plot(\x,{cos(ln(\x) r)});\end{scope}
\end{tikzpicture}
\caption{方程$t^2y^{\prime\prime}+ty^\prime+y=0$的一支解，定义域为$t>0$。}\label{fig:c2-h11-s}
\end{figure}
\end{solution}

\Needspace{8\baselineskip}
\begin{exercise}[由衰减峰值确定阻尼]\label{ex:c2-h12}
实测欠阻尼振子的相邻同侧位移峰值为$A_n,A_{n+1}$，间隔为$T_d$。由这些量求$\gamma,\omega_0$；若记录的是相邻异侧峰值的绝对值，公式如何改变？ 见图~\ref{fig:c2-h12-q}。
\end{exercise}

\begin{figure}[H]\centering
\begin{tikzpicture}[>=Stealth,line cap=round,line join=round]
\begin{scope}[x=1.00000cm,y=1.00000cm]\draw[axis](0,0)--(7.25,0)node[right]{$\omega_dt$};\draw[axis](0,-1.1)--(0,1.2)node[above]{$x/A_0$};\draw[curve,domain=0:7,samples=101]plot(\x,{exp(-.15*\x)*(cos(\x r)+.15*sin(\x r))});\end{scope}
\end{tikzpicture}
\caption{衰减包络限制峰值，实测峰间隔须注明同侧或异侧。}\label{fig:c2-h12-q}
\end{figure}
\begin{solution}
阻尼振子的真实峰点略偏离简单余弦的峰点，但每隔一个阻尼周期，峰点相位会重复。因此同侧相邻峰的比值仍只含指数包络。图~\ref{fig:c2-h12-q}的横轴使用阻尼角频率，先从测得周期求$\omega_d$，再恢复$\omega_0$。


自由运动写成 $x=C\ee^{-\gamma t}\cos\psi$，其中 $\psi=\omega_dt-\varphi$。对乘积求导，$\dot x=-C\ee^{-\gamma t}(\gamma\cos\psi+\omega_d\sin\psi)$，所以峰值条件为 $\tan\psi=-\gamma/\omega_d$。同侧相邻峰的相位差为 $2\pi$，余弦因子相同；两个峰值相除，$A_n/A_{n+1}=\ee^{-\gamma t_n}/\ee^{-\gamma(t_n+T_d)}=\ee^{\gamma T_d}$，取对数得
\[\Lambda=\ln\frac{A_n}{A_{n+1}}=\gamma T_d,\qquad
\omega_d=\frac{2\pi}{T_d}.\]
由图~\ref{fig:c2-h12-s}标出的两个同侧峰，得到
\[\boxed{\gamma=\Lambda/T_d,\qquad
\omega_0=\sqrt{(2\pi/T_d)^2+(\Lambda/T_d)^2}}.\]
若记录相邻异侧峰，相位差为$\pi$，余弦变号而绝对值相同。设时间间隔为$\Delta t$、绝对峰值对数比为$\Lambda_{1/2}$，仍有$|A_n|/|A_{n+1}|=\ee^{\gamma\Delta t}$，故$\gamma=\Lambda_{1/2}/\Delta t$，但$\omega_d=\pi/\Delta t$。阻尼率的比值形式不变，频率关系出现因子二。

实际数据可对多个同侧峰拟合 $\ln|A_n|$ 对 $t_n$ 的直线，其斜率为 $-\gamma$；取绝对值使负侧峰也能使用。峰值较小时若背景噪声占比增大，应说明使用的数据区间；这个判断来自数据质量，不改变上述解析关系。


图~\ref{fig:c2-h12-s}标出的同侧峰间隔为$T_d=2\pi/\omega_d$。确定$\gamma$与$\omega_d$后，再用$\omega_0^2=\omega_d^2+\gamma^2$，这是从根的实部、虚部恢复原方程系数。若把$\omega_d$直接叫作无阻尼固有频率，就会系统低估$\omega_0$。

\begin{figure}[H]\centering
\begin{tikzpicture}[>=Stealth,line cap=round,line join=round]
\begin{scope}[x=1.00000cm,y=1.00000cm]\draw[axis](0,0)--(7.25,0)node[right]{$\omega_dt$};\draw[axis](0,-1.1)--(0,1.2)node[above]{$x/A_0$};\draw[curve,domain=0:7,samples=101]plot(\x,{exp(-.15*\x)*(cos(\x r)+.15*sin(\x r))});\draw[dashed,domain=0:7]plot(\x,{exp(-.15*\x)});\draw[<->](0,1.45)--(6.283185,1.45)node[midway,above]{$\omega_dT_d=2\pi$};\fill(0,1)circle(1.5pt);\fill(6.283185,.389661)circle(1.5pt);\end{scope}
\end{tikzpicture}
\caption{同侧峰间隔为$T_d$；虚线经过各峰，其高度按指数衰减。}\label{fig:c2-h12-s}
\end{figure}
\end{solution}

\Needspace{8\baselineskip}
\begin{exercise}[临界阻尼是否穿越平衡点]\label{ex:c2-h13}
临界阻尼振子初态为$x(0)=a>0,\dot x(0)=-u$，$u\ge0$。求首次到达$x=0$的条件与时间，讨论临界边界。 见图~\ref{fig:c2-h13-q}。
\end{exercise}

\begin{figure}[H]\centering
\begin{tikzpicture}[>=Stealth,line cap=round,line join=round]
\draw[curve](0,0)--(.2,0)--(.4,.2)--(.6,-.2)--(.8,.2)--(1,-.2)--(1.2,.2)--(1.4,-.2)--(1.6,.2)--(1.8,-.2)--(2,0)--(2.2,0);
\draw[thick](0,-.35)--(0,.35);\draw(2.2,-.3)rectangle(2.8,.3);
\draw[axis](2.4,.65)--(1.3,.65)node[midway,above]{$u$};
\draw[axis](.7,-.8)--(4,-.8)node[right]{$x$};
\draw(1,-.74)--(1,-.86)node[below]{$0$};
\draw[guide](2.5,-.3)--(2.5,-.8);\draw(2.5,-.74)--(2.5,-.86)node[below]{$a$};
\end{tikzpicture}
\caption{位移原点取在平衡位置，初态$x=a>0$，初速度向左；弹簧固定端与位移原点分开标示。}\label{fig:c2-h13-q}
\end{figure}
\begin{solution}
图~\ref{fig:c2-h13-q}中，质点从平衡点右侧出发，初速度向左。临界阻尼方程的重复根为$-\omega_0$，通解为$x=(C_1+C_2t)\ee^{-\omega_0t}$。由$x(0)=a$及$\dot x(0)=C_2-\omega_0C_1=-u$，得$C_1=a$、$C_2=\omega_0a-u$，所以
\[x(t)=\ee^{-\omega_0t}[a+(\omega_0a-u)t].\]
指数因子在有限时间内严格为正。图~\ref{fig:c2-h13-s}中的穿越，只可能由方括号变号引起。因此
\[\boxed{u>\omega_0a\quad\Longleftrightarrow\quad
 t_0=\frac a{u-\omega_0a}<\infty}.\]
在$u>\omega_0a$的情形，对位移求导，并利用$t=t_0$时方括号为零，得到穿越速度
\[\dot x(t_0)=(\omega_0a-u)\ee^{-\omega_0t_0}<0.\]
质点以非零速度穿过原点。方括号是一次函数，只有一个零点，因而正时间内最多穿越一次。

当$u=\omega_0a$时，$x=a\ee^{-\omega_0t}>0$，对应图~\ref{fig:c2-h13-s}中的虚线，只在$t\to\infty$时趋于平衡点。当$0\le u<\omega_0a$时，方括号从$a$开始增加，也始终为正。由穿越一侧令$u\downarrow\omega_0a$，可见$t_0\to\infty$；这给出了有限时间穿越与渐近趋近之间的临界边界。临界阻尼的重根本身并不排除一次穿越，还须结合初值判断。

\begin{figure}[H]\centering
\begin{tikzpicture}[>=Stealth,line cap=round,line join=round]
\begin{scope}[x=1.15cm,y=2.2cm]
\draw[axis](0,0)--(5.25,0)node[right]{$\omega_0t$};\draw[axis](0,-.22)--(0,1.15)node[above]{$x/a$};
\draw[curve,domain=0:5,samples=121]plot(\x,{(1-.6*\x)*exp(-\x)});
\draw[dashed,domain=0:5,samples=121]plot(\x,{exp(-\x)});
\draw[dotted,domain=0:5,samples=121]plot(\x,{(1+.5*\x)*exp(-\x)});
\fill({5/3},0)circle(1.6pt)node[below]{$5/3$};
\node[align=left,right]at(2.8,.85){$u/(\omega_0a)$：\\实线 $1.6$\\虚线 $1$，点线 $0.5$};\end{scope}
\end{tikzpicture}
\caption{实线在$\omega_0t=5/3$穿越原点；另外两条始终为正。扩大负位移区使一次穿越可辨。}\label{fig:c2-h13-s}
\end{figure}
\end{solution}

\Needspace{8\baselineskip}
\begin{exercise}[带饱和项的一阶增长]\label{ex:c2-h14}
状态量$y>0$满足$\dot y=ay(1-y/K)$，$a,K>0$，$y(0)=y_0>0$。求解并讨论$y_0<K$与$y_0>K$两侧的趋近过程。 见图~\ref{fig:c2-h14-q}。
\end{exercise}

\begin{figure}[H]\centering
\begin{tikzpicture}[>=Stealth,line cap=round,line join=round]
\begin{scope}[x=1.60000cm,y=1.38462cm]\draw[axis](0,0)--(1.75,0)node[right]{$y/K$};\draw[axis](0,-0.9)--(0,0.4)node[above]{$\dot y/(aK)$};\draw[curve,domain=0:1.5,samples=101]plot(\x,{\x*(1-\x)});\end{scope}
\end{tikzpicture}
\caption{正根$y=K$两侧的变化率异号。}\label{fig:c2-h14-q}
\end{figure}
\begin{solution}
两个平衡值为$0$与$K$，正值状态在它们之间增长，在$K$上方下降。图~\ref{fig:c2-h14-q}的抛物线给出这三个符号区间。分离积分时先保留绝对值，才可在$y_0<K$和$y_0>K$两侧使用同一最终表达式。


若 $y_0=K$，直接有平衡解 $y=K$。以下先取 $y_0\ne K$，分离变量并分解为 $1/[y(1-y/K)]=1/y+1/(K-y)$。从初态积到当前状态，
\[
\int_{y_0}^{y(t)}\left(\frac1u+\frac1{K-u}\right)\dd u
=[\ln|u|-\ln|K-u|]_{y_0}^{y(t)}=at.
\]
因此 $|y/(K-y)|=|y_0/(K-y_0)|\ee^{at}$。解在正时间不穿越平衡值 $K$，比值符号由初态确定，故可写 $y/(K-y)=[y_0/(K-y_0)]\ee^{at}$。交叉相乘并把含 $y$ 的项合并，得 $y[K-y_0+y_0\ee^{at}]=Ky_0\ee^{at}$，所以
\[\boxed{y(t)=\frac K{1+(K/y_0-1)\ee^{-at}}\qquad(t\ge0)}.
\]
当$y_0<K$，分母从$K/y_0>1$下降到$1$，状态单调上升；$y_0>K$时分母从小于$1$的正数上升到$1$，状态单调下降。图~\ref{fig:c2-h14-s}中两支都趋于$K$而不在有限时间穿越。平衡解$y=0$和$y=K$须另行列出，其中$K$稳定、$0$对正扰动不稳定。

令 $y=K+\eta$，原方程右侧为 $a(K+\eta)(-\eta/K)=-a\eta-a\eta^2/K$，所以线性项为 $\dot\eta=-a\eta$，长时偏差按 $\ee^{-at}$ 衰减，与精确解展开一致。若分离时只取对数而丢掉绝对值，容易误以为只适用于 $y<K$；最终有理式在两侧都成立。


对于 $y_0>K$，解的分母满足
\[1+(K/y_0-1)\ee^{-at}\ge K/y_0>0\qquad(t\ge0),\]
所以本题正时间内没有极点。长时有$y-K\sim K(1-K/y_0)\ee^{-at}$，系数正负分别对应图中从上、下两侧趋近的曲线。

在$0<y<K$内，增长率$\dot y=ay(1-y/K)$在$y=K/2$最大。若初态低于$K/2$，时间曲线先加速增长再减速；若初态已高于$K/2$，则从一开始减速。图~\ref{fig:c2-h14-s}中的趋近过程具有同一稳定终值，初期曲率则由初态所在区间决定。

\begin{figure}[H]\centering
\begin{tikzpicture}[>=Stealth,line cap=round,line join=round]
\begin{scope}[x=1.00000cm,y=1.00000cm]\draw[axis](0,0)--(5.25,0)node[right]{$at$};\draw[axis](0,0)--(0,2)node[above]{$y/K$};\draw[curve,domain=0:5,samples=101]plot(\x,{1/(1+2*exp(-\x))});\draw[dashed,domain=0:5]plot(\x,{1/(1-.4*exp(-\x))});\draw[guide](0,1)--(5,1);\end{scope}
\end{tikzpicture}
\caption{分别从下方和上方趋近同一稳定平衡。}\label{fig:c2-h14-s}
\end{figure}
\end{solution}

\Needspace{8\baselineskip}
\begin{exercise}[线性化失效时的幂律弛豫]\label{ex:c2-h15}
求$\dot y=-ky^3$在$k>0,y(0)=y_0\ne0$下的解，判断零平衡的稳定性，并与$\dot y=-\lambda y$的衰减比较。 见图~\ref{fig:c2-h15-q}。
\end{exercise}

\begin{figure}[H]\centering
\begin{tikzpicture}[>=Stealth,line cap=round,line join=round]
\begin{scope}[x=1.05660cm,y=1.00000cm]\draw[axis](-1.2,0)--(1.45,0)node[right]{$y$};\draw[axis](0,-2)--(0,2)node[above]{$\dot y/k$};\draw[curve,domain=-1.2:1.2,samples=101]plot(\x,{-(\x)^3});\end{scope}
\end{tikzpicture}
\caption{两侧箭头都指向零点，但在零点处线性项消失。}\label{fig:c2-h15-q}
\end{figure}
\begin{solution}
线性化在原点得到零系数，不能据此判断稳定性。图~\ref{fig:c2-h15-q}中的三次变化率对正$y$为负、对负$y$为正，已经表明两侧都向零运动。分离积分后还要用初始符号选择平方根分支。


在不为零的分支分离，
\[
\int_{y_0}^{y(t)}u^{-3}\dd u=-k\int_0^t\dd s,\qquad
\left[-\frac1{2u^2}\right]_{y_0}^{y(t)}=-kt.
\]
所以 $1/y^2=1/y_0^2+2kt$，结合初始符号取平方根，得到
\[\boxed{y(t)=\frac{y_0}{\sqrt{1+2ky_0^2t}}}.
\]
正初态递减，负初态递增，二者都趋于零，故零点渐近稳定。图~\ref{fig:c2-h15-s}中长时$|y|\sim(2kt)^{-1/2}$；它慢于任意固定衰减率的指数。

$f'(0)=0$使一阶线性化成为$\dot\eta=0$，只能说明线性项没有提供信息，不能据此宣称原系统中性稳定。三次项的符号才决定本题的稳定性。恒零解在分离变量时被除去，必须补回；各非零解只在无穷远接近零，不发生有限时间并合。


若希望幅值降至初值的$\epsilon$倍，所需时间为
\[t_\epsilon=\frac{\epsilon^{-2}-1}{2ky_0^2}.\]
线性指数衰减所需时间则为$\ln(1/\epsilon)/\lambda$。图~\ref{fig:c2-h15-s}比较的长时尾部因此明显不同；三次系统接近零后恢复越来越弱。稳定性可由$V=y^2/2$说明：$\dot V=-ky^4<0$对非零$y$成立，并结合显式解得到渐近趋零。

\begin{figure}[H]\centering
\begin{tikzpicture}[>=Stealth,line cap=round,line join=round]
\begin{scope}[x=1.00000cm,y=1.38462cm]\draw[axis](0,0)--(5.25,0)node[right]{$\tau$};\draw[axis](0,0)--(0,1.3)node[above]{$y/y_0$};\draw[curve,domain=0:5,samples=101]plot(\x,{1/sqrt(1+2*\x)});\draw[dashed,domain=0:5]plot(\x,{exp(-\x)});\end{scope}
\end{tikzpicture}
\caption{取$\tau=ky_0^2t$；实线为幂律衰减，虚线为指数衰减。}\label{fig:c2-h15-s}
\end{figure}
\end{solution}

\Needspace{8\baselineskip}
\begin{exercise}[速度作为位置函数的降阶]\label{ex:c2-h16}
质点满足$\ddot x=-\omega_0^2x$，$x(0)=0,\dot x(0)=v_0>0$。用$v\dd v/\dd x$求相轨道，说明为何一条能量曲线还不足以确定带方向的时间解。 见图~\ref{fig:c2-h16-q}。
\end{exercise}

\begin{figure}[H]\centering
\begin{tikzpicture}[>=Stealth,line cap=round,line join=round]
\draw[axis](-1.8,0)--(1.8,0)node[right]{$x$};\draw[axis](0,-1.7)--(0,1.8)node[above]{$v$};\draw[curve](0,0)ellipse(1.4 and 1.4);\fill(0,1.4)circle(1.5pt)node[right]{初态};
\end{tikzpicture}
\caption{相轨道是能量椭圆，运动方向需要由方程补充。}\label{fig:c2-h16-q}
\end{figure}
\begin{solution}
不显含时间的简谐方程可用$v(x)$降阶，但$x$在一周内并不单调，所以$v(x)$只能分段取单值。图~\ref{fig:c2-h16-q}的相轨道把向右、向左运动同时画在同一椭圆上，初态给出应从哪一点沿哪个方向出发。


在速度可看作位置函数的单调段，$v\dd v=-\omega_0^2x\dd x$。从初态积分，$[v^2/2]_{v_0}^{v}=-\omega_0^2[x^2/2]_0^x$，即 $v^2+\omega_0^2x^2=v_0^2$。初始向右，第一段取正根。为计算时间，令 $\xi=(v_0/\omega_0)\sin u$、$\dd\xi=(v_0/\omega_0)\cos u\dd u$，并取 $0\le u\le\pi/2$。根式为 $v_0\cos u$，约去共同因子后，
\[t=\int_0^x\frac{\dd\xi}{\sqrt{v_0^2-\omega_0^2\xi^2}}
=\frac1{\omega_0}\int_0^{\arcsin(\omega_0x/v_0)}\dd u
=\omega_0^{-1}\arcsin(\omega_0x/v_0)
\]
恢复$x=(v_0/\omega_0)\sin\omega_0t$，直到第一次转向。图~\ref{fig:c2-h16-s}在右端转向后速度改为负根，继续沿椭圆顺时针运动，完整解仍由正弦式给出。

平方后的相轨道包含两个速度分支。初始方向、经过转向点后的连接规则以及时间原点，都是恢复时间解所需的信息。只写$\dot x^2=v_0^2-\omega_0^2x^2$而不说明分支，尚未完成初值问题。


第一段速度为正，积分出的反正弦取$[-\pi/2,\pi/2]$主支，故只直接描述$0\le t\le\pi/(2\omega_0)$。到达$x=v_0/\omega_0$时，速度为零而加速度为$-\omega_0v_0<0$，随后应接负根，返回平衡点。

图~\ref{fig:c2-h16-s}中相空间速度为$(\dot x,\dot v)=(v,-\omega_0^2x)$。在椭圆顶部它向右，在右端点它向下，所以轨道顺时针遍历。能量椭圆给出允许运动集合，向量场给出方向和时间参数；两者一起才恢复完整的初值解，而不是在每个转向点重新任取正根。

\begin{figure}[H]\centering
\begin{tikzpicture}[>=Stealth,line cap=round,line join=round]
\draw[axis](-1.8,0)--(1.8,0)node[right]{$\omega_0x/v_0$};\draw[axis](0,-1.7)--(0,1.8)node[above]{$v/v_0$};\draw[curve](0,0)circle(1.2);\draw[axis](-.2,1.2)--(.2,1.2);\draw[axis](1.2,.2)--(1.2,-.2);\draw[axis](.2,-1.2)--(-.2,-1.2);\node[right]at(1.2,-.3){转向};
\end{tikzpicture}
\caption{归一化相图为圆，箭头明确选择了实际运动的方向。}\label{fig:c2-h16-s}
\end{figure}
\end{solution}

\Needspace{8\baselineskip}
\begin{exercise}[拍与共振的连续极限]\label{ex:c2-h17}
无阻尼振子初始静止，受单位质量外力$f_0\cos\Omega t$，$\Omega\ne\omega_0$。求解并在$|\Omega-\omega_0|\ll\omega_0$时写出拍包络，说明何时可用精确共振的线性包络近似。 见图~\ref{fig:c2-h17-q}。
\end{exercise}

\begin{figure}[H]\centering
\begin{tikzpicture}[>=Stealth,line cap=round,line join=round]
\draw[axis](0,0)--(6,0)node[right]{$t$};\draw[axis](0,-1.4)--(0,1.5)node[above]{$f$};\draw[curve,domain=0:5.8,samples=151]plot(\x,{cos(4*\x r)});\draw[dashed,domain=0:5.8,samples=151]plot(\x,{cos(4.4*\x r)});
\end{tikzpicture}
\caption{驱动与固有频率接近，二者相位差缓慢积累。}\label{fig:c2-h17-q}
\end{figure}
\begin{solution}
为满足零初态，须在周期特解上减去自由余弦，响应中便出现两个相近频率。图~\ref{fig:c2-h17-q}中外源与固有振荡的相位差缓慢积累，形成慢包络；本题没有阻尼耗散。


先取周期特解 $x_p=C\cos\Omega t$，代入方程得 $(\omega_0^2-\Omega^2)C=f_0$，故 $C=f_0/(\omega_0^2-\Omega^2)$。加上齐次通解，$x=x_p+C_1\cos\omega_0t+C_2\sin\omega_0t$；初位移为零给出 $C_1=-C$，初速度为零给出 $\omega_0C_2=0$。因此
\[x=\frac{f_0}{\omega_0^2-\Omega^2}(\cos\Omega t-\cos\omega_0t).\]
令$\Delta=\Omega-\omega_0$，和差化积给出
\[
 \cos\Omega t-\cos\omega_0t
 =-2\sin\frac{(\Omega+\omega_0)t}{2}\sin\frac{\Delta t}{2}.
\]
分母为$-\Delta(\Omega+\omega_0)$，两个负号相消，得到
\[
 x=\frac{2f_0}{\Delta(\Omega+\omega_0)}
 \sin\frac{(\Omega+\omega_0)t}{2}\sin\frac{\Delta t}{2}.
\]
图~\ref{fig:c2-h17-s}中的快振荡频率接近$\omega_0$，幅度包络为$2|f_0\sin(\Delta t/2)|/[|\Delta|(\Omega+\omega_0)]$，其中$\Delta=\Omega-\omega_0$。

当$|\Delta|t\ll1$，慢正弦可近似$\Delta t/2$，于是$x\simeq f_0t\sin\omega_0t/(2\omega_0)$。只有在这个时间窗口内，近共振响应才近似精确共振的线性增长。长到$t\sim1/|\Delta|$，包络会弯曲并重新归零；不能把固定时间的频率极限当成对所有时间一致成立的近似。


包络第一次重新为零的正时刻为$2\pi/|\Delta|$，与振荡快周期约$2\pi/\omega_0$分离。

图~\ref{fig:c2-h17-s}的包络在初期近似直线，但完整包络有界，最大幅值为$2|f_0|/[|\Delta|(\Omega+\omega_0)]$。固定$t$再令$\Delta\to0$得到共振增长；若$t$也随$1/|\Delta|$增长，就必须保留慢正弦，不能作同一展开。

\begin{figure}[H]\centering
\begin{tikzpicture}[>=Stealth,line cap=round,line join=round]
\begin{scope}[x=0.65000cm,y=0.66667cm]\draw[axis](0,0)--(9.25,0)node[right]{$t$};\draw[axis](0,-1.3)--(0,1.4)node[above]{$x/X_*$};\draw[curve,domain=0:9,samples=101]plot(\x,{sin(4*\x r)*sin(.35*\x r)});\draw[dashed,domain=0:9,samples=101]plot(\x,{sin(.35*\x r)});\draw[dashed,domain=0:9,samples=101]plot(\x,{-sin(.35*\x r)});\end{scope}
\end{tikzpicture}
\caption{快振荡被慢包络调制，包络的初始小段近似直线。}\label{fig:c2-h17-s}
\end{figure}
\end{solution}

\Needspace{8\baselineskip}
\begin{exercise}[线性化的摆与受驱平衡]\label{ex:c2-h18}
单摆满足$\ddot\theta+2\gamma\dot\theta+(g/l)\sin\theta=\tau/(ml^2)$，$\tau$为恒定外力矩，$\gamma\ge0$。求静平衡及其附近小扰动频率，讨论静平衡不存在的条件。 见图~\ref{fig:c2-h18-q}。
\end{exercise}

\begin{figure}[H]\centering
\begin{tikzpicture}[>=Stealth,line cap=round,line join=round]
\fill(0,0)circle(2pt)node[above]{$O$};\draw[guide](0,0)--(0,-2.8);\draw[curve](0,0)--(1.4,-2.2);\fill(1.4,-2.2)circle(3pt)node[right]{$m$};\draw[axis](.5,.1)arc(10:-70:.5);\node at(.9,-.3){$\tau$};
\end{tikzpicture}
\caption{恒力矩使平衡位置偏离竖直向下方向。}\label{fig:c2-h18-q}
\end{figure}
\begin{solution}
恒定外力矩改变平衡角，微小振动应围绕新平衡展开。图~\ref{fig:c2-h18-q}中的力矩方向固定，因此平衡条件是重力矩与外力矩相等；直接以$\theta=0$为中心作小角近似，会漏掉偏置平衡与失稳阈值。


平衡要求 $\sin\theta_*=\tau/(mgl)=s$。若 $|s|>1$，不存在静平衡。若 $|s|<1$，每周期有两支：$\cos\theta_*>0$ 的下方支和 $\cos\theta_*<0$ 的上方支。令 $\theta=\theta_*+\eta$，展开 $\sin(\theta_*+\eta)=\sin\theta_*+\eta\cos\theta_*+O(\eta^2)$；常数项由平衡方程抵消，保留一阶得
\[\ddot\eta+2\gamma\dot\eta+\omega_*^2\eta=0,
\qquad \boxed{\omega_*^2=(g/l)\cos\theta_*}.\]
图~\ref{fig:c2-h18-s}中下方支有正刚度，无阻尼频率$\omega_*$；若欠阻尼，实际衰减振荡频率为$\sqrt{\omega_*^2-\gamma^2}$。上方支刚度为负，总有一个增长模态，阻尼不能把它变成局部势阱。

在$|s|=1$处刚度为零，线性方程退化，须保留非线性项讨论临界失稳。直接把$\sin\theta$替成$\theta$，只适用于原点附近，不能用于任意偏置平衡。


当$|s|<1$，可把两支明确写成
\[\theta_s=\arcsin s+2j\pi,\qquad
\theta_u=\pi-\arcsin s+2j\pi.\]
前者有$\cos\theta_s=\sqrt{1-s^2}>0$。若无阻尼，$\omega_* =\sqrt{g/l}\,(1-s^2)^{1/4}$；它随$|s|\to1$趋零。后者的特征方程为$\lambda^2+2\gamma\lambda-|\omega_*^2|=0$，根$-\gamma+\sqrt{\gamma^2+|\omega_*^2|}>0$，故确有增长模态。

在$s=1,\theta_*=\pi/2$处，$\sin(\theta_*+\eta)=1-\eta^2/2+O(\eta^4)$，因此
\[\ddot\eta+2\gamma\dot\eta=\frac g{2l}\eta^2+O(\eta^4).\]
任意小的正向静止扰动都会被推离，故该临界平衡也不稳定；$s=-1$时方向相反。解答图中水平线与正弦曲线相切，表示两支平衡在此合并，而不是形成有限频率的小振动。

一阶刚度 $(g/l)\cos\theta_*$ 的正负，直接对应图~\ref{fig:c2-h18-s}中正弦交点处的斜率。

下方支在$\gamma>0$时渐近稳定；$\gamma=0$时，小扰动振荡而不衰减。

\begin{figure}[H]\centering
\begin{tikzpicture}[>=Stealth,line cap=round,line join=round]
\begin{scope}[x=0.75000cm,y=0.75000cm]\draw[axis](-1.6,0)--(5.05,0)node[right]{$\theta$};\draw[axis](0,-1.3)--(0,1.4)node[above]{$\sin\theta$};\draw[curve,domain=-1.6:4.8,samples=101]plot(\x,{sin(\x r)});\draw[guide](-1.6,.6)--(4.8,.6);\fill(.643501,.6)circle(1.5pt)node[above left]{稳定支};\fill(2.498092,.6)circle(1.5pt)node[above right]{不稳定支};\end{scope}
\end{tikzpicture}
\caption{同一恒力矩可对应稳定与不稳定两支平衡，斜率符号决定局部刚度。}\label{fig:c2-h18-s}
\end{figure}
\end{solution}

\Needspace{8\baselineskip}
\begin{exercise}[由已知解构造第二解]\label{ex:c2-h19}
在$t>0$上求$t^2y^{\prime\prime}-2ty^{\prime}+2y=0$的两个独立解。先验证$y_1=t$，再用降阶公式求$y_2$，以Wronskian确认独立性。 见图~\ref{fig:c2-h19-q}。
\end{exercise}

\begin{figure}[H]\centering
\begin{tikzpicture}[>=Stealth,line cap=round,line join=round]
\draw[axis](0,0)--(3,0)node[right]{$t$};\draw[axis](0,0)--(0,2.5)node[above]{$y$};\draw[curve](0,0)--(2,2);\node[right]at(1.8,1.6){$y_1=t$};
\end{tikzpicture}
\caption{已知一个解后，第二解仍需有不同的时间依赖。}\label{fig:c2-h19-q}
\end{figure}
\begin{solution}
把图~\ref{fig:c2-h19-q}中的$y_1=t$代入原式，得到$0-2t+2t=0$，故它确为一个解。在$t>0$上除以$t^2$，标准形式为
\[y''-\frac2t y'+\frac2{t^2}y=0.\]
其中 $p(t)=-2/t$，原函数为 $\int p(t)\dd t=-2\ln t+C$，所以 $\exp(-\int p\dd t)$ 是 $t^2$ 的非零常数倍，该常数可并入第二解。因为 $y_1$ 在此区间不为零，降阶公式给出
\[y_2=t\int1\,\dd t=t(t+C)=t^2+Ct.
\]
积分常数只产生已有解$t$的倍数，不必重复计入第二个独立解。因此
\[\boxed{y=C_1t+C_2t^2}.
\]
图~\ref{fig:c2-h19-s}中，$t^2/t=t$随自变量变化，两曲线不是恒定倍数。相应的Wronskian为
\[W(t)=\begin{vmatrix}t&t^2\\1&2t\end{vmatrix}=t^2>0.\]
这确认了线性独立；$W'=2t=(-p)W$也符合Abel关系。

降阶公式还可由直接代换得到。设$y=tu(t)$，则$y'=u+tu'$、$y''=2u'+tu''$，代入原方程为
\[t^2(2u'+tu'')-2t(u+tu')+2tu=t^3u''=0.\]
在$t>0$上有$u''=0$。两次积分得$u'=C_2$、$u=C_1+C_2t$，恢复同一通解；两个积分常数分别对应$t$和$t^2$。

任给$t_0>0$处的初值$y(t_0)=y_0$、$y'(t_0)=v_0$，系数满足
\[\begin{pmatrix}t_0&t_0^2\\1&2t_0\end{pmatrix}
\begin{pmatrix}C_1\\C_2\end{pmatrix}=\begin{pmatrix}y_0\\v_0\end{pmatrix}.\]
第二行减去第一行的 $1/t_0$，得到 $t_0C_2=v_0-y_0/t_0$，所以 $C_2=(t_0v_0-y_0)/t_0^2$，再由第一行得到 $C_1=(2y_0-t_0v_0)/t_0$。行列式为 $t_0^2\ne0$，两常数唯一确定。$t=0$ 是原方程的奇点，虽然多项式解可以延拓到零点，常规初值存在唯一性定理的条件在此并不成立。

\begin{figure}[H]\centering
\begin{tikzpicture}[>=Stealth,line cap=round,line join=round]
\begin{scope}[x=1.24444cm,y=1.00000cm]\draw[axis](0,0)--(2.25,0)node[right]{$t$};\draw[axis](0,0)--(0,4.5)node[above]{$y$};\draw[curve,domain=0:2,samples=101]plot(\x,{(\x)^2});\draw[dashed](0,0)--(2,2);\node[right]at(2,3.6){实线$t^2$};\node[right]at(2,1.8){虚线$t$};\end{scope}
\end{tikzpicture}
\caption{两个独立解分别按一次和二次幂增长。}\label{fig:c2-h19-s}
\end{figure}
\end{solution}

\Needspace{8\baselineskip}
\begin{exercise}[单位冲量的连接条件]\label{ex:c2-h20}
设因果函数$G$满足$G^{\prime\prime}+2\gamma G^{\prime}+\omega_0^2G=\delta(t)$，负时间为零。推导$G(0^+)$与$G^{\prime}(0^+)$，解释为什么输入冲量改变速度而不改变位置。 见图~\ref{fig:c2-h20-q}。
\end{exercise}

\begin{figure}[H]\centering
\begin{tikzpicture}[>=Stealth,line cap=round,line join=round]
\draw[axis](-1.5,0)--(2.5,0)node[right]{$t$};\draw[axis](0,0)--(0,2)node[above]{$\delta(t)$};\draw[axis](0,0)--(0,1.5);\node[right]at(.2,1.4){单位面积冲量};
\end{tikzpicture}
\caption{冲量的高度不能按普通函数解释，应使用跨越零点的积分。}\label{fig:c2-h20-q}
\end{figure}
\begin{solution}
图~\ref{fig:c2-h20-q}中的$\delta(t)$表示单位面积的冲量。写$G=Hg$，其中$g$在正时间侧光滑，则分布导数为
\[
 G'=Hg'+g(0^+)\delta,\qquad
 G''=Hg''+g'(0^+)\delta+g(0^+)\delta'.
\]
将它们代回$G''+2\gamma G'+\omega_0^2G=\delta$，奇异项分别为
\[
 g(0^+)\delta'
 +[g'(0^+)+2\gamma g(0^+)]\delta.
\]
右端没有$\delta'$，故先得$g(0^+)=0$；再比较$\delta$的系数，得$g'(0^+)=1$。由于负时间为零，
\[
 \boxed{G(0^+)=G(0^-)=0,\qquad
 G'(0^+)-G'(0^-)=1}.
\]
\begin{figure}[H]\centering
\begin{tikzpicture}[>=Stealth,line cap=round,line join=round]
\begin{scope}[x=1.00000cm,y=1cm]\draw[axis](0,0)--(6.25,0)node[right]{$t$};\draw[axis](0,-0.8)--(0,1.2)node[above]{$G$};\draw[curve,domain=0:6,samples=101]plot(\x,{exp(-.2*\x)*sin(.9797959*\x r)/.9797959});\draw[curve](-1,0)--(0,0);\draw[dashed](0,0)--(.7,.7);\node[right]at(.8,1.1){$G^\prime(0^+)=1$};\end{scope}
\end{tikzpicture}
\caption{位置连续，正侧初斜率为单位值，负侧初斜率为零。}\label{fig:c2-h20-s}
\end{figure}
也可在位置连续已经确定后，跨越零点积分原方程：
\[
 [G']_{0^-}^{0^+}+2\gamma[G]_{0^-}^{0^+}
 +\lim_{\epsilon\to0}\omega_0^2
 \int_{-\epsilon}^{\epsilon}G(t)\dd t=1.
\]
后两项为零，留下同一速度连接条件。这个积分法确定速度跃变，前面的$\delta'$比较则说明位置为什么必须连续。
正时间上解齐次方程。欠阻尼时沿用其通解 $G=\ee^{-\gamma t}(A\cos\omega_dt+B\sin\omega_dt)$，零时刻条件给出 $A=0$、$-\gamma A+\omega_dB=1$，所以 $B=1/\omega_d$，得到 $G(t)=\ee^{-\gamma t}\sin(\omega_dt)/\omega_d$。

图~\ref{fig:c2-h20-s}中的曲线因此从原点以单位斜率起步。质量$m$受到冲量$J$时，物理位移为$x=(J/m)G$，故$m[\dot x]_{0^-}^{0^+}=J$。弹簧力和有限阻尼力在趋零时间内的冲量为零；普通力冲量改变动量，不能造成位置瞬时跳变。若要求位置跳变，方程中还需更奇异的$\delta'$型外源。
\end{solution}


\begin{exercise}[有阻力的弹簧下落]\label{ex:c2-v2}

质量$m$的小球挂在竖直轻弹簧下端，弹簧劲度系数$k$。从弹簧原长位置静止释放，以向下伸长量$x$为坐标。空气阻力大小为$\beta v^2$，方向总与速度相反，$\beta>0$。求首次下落阶段$v$与$x$的关系，确定最低点的方程，并讨论是否能把同一方程继续用于回升。 见图~\ref{fig:c2-v2-q}。
\end{exercise}

\begin{figure}[H]\centering
\begin{tikzpicture}[>=Stealth,line cap=round,line join=round]

\draw[thick](-.6,2.6)--(.6,2.6);\draw[curve](0,2.6)--(.2,2.35)--(-.2,2.1)--(.2,1.85)--(-.2,1.6)--(0,1.35);
\fill(0,1.1)circle(.16);\draw[axis](.65,1.1)--(.65,.2)node[below]{$x,\ v>0$};
\draw[axis](0,.9)--(0,.15)node[left]{$mg$};\draw[axis](-.35,1.1)--(-.35,2)node[left]{$kx+\beta v^2$};
\node[right]at(1.6,1.7){$x(0)=v(0)=0$};

\end{tikzpicture}
\caption{首次下落时弹力和阻力均向上；回升后只有阻力改变方向。}\label{fig:c2-v2-q}
\end{figure}
\begin{solution}
图~\ref{fig:c2-v2-q}取向下为正，首次下落有$v\ge0$，故
\[m\ddot x=mg-kx-\beta v^2.\]
所求的是速度与位置的关系，方程又不显含时间，宜先把$v$看作$x$的函数。令$w(x)=v^2$，在$v>0$段由$\ddot x=v\dd v/\dd x=w'/2$得
\[w'+\lambda w=2g-\frac{2k}{m}x,\qquad
\lambda=\frac{2\beta}{m},\quad w(0)=0.\]
这是关于$w$的一阶线性初值问题。乘积分因子$\ee^{\lambda x}$，从$0$积到$x$：
\[w(x)=\ee^{-\lambda x}\int_0^x
\left(2g-\frac{2k}{m}s\right)\ee^{\lambda s}\dd s.\]
对积分中的 $s\ee^{\lambda s}$ 作分部积分，
\[
\int s\ee^{\lambda s}\dd s
=\frac{s\ee^{\lambda s}}\lambda-\frac1\lambda\int\ee^{\lambda s}\dd s
=\ee^{\lambda s}\left(\frac s\lambda-\frac1{\lambda^2}\right)+C.
\]
因此定积分可写为
\[
w(x)=\ee^{-\lambda x}\left[
\ee^{\lambda s}\left(\frac{2g}{\lambda}-\frac{2ks}{m\lambda}+\frac{2k}{m\lambda^2}\right)\right]_0^x.
\]
代入两端并约去上端的指数因子，得到
\[\boxed{w(x)=\frac{2g}{\lambda}-\frac{2kx}{m\lambda}
+\frac{2k}{m\lambda^2}
-\left(\frac{2g}{\lambda}+\frac{2k}{m\lambda^2}\right)\ee^{-\lambda x}},\quad
v=+\sqrt{w(x)}.\]
最低点 $x_*$ 是除 $x=0$ 外的第一个正根 $w(x_*)=0$。积分核 $(2g-2ks/m)\ee^{\lambda s}$ 在 $s<mg/k$ 为正、在其后为负，所以积分先增加、再严格减小；显式式子又在 $x\to\infty$ 时趋负，故正根唯一且 $x_*>mg/k$。在该根处 $w'(x_*)=2g-2kx_*/m<0$，为简单转向点。

图~\ref{fig:c2-v2-s}取$B=2\beta g/k=2$，用$X=kx/(mg)$、$W=kv^2/(mg^2)$表示，曲线为$W=3/2-X-(3/2)\ee^{-2X}$。第一正根约为$X_*=1.41072$；只取$0\le X\le X_*$这一物理分支。到达最低点的时间为$\int_0^{x_*}\dd x/\sqrt{w(x)}$，两端均是可积的平方根奇性，故时间有限。

回升时 $v<0$，阻力变为向下，应重新使用 $m\ddot x=mg-kx+\beta v^2$，相应方程是 $w'-\lambda w=2g-2kx/m$。直接沿用下落的 $w(x)$ 并取负平方根，会把耗散过程误作可逆运动。

最后考察 $\beta\to0$、即 $\lambda\to0$ 的极限。先将显式解重排为
\[
w=\left(\frac{2g}{\lambda}+\frac{2k}{m\lambda^2}\right)(1-\ee^{-\lambda x})
-\frac{2kx}{m\lambda}.
\]
用 $1-\ee^{-\lambda x}=\lambda x-\lambda^2x^2/2+\lambda^3x^3/6+O(\lambda^4)$，第一项产生 $2gx+2kx/(m\lambda)-kx^2/m+O(\lambda)$；其中发散项与末项相消，留下 $w=2gx-(k/m)x^2$、$x_*=2mg/k$，与无阻力能量式相符。
\begin{figure}[H]\centering
\begin{tikzpicture}[>=Stealth,line cap=round,line join=round]

\begin{scope}[x=3cm,y=2cm]
\draw[axis](0,0)--(1.8,0)node[right]{$X$};\draw[axis](0,-.15)--(0,.65)node[above]{$W$};
\draw[curve,domain=0:1.41072,samples=101]plot(\x,{1.5-\x-1.5*exp(-2*\x)});
\draw[guide](1,0)--(1,.5);\node[below]at(1,0){$1$};\fill(1.41072,0)circle(1.5pt)node[below]{$X_*$};
\node[right]at(1.85,.4){$B=2$};
\end{scope}

\end{tikzpicture}
\caption{速度平方在首次下落区间非负；最低点位于静力平衡位置$X=1$的下方。}\label{fig:c2-v2-s}
\end{figure}
\end{solution}
